% ============================================================
% Relatório: Algoritmo de Kruskal
% Trabalho final de Estrutura de Dados e Algoritmos em Grafo
% Template: meutrabalho (v1.6)
% Compilar: pdflatex -> biber -> pdflatex -> pdflatex (compile.bat)
% ============================================================
\documentclass[a4paper]{article}
\usepackage{meutrabalho}

% ---------- Ajustes deste trabalho ----------
\usetikzlibrary{shapes.misc, backgrounds, fit}
\definecolor{grupoA}{HTML}{174A7E}   % azul do template
\definecolor{grupoB}{HTML}{2E8B57}   % verde
\definecolor{grupoC}{HTML}{C8641E}   % laranja
\definecolor{aceita}{HTML}{2E9E44}   % aresta aceita no passo

\tikzset{
    vert/.style={circle, draw=primary, fill=primary!12, thick, minimum size=7mm, inner sep=0pt, font=\small\bfseries},
    vertg/.style={circle, draw=black!45, fill=white, thick, minimum size=6mm, inner sep=0pt, font=\small\bfseries},
    peso/.style={font=\footnotesize, fill=white, inner sep=1pt},
    pino/.style={circle, fill=#1, draw=white, line width=0.8pt, minimum size=2.8mm, inner sep=0pt},
}

% Figura no padrão ABNT: título acima, fonte abaixo.
% \begin{figura}{título}{label}{fonte} ... \end{figura}
\NewDocumentEnvironment{figura}{m m m}{%
    \begin{figure}[H]
    \centering
    \caption{#1}\label{#2}
}{%
    \par\smallskip{\footnotesize Fonte: #3}
    \end{figure}
}
\newcommand{\fonte}[1]{\par\smallskip{\centering\footnotesize Fonte: #1\par}}

% Grafo do exemplo ilustrativo (Seção 4)
\newcommand{\grafobase}{%
    \coordinate (A) at (0,2.2);   \coordinate (B) at (1.5,2.45);
    \coordinate (C) at (3.1,2.15); \coordinate (D) at (0.1,0);
    \coordinate (E) at (1.6,0.35); \coordinate (F) at (3.2,0);
}
\newcommand{\arestasbase}{%
    \foreach \u/\v/\w in {A/B/4,B/C/6,A/D/3,A/E/5,B/E/2,C/E/7,C/F/3,D/E/6,E/F/4}
        \draw[black!30, dashed] (\u) -- node[peso, text=black!55] {\w} (\v);
}
% \quadro{título}{aceitas antes}{nova aceita}{vértice/cor}{nota}
\newcommand{\quadro}[5]{%
\begin{tikzpicture}[scale=1.05]
    \grafobase
    \arestasbase
    \foreach \u/\v in {#2} \draw[black, line width=1.2pt] (\u) -- (\v);
    \foreach \u/\v in {#3} \draw[aceita, line width=2pt] (\u) -- (\v);
    \foreach \n/\c in {#4} \node[circle, draw=\c!80!black, fill=\c!30, thick, minimum size=6mm, inner sep=0pt, font=\small\bfseries] at (\n) {\n};
    \node[anchor=south west, font=\small\bfseries] at (-0.4,2.75) {#1};
    \node[anchor=north west, font=\scriptsize, text width=3.9cm] at (-0.4,-0.4) {#5};
\end{tikzpicture}}

% ---------- Metadados ----------
\title{Algoritmo de Kruskal: teoria, implementação e aplicações}
\author{Nícolas Terleski de Oliveira}
\university{Universidade Estadual do Centro-Oeste --- UNICENTRO}
\course{Matemática Aplicada e Computacional}
\subject{Estrutura de Dados e Algoritmos em Grafo}
\professor{Prof. Me. João Duda}
\date{Guarapuava -- PR \\ 2026}

\begin{document}

\makecover

\resumo{
Este trabalho apresenta o algoritmo de Kruskal para o problema da árvore
geradora mínima: os conceitos de grafos em que ele se apoia, o funcionamento
passo a passo, a estrutura Union-Find, a prova de que a escolha gulosa sempre
leva ao custo mínimo e a análise de complexidade. O algoritmo foi implementado
em C++ e aplicado a um problema real, o planejamento de uma rede elétrica entre
11 cidades da região de Guarapuava (PR), e a uma extensão, a divisão de pedidos
de entrega entre motoboys.
}
\palavraschave{árvore geradora mínima; algoritmo de Kruskal; Union-Find; grafos; algoritmos gulosos}
\imprimirresumo

\sumario
\listadefiguras
\listadetabelas
\listadecodigos

% ============================================================
\section{Introdução}
% ============================================================
O algoritmo de Kruskal encontra a forma mais barata de ligar todos os pontos de
uma rede, sem nenhuma ligação sobrando. Ele foi publicado por Joseph Kruskal em
1956 \cite{kruskal1956} e resolve o problema da \textbf{árvore geradora mínima}.

O problema da árvore geradora mínima aparece em toda situação que tenha estas
três características \cite[cap.~23]{cormen2012}:
\begin{enumerate}
    \item há um conjunto de pontos que precisam ficar conectados;
    \item entre alguns pares de pontos existe uma ligação possível, e cada ligação tem um custo;
    \item basta que cada ponto consiga alcançar os outros, mesmo passando por pontos intermediários.
\end{enumerate}

Cidades numa rede elétrica, prédios numa rede de computadores e antenas num
sistema de telecomunicações são exemplos clássicos; a \cref{sec:aplicacoes}
apresenta outros, menos óbvios, como o agrupamento de dados e a segmentação de
imagens. Em todos eles o objetivo é o mesmo: escolher o conjunto de ligações que
deixa todos os pontos conectados com a menor soma de custos.

Esse problema não deve ser confundido com o do caminho mínimo. A árvore geradora
mínima minimiza o custo total da rede, mas não garante que o trajeto entre dois
pontos específicos seja o mais curto possível. Quando a pergunta é essa, por
exemplo a rota mais curta entre duas cidades, usam-se outros algoritmos, como
Dijkstra, Bellman-Ford ou Floyd-Warshall.

A ideia do Kruskal é simples: percorrer as ligações da mais barata para a mais
cara e aceitar cada uma, desde que ela não feche um ciclo com as ligações já
escolhidas. Este relatório apresenta os conceitos por trás do método, mostra por
que essa estratégia sempre dá o melhor resultado, descreve a implementação em C++
e aplica o algoritmo a um problema real: o planejamento de uma rede elétrica
entre cidades da região de Guarapuava (PR).

% ============================================================
\section{Fundamentação teórica}\label{sec:fundamentacao}
% ============================================================
O Kruskal se apoia em cinco conceitos de teoria dos grafos. Cada um vem com um
exemplo do grafo da \cref{sec:exemplo}.

\textbf{Estratégia gulosa.} Além dos conceitos de grafo, o Kruskal usa uma ideia
de projeto de algoritmos: ele é um algoritmo \textbf{guloso}. Um algoritmo
guloso monta a solução passo a passo e, em cada passo, faz a escolha que parece
melhor naquele momento, sem nunca voltar atrás. Isso nem sempre leva ao melhor
resultado: para dar 6 centavos de troco com moedas de 1, 3 e 4 centavos, pegar
sempre a maior moeda possível dá $4 + 1 + 1$ (três moedas), enquanto $3 + 3$ usa
só duas. Por isso, todo algoritmo guloso precisa de uma prova de que funciona; a
do Kruskal está na \cref{sec:prova}.

\begin{definicao}[Grafo ponderado]
Um grafo $G = (V, E)$ é formado por um conjunto de vértices $V$ e um conjunto de
arestas $E$, em que cada aresta liga dois vértices. O grafo é ponderado quando
cada aresta $e$ tem um número $w(e)$, chamado peso, que representa um custo ou
uma distância. Exemplo: a aresta B--E tem peso 2. Neste trabalho os grafos são
\textbf{não direcionados}: a aresta B--E serve nos dois sentidos.
\end{definicao}

\begin{figura}{Grafo ponderado: cada aresta liga dois vértices e tem um peso}{fig:ponderado}{elaborada pelo autor (2026).}
\begin{tikzpicture}[x=1cm,y=1cm]
    \node[vert] (A) at (0,0) {A};
    \node[vert] (B) at (4,0.6) {B};
    \node[vert] (C) at (8,0) {C};
    \node[vert] (D) at (4,-2.2) {D};
    \draw[thick, black!60] (A) -- node[peso, above] {4} (B);
    \draw[thick, black!60] (B) -- node[peso, above] {6} (C);
    \draw[thick, black!60] (A) -- node[peso, below left] {3} (D);
    \draw[thick, black!60] (B) -- node[peso, right] {2} (D);
    \draw[thick, black!60] (C) -- node[peso, below right] {5} (D);
\end{tikzpicture}
\end{figura}

\begin{definicao}[Grafo conexo]
Um grafo é conexo quando, para quaisquer dois vértices $u$ e $v$, existe um
caminho de $u$ até $v$. Não é preciso haver aresta direta: não existe aresta
D--F, mas o caminho D--E--F liga os dois.
\end{definicao}

\begin{figura}{Grafo conexo e grafo desconexo}{fig:conexo}{elaborada pelo autor (2026).}
\begin{tikzpicture}[x=1cm,y=1cm]
    \begin{scope}
        \node[font=\small\bfseries] at (1.8,1.4) {Conexo};
        \node[vert] (A) at (0,0) {A}; \node[vert] (B) at (1.8,0.5) {B};
        \node[vert] (C) at (3.6,0) {C}; \node[vert] (D) at (1.8,-1.7) {D};
        \draw[thick, black!50] (A) -- (B) -- (C) -- (D);
        \node[font=\footnotesize, text width=4.6cm, align=center] at (1.8,-2.6) {Sem aresta D--A, mas o caminho D--C--B--A liga os dois};
    \end{scope}
    \draw[black!20] (5,1.5) -- (5,-2.9);
    \begin{scope}[xshift=6.4cm]
        \node[font=\small\bfseries] at (1.8,1.4) {Desconexo};
        \node[vert] (A) at (0,0) {A}; \node[vert] (B) at (1.8,0.5) {B};
        \node[vert] (C) at (3.6,0) {C}; \node[vert] (D) at (1.8,-1.7) {D};
        \draw[thick, black!50] (A) -- (B); \draw[thick, black!50] (C) -- (D);
        \node[font=\footnotesize, text width=4.6cm, align=center] at (1.8,-2.6) {Nenhum caminho vai de A até C: são dois pedaços};
    \end{scope}
\end{tikzpicture}
\end{figura}

\begin{definicao}[Ciclo]
Um ciclo é um caminho que sai de um vértice, passa por arestas sem repetir
nenhuma e volta ao vértice de partida. Exemplo: A--B--E--A. Numa rede física, um
ciclo significa uma ligação sobrando: tirando qualquer aresta do ciclo, todos
continuam conectados.
\end{definicao}

\begin{figura}{Ciclo A--B--E--A: tirando qualquer aresta dele, todos seguem ligados}{fig:ciclo}{elaborada pelo autor (2026).}
\begin{tikzpicture}[x=1cm,y=1cm]
    \node[vert] (A) at (0,0) {A}; \node[vert] (B) at (4,0) {B};
    \node[vert] (E) at (2,-2) {E}; \node[vert] (D) at (-2,-2) {D};
    \draw[primary, line width=1.8pt] (A) -- (B) -- (E) -- (A);
    \draw[thick, black!40] (A) -- (D);
    \node[anchor=west, font=\footnotesize, text width=5cm] at (5,-0.8)
        {Em destaque: sai de A e volta a A sem repetir nenhuma aresta.\\[2pt]
         \textcolor{secondary}{A aresta A--D não faz parte do ciclo.}};
\end{tikzpicture}
\end{figura}

\begin{definicao}[Árvore]
Uma árvore é um grafo conexo e sem ciclos.
\end{definicao}

\begin{figura}{Árvore e não árvore com os mesmos 5 vértices}{fig:arvore}{elaborada pelo autor (2026).}
\begin{tikzpicture}[x=1cm,y=1cm]
    \begin{scope}
        \node[font=\small\bfseries] at (1.9,1.0) {Árvore};
        \node[vert] (A) at (1,0) {A}; \node[vert] (B) at (3.2,0) {B};
        \node[vert] (C) at (0.2,-1.8) {C}; \node[vert] (D) at (1.9,-1.8) {D};
        \node[vert] (E) at (3.8,-1.8) {E};
        \draw[thick, black!40] (A) -- (B) (A) -- (C) (A) -- (D) (B) -- (E);
        \node[font=\footnotesize] at (1.9,-2.7) {5 vértices, 4 arestas, nenhum ciclo};
    \end{scope}
    \draw[black!20] (5,1.2) -- (5,-3);
    \begin{scope}[xshift=6.2cm]
        \node[font=\small\bfseries] at (1.9,1.0) {Não é árvore};
        \node[vert] (A) at (1,0) {A}; \node[vert] (B) at (3.2,0) {B};
        \node[vert] (C) at (0.2,-1.8) {C}; \node[vert] (D) at (1.9,-1.8) {D};
        \node[vert] (E) at (3.8,-1.8) {E};
        \draw[thick, black!40] (A) -- (B) (A) -- (C) (A) -- (D) (B) -- (E);
        \draw[red!75!black, line width=1.8pt] (D) -- (E);
        \node[font=\footnotesize] at (1.9,-2.7) {A aresta a mais, D--E, fecha o ciclo A--B--E--D--A};
    \end{scope}
\end{tikzpicture}
\end{figura}

\begin{definicao}[Árvore geradora mínima]
Dado um grafo conexo e ponderado, uma árvore geradora é um subconjunto das
arestas que liga \textbf{todos} os vértices e não forma ciclo. A árvore geradora
mínima (em inglês, \emph{minimum spanning tree}, ou MST) é a árvore geradora com a
menor soma de pesos. Pode existir mais de uma quando há pesos empatados, mas
todas têm o mesmo custo.
\end{definicao}

\begin{figura}{Árvore geradora mínima do grafo da \cref{sec:exemplo}: 5 arestas, custo 16}{fig:mst}{elaborada pelo autor (2026).}
\begin{tikzpicture}[scale=2.6]
    \grafobase
    \foreach \u/\v/\w in {B/C/6,A/E/5,C/E/7,D/E/6}
        \draw[black!30, dashed] (\u) -- node[peso, text=black!55] {\w} (\v);
    \foreach \u/\v/\w in {A/B/4,A/D/3,B/E/2,C/F/3,E/F/4}
        \draw[primary, line width=1.8pt] (\u) -- node[peso, font=\footnotesize\bfseries] {\w} (\v);
    \foreach \n in {A,B,C,D,E,F} \node[vert] at (\n) {\n};
\end{tikzpicture}
\par\smallskip{\footnotesize Linhas grossas: a árvore ($2+3+3+4+4=16$). Tracejadas: arestas que ficaram de fora.}
\end{figura}

\begin{teorema}\label{teo:n-1}
Toda árvore com $n$ vértices tem exatamente $n-1$ arestas.
\end{teorema}

\begin{proof}
\textbf{Hipótese:} $T$ é uma árvore com $n$ vértices, isto é, $T$ é conexo e não
tem ciclos (Definição 4). Seja $m$ o número de arestas de $T$.
\textbf{Tese:} $m = n-1$.

Vamos montar $T$ a partir dos seus $n$ vértices, sem nenhuma aresta,
acrescentando as $m$ arestas de $T$ uma de cada vez, em qualquer ordem, e
acompanhar o número de grupos (pedaços conexos). No início há $n$ grupos, um para
cada vértice.

\emph{Afirmação:} cada aresta acrescentada diminui o número de grupos em
exatamente 1. Seja $u$--$v$ a aresta acrescentada. Se $u$ e $v$ já estivessem no
mesmo grupo, haveria um caminho entre eles formado por arestas de $T$ já
colocadas, e esse caminho junto com $u$--$v$ formaria um ciclo em $T$, o que
contradiz a hipótese. Logo $u$ e $v$ estão em grupos diferentes. Como a aresta só
tem essas duas pontas, ela une exatamente esses dois grupos e não altera os
demais: o número de grupos cai exatamente 1.

Pela afirmação, depois das $m$ arestas há $n-m$ grupos. Por outro lado, o grafo
montado ao final é o próprio $T$, que é conexo pela hipótese e, portanto, forma
exatamente 1 grupo. Assim, $n - m = 1$, ou seja, $m = n-1$, que é a tese.
\end{proof}

O \cref{teo:n-1} diz ao algoritmo quando parar: com $n-1$ arestas aceitas, a
árvore está completa.

% ============================================================
\section{Funcionamento}\label{sec:funcionamento}
% ============================================================
Como visto na \cref{sec:fundamentacao}, o Kruskal é um algoritmo guloso. Dado um
grafo conexo e ponderado com $n$ vértices, ele funciona assim:
\begin{enumerate}
    \item Ordenar todas as arestas do menor peso para o maior.
    \item Começar com uma árvore $T$ vazia. Cada vértice forma um grupo sozinho.
    \item Percorrer as arestas na ordem. Para cada aresta $u$--$v$:
    \begin{itemize}
        \item se $u$ e $v$ estão em grupos diferentes, colocar a aresta em $T$ e juntar os dois grupos;
        \item se estão no mesmo grupo, descartar a aresta, porque ela fecharia um ciclo.
    \end{itemize}
    \item Parar quando $T$ tiver $n-1$ arestas (\cref{teo:n-1}).
\end{enumerate}

Esses passos ficam mais claros na prática: na \cref{sec:exemplo} eles são
aplicados, quadro a quadro, a um grafo de exemplo.

\subsection{Como saber se dois vértices estão no mesmo grupo: Union-Find}

O passo 3 exige responder rápido ``$u$ e $v$ estão no mesmo grupo?'' e ``junte o
grupo de $u$ com o de $v$''. A estrutura usada para isso é o
\textbf{Union-Find} (conjuntos disjuntos, isto é, conjuntos que não têm nenhum
elemento em comum, como os grupos do Kruskal, em que cada vértice pertence a
exatamente um grupo). Ela guarda um vetor \texttt{pai}, em que cada vértice
aponta para outro do mesmo grupo. O vértice que aponta para si mesmo
(\texttt{pai[x] == x}) é a \textbf{raiz}, o representante do grupo. No início,
\texttt{pai[i] = i} para todo $i$: cada vértice é raiz de si mesmo.

O Union-Find é a estrutura usada na implementação clássica do Kruskal:
\textcite{cormen2012} apresentam o algoritmo com ela no capítulo 23 e dedicam o
capítulo 21 inteiro a essa estrutura. Ela não é, porém, a única forma de fazer o
teste de ciclo. Também dá para fazer uma busca em largura ou em profundidade (BFS
ou DFS) nas arestas já escolhidas, verificando se existe caminho de $u$ até $v$,
ou guardar um rótulo de grupo por vértice e renomear todos os rótulos de um grupo
a cada união. As duas alternativas dão o resultado correto, mas são mais lentas;
a \cref{sec:complexidade} compara os custos e justifica a escolha do Union-Find.

\textbf{find(x)}, que descobre a raiz do grupo de $x$:
\begin{enumerate}
    \item Olhar o pai de $x$.
    \item Se $x$ é pai de si mesmo, $x$ é a raiz e a busca termina.
    \item Se não, subir para o pai e repetir.
    \item No fim, fazer cada vértice visitado apontar direto para a raiz (compressão de caminho).
\end{enumerate}
Em poucas palavras: o find sobe pelas setas até encontrar quem aponta para si
mesmo, e aproveita a subida para encurtar o caminho da próxima vez.

\textbf{unir(a, b)}, que junta os grupos de $a$ e de $b$:
\begin{enumerate}
    \item Achar a raiz de $a$ e a raiz de $b$ com o find.
    \item Se forem a mesma, $a$ e $b$ já estão no mesmo grupo: não juntar nada, porque seria ciclo.
    \item Se forem diferentes, colocar a raiz do grupo menor embaixo da raiz do grupo maior.
\end{enumerate}
Em poucas palavras: o unir só junta grupos diferentes, e sempre ligando raiz com raiz.

\begin{figura}{Union-Find no exemplo da \cref{sec:exemplo}: árvore de pais e vetor \texttt{pai}}{fig:unionfind}{elaborada pelo autor (2026).}
\resizebox{\textwidth}{!}{\begin{tikzpicture}[x=1cm,y=1cm, raiz/.style={vert}, seta/.style={-{Stealth[length=5pt]}, black!40, thick},
    casa/.style={draw=black!30, minimum width=0.72cm, minimum height=0.85cm, align=center, font=\scriptsize, inner sep=1pt}]
    % ---- quadro 1
    \begin{scope}
        \draw[black!20, rounded corners] (-0.4,1.1) rectangle (4.7,-5.4);
        \node[anchor=west, font=\footnotesize\bfseries] at (-0.2,0.7) {1. Depois de B--E, A--D e C--F};
        \node[raiz] (A) at (0.6,-0.4) {A}; \node[raiz] (B) at (2.1,-0.4) {B}; \node[raiz] (C) at (3.6,-0.4) {C};
        \node[vertg] (D) at (0.6,-2) {D}; \node[vertg] (E) at (2.1,-2) {E}; \node[vertg] (F) at (3.6,-2) {F};
        \draw[seta] (D) -- (A); \draw[seta] (E) -- (B); \draw[seta] (F) -- (C);
        \foreach \i/\l/\p in {0/A/0,1/B/1,2/C/2,3/D/0,4/E/1,5/F/2}
            \node[casa] at (\i*0.72+0.25,-3.3) {\textcolor{secondary}{\l\,\i}\\[-1pt]\textbf{\p}};
        \node[anchor=west, font=\scriptsize, text width=4.8cm] at (-0.2,-4.4) {Três grupos; raízes A, B e C.\\\textcolor{secondary}{Raiz = quem aponta para si mesmo.}};
    \end{scope}
    % ---- quadro 2
    \begin{scope}[xshift=5.4cm]
        \draw[black!20, rounded corners] (-0.4,1.1) rectangle (4.7,-5.4);
        \node[anchor=west, font=\footnotesize\bfseries] at (-0.2,0.7) {2. find(E), depois de A--B};
        \node[raiz] (A) at (1.4,-0.2) {A}; \node[raiz] (C) at (3.8,-0.2) {C};
        \node[vertg] (B) at (0.5,-1.3) {B}; \node[vertg] (D) at (2.3,-1.3) {D}; \node[vertg] (F) at (3.8,-1.3) {F};
        \node[vertg] (E) at (0.5,-2.4) {E};
        \draw[seta] (D) -- (A); \draw[seta] (F) -- (C);
        \draw[-{Stealth[length=6pt]}, primary, line width=1.4pt] (E) -- (B);
        \draw[-{Stealth[length=6pt]}, primary, line width=1.4pt] (B) -- (A);
        \draw[-{Stealth[length=5pt]}, aceita, thick, dashed] (E) to[bend right=25] (A);
        \foreach \i/\l/\p in {0/A/0,1/B/0,2/C/2,3/D/0,4/E/0,5/F/2}
            \node[casa] (c\i) at (\i*0.72+0.25,-3.3) {\textcolor{secondary}{\l\,\i}\\[-1pt]\textbf{\p}};
        \node[casa, draw=aceita, very thick, fill=aceita!15] at (c4) {\textcolor{secondary}{E\,4}\\[-1pt]\textbf{0}};
        \node[anchor=west, font=\scriptsize, text width=4.8cm] at (-0.2,-4.4) {E $\to$ B $\to$ A: a raiz de E é A.\\\textcolor{secondary}{Compressão: \texttt{pai[4]} vai de 1 para 0.}};
    \end{scope}
    % ---- quadro 3
    \begin{scope}[xshift=10.8cm]
        \draw[black!20, rounded corners] (-0.4,1.1) rectangle (4.7,-5.4);
        \node[anchor=west, font=\footnotesize\bfseries] at (-0.2,0.7) {3. unir(E, F): raiz sob raiz};
        \node[raiz] (A) at (2.1,-0.2) {A};
        \node[vertg] (B) at (0.3,-1.3) {B}; \node[vertg] (D) at (1.4,-1.3) {D}; \node[vertg] (E) at (2.5,-1.3) {E};
        \node[vertg] (C) at (3.8,-1.3) {C}; \node[vertg] (F) at (3.8,-2.4) {F};
        \draw[seta] (B) -- (A); \draw[seta] (D) -- (A); \draw[seta] (E) -- (A); \draw[seta] (F) -- (C);
        \draw[-{Stealth[length=6pt]}, aceita, line width=1.4pt] (C) -- (A);
        \foreach \i/\l/\p in {0/A/0,1/B/0,2/C/0,3/D/0,4/E/0,5/F/2}
            \node[casa] (d\i) at (\i*0.72+0.25,-3.3) {\textcolor{secondary}{\l\,\i}\\[-1pt]\textbf{\p}};
        \node[casa, draw=aceita, very thick, fill=aceita!15] at (d2) {\textcolor{secondary}{C\,2}\\[-1pt]\textbf{0}};
        \node[anchor=west, font=\scriptsize, text width=4.8cm] at (-0.2,-4.4) {Raízes A e C são diferentes: une.\\\textcolor{secondary}{\texttt{pai[2] = 0}: C vai para baixo de A.}};
    \end{scope}
\end{tikzpicture}}
\end{figura}

Cada seta liga um vértice ao seu pai, e os círculos azuis são raízes. No
quadro 2, o find(E) sobe pelas setas em destaque até A, e a compressão
(tracejado) faz E apontar direto para A. No quadro 3, o unir(E, F) acha as raízes
A e C e coloca a raiz C, do grupo menor, embaixo de A. A casa do vetor que mudou
aparece em verde.

Um cuidado importante: quem recebe um novo pai é sempre a \textbf{raiz} do grupo,
nunca o vértice que apareceu na aresta. Se o vértice fosse ligado diretamente,
ele sairia do grupo antigo e quebraria a ligação com os outros membros.

Duas otimizações deixam o Union-Find praticamente instantâneo:
\begin{itemize}
    \item \textbf{Compressão de caminho:} durante o find, cada vértice visitado
    passa a apontar direto para a raiz, encurtando as buscas seguintes.
    \item \textbf{União por tamanho:} ao juntar dois grupos, a raiz do grupo menor
    fica embaixo da raiz do maior, o que impede a formação de cadeias longas.
\end{itemize}

Com as duas juntas, cada operação custa, em média ao longo de uma sequência de
operações (custo amortizado), $O(\alpha(n))$, em que $\alpha$ é a inversa da
função de Ackermann \cite{tarjan1975}; \cite[cap.~21]{cormen2012}. Ela cresce tão
devagar que vale no máximo 4 para qualquer $n$ de uso prático. O retorno de
\texttt{unir} já é a decisão do Kruskal: verdadeiro significa aceitar a aresta,
falso significa descartá-la.

% ============================================================
\section{Exemplo ilustrativo}\label{sec:exemplo}
% ============================================================
Até aqui o Kruskal é só uma lista de regras. Será que, seguindo apenas essas
regras, sem olhar o desenho e sem testar combinações, chegamos mesmo na forma
mais barata de ligar tudo? E quanto trabalho isso dá? Para responder, vamos rodar
o algoritmo à mão num grafo pequeno, com 6 vértices e 9 arestas
(\cref{tab:arestas}).

\begin{table}[H]
\centering
\caption{Arestas do grafo de exemplo}\label{tab:arestas}
\begin{tabular}{l*{9}{c}}
\toprule
Aresta & A--B & B--C & A--D & A--E & B--E & C--E & C--F & D--E & E--F \\
Peso   & 4 & 6 & 3 & 5 & 2 & 7 & 3 & 6 & 4 \\
\bottomrule
\end{tabular}
\fonte{elaborada pelo autor (2026).}
\end{table}

\textbf{Passo 1: ordenar.} As mesmas arestas, agora do menor peso para o maior
(\cref{tab:ordenadas}).

\begin{table}[H]
\centering
\caption{Arestas em ordem crescente de peso}\label{tab:ordenadas}
\begin{tabular}{l*{9}{c}}
\toprule
Aresta & B--E & A--D & C--F & A--B & E--F & A--E & B--C & D--E & C--E \\
Peso   & 2 & 3 & 3 & 4 & 4 & 5 & 6 & 6 & 7 \\
\bottomrule
\end{tabular}
\fonte{elaborada pelo autor (2026).}
\end{table}

\begin{figura}{Execução do Kruskal no grafo de exemplo, quadro a quadro}{fig:execucao}{elaborada pelo autor (2026).}
\quadro{Início}{}{}{A/white,B/white,C/white,D/white,E/white,F/white}{Cada vértice sozinho.\\\textcolor{secondary}{6 grupos, 0 arestas}}\hfill
\quadro{1. B--E (2): aceita}{}{B/E}{A/white,B/grupoA,C/white,D/white,E/grupoA,F/white}{B e E no mesmo grupo.\\\textcolor{secondary}{5 grupos, 1 aresta}}\hfill
\quadro{2. A--D (3): aceita}{B/E}{A/D}{A/grupoC,B/grupoA,C/white,D/grupoC,E/grupoA,F/white}{Surge o grupo de A e D.\\\textcolor{secondary}{4 grupos, 2 arestas}}

\bigskip
\quadro{3. C--F (3): aceita}{B/E,A/D}{C/F}{A/grupoC,B/grupoA,C/grupoB,D/grupoC,E/grupoA,F/grupoB}{Surge o grupo de C e F.\\\textcolor{secondary}{3 grupos, 3 arestas}}\hfill
\quadro{4. A--B (4): aceita}{B/E,A/D,C/F}{A/B}{A/grupoA,B/grupoA,C/grupoB,D/grupoA,E/grupoA,F/grupoB}{A--B junta dois grupos de 2.\\\textcolor{secondary}{2 grupos, 4 arestas}}\hfill
\quadro{5. E--F (4): aceita}{B/E,A/D,C/F,A/B}{E/F}{A/grupoA,B/grupoA,C/grupoA,D/grupoA,E/grupoA,F/grupoA}{Tudo num grupo: árvore pronta.\\\textcolor{secondary}{1 grupo, 5 arestas, custo 16}}
\end{figura}

Em cada quadro, a linha verde é a aresta aceita naquele passo, as linhas escuras
são as aceitas antes e as tracejadas ainda não entraram. Vértices da mesma cor
estão no mesmo grupo: a cada aresta aceita, dois grupos viram um.

Depois do passo 5 a árvore tem $5 = n-1$ arestas e o algoritmo para. As arestas
A--E, B--C, D--E e C--E nem são examinadas. Se fossem, todas seriam descartadas:
por exemplo, A e E já estão ligados pelo caminho A--B--E, então A--E fecharia o
ciclo A--B--E--A.

\textbf{Resultado:} B--E, A--D, C--F, A--B e E--F, com custo
$2 + 3 + 3 + 4 + 4 = \mathbf{16}$. Como há empates (A--D e C--F pesam 3; A--B e
E--F pesam 4), a ordem dentro de cada empate pode mudar, mas o custo final é
sempre 16.

Respondendo à pergunta do início: sim. Seguindo só as regras, o Kruskal chegou ao
custo 16, que é o mínimo possível (a \cref{sec:prova} prova que isso sempre
acontece), e precisou olhar apenas 5 das 9 arestas.

Vale ressaltar que o Kruskal não é o único algoritmo guloso para a árvore
geradora mínima: o algoritmo de Prim resolve o mesmo problema e também escolhe
sempre a aresta mais barata. Os dois aparecem lado a lado na literatura como
variações do mesmo método guloso \cite[cap.~23]{cormen2012}. A diferença está em
\emph{quais} arestas cada um considera a cada passo. O Prim parte de um vértice
inicial e faz crescer uma única árvore, escolhendo sempre a aresta mais barata
que liga essa árvore a um vértice ainda de fora. O Kruskal olha o grafo inteiro:
escolhe a aresta mais barata de todas que não feche ciclo, mesmo que ela fique
longe das que já foram escolhidas. Assim ele deixa várias ``ilhas'' crescerem em
lugares diferentes e vai juntando as ilhas até sobrar uma só. Neste exemplo, logo
depois de B--E o Kruskal aceita A--D, que não encosta em B nem em E; o Prim,
começando por B, não faria essa escolha naquele momento.

% ============================================================
\section{Por que o Kruskal sempre encontra o mínimo}\label{sec:prova}
% ============================================================
Estratégias gulosas costumam falhar em outros problemas, mas aqui ela funciona
por causa de um \textbf{argumento de troca}: qualquer árvore que não use as
escolhas do Kruskal pode ser trocada por uma que use, sem ficar mais cara.

\textbf{A ideia no exemplo.} Quando o Kruskal aceita A--B (peso 4), o grupo
$\{A, D\}$ já está formado. Suponha uma árvore ``rival'' que não usa A--B. Ela
também precisa ligar $\{A, D\}$ ao resto do grafo, então usa algum cabo saindo
desse grupo. Os únicos são A--B (4), A--E (5) e D--E (6), e A--B é o mais barato.
Trocando o cabo da rival por A--B, a árvore continua ligando todos os vértices e
o custo não aumenta. Logo, a rival nunca é melhor.

\begin{figura}{O argumento da troca: sair do grupo de A e D por A--B nunca custa mais}{fig:troca}{elaborada pelo autor (2026).}
\begin{tikzpicture}[scale=1.45]
    \begin{scope}
        \grafobase
        \node[anchor=south west, font=\small\bfseries] at (-0.5,2.8) {Árvore rival, sem A--B};
        \draw[primary, dashed, fill=primary!8] (0.05,1.1) ellipse (0.55 and 1.55);
        \foreach \u/\v/\w in {A/B/4,B/C/6,C/E/7,D/E/6}
            \draw[black!25, dashed] (\u) -- node[peso, text=black!50] {\w} (\v);
        \foreach \u/\v/\w in {A/D/3,B/E/2,C/F/3,E/F/4}
            \draw[black, line width=1.2pt] (\u) -- node[peso, font=\footnotesize\bfseries] {\w} (\v);
        \draw[primary, line width=2pt] (A) -- node[peso, font=\footnotesize\bfseries] {5} (E);
        \foreach \n in {A,B,C,D,E,F} \node[vertg] at (\n) {\n};
        \node[anchor=north west, font=\scriptsize, text width=5cm] at (-0.5,-0.4) {Sai do grupo de A e D pela aresta A--E (5).\\\textcolor{secondary}{Custo: $2+3+3+4+5=17$}};
    \end{scope}
    \begin{scope}[xshift=4.6cm]
        \grafobase
        \node[anchor=south west, font=\small\bfseries] at (-0.5,2.8) {Depois da troca: A--E sai, A--B entra};
        \draw[primary, dashed, fill=primary!8] (0.05,1.1) ellipse (0.55 and 1.55);
        \foreach \u/\v/\w in {A/E/5,B/C/6,C/E/7,D/E/6}
            \draw[black!25, dashed] (\u) -- node[peso, text=black!50] {\w} (\v);
        \foreach \u/\v/\w in {A/D/3,B/E/2,C/F/3,E/F/4}
            \draw[black, line width=1.2pt] (\u) -- node[peso, font=\footnotesize\bfseries] {\w} (\v);
        \draw[aceita, line width=2pt] (A) -- node[peso, font=\footnotesize\bfseries] {4} (B);
        \foreach \n in {A,B,C,D,E,F} \node[vertg] at (\n) {\n};
        \node[anchor=north west, font=\scriptsize, text width=5cm] at (-0.5,-0.4) {Sai do grupo de A e D pela aresta A--B (4).\\\textcolor{secondary}{Custo: $2+3+3+4+4=16$}};
    \end{scope}
\end{tikzpicture}
\end{figura}

A área tracejada é o grupo $\{A, D\}$. A árvore rival sai dele por A--E; trocando
essa aresta por A--B, a árvore continua ligando todos os vértices e o custo cai
de 17 para 16.

\begin{teorema}\label{teo:kruskal}
Num grafo conexo e ponderado, a árvore $T_K$ produzida pelo Kruskal é uma árvore
geradora mínima.
\end{teorema}

\begin{proof}
\textbf{Hipótese:} $G$ é um grafo conexo e ponderado com $n$ vértices, e o
Kruskal examina as arestas de $G$ em ordem crescente de peso, aceitando cada uma
cujas pontas estejam em grupos diferentes. \textbf{Tese:} a árvore $T_K$ formada
pelas arestas aceitas é uma árvore geradora mínima de $G$.

Vamos acompanhar o algoritmo e mostrar que ele nunca erra uma escolha. A
propriedade que vamos seguir, chamada (P), é: \emph{as arestas aceitas até o
momento estão todas contidas em alguma árvore geradora mínima de $G$}.

\emph{Afirmação 1:} (P) vale no início. Nenhuma aresta foi aceita e, como $G$ é
conexo, existe pelo menos uma árvore geradora mínima; o conjunto vazio está
contido nela.

\emph{Afirmação 2:} se (P) vale antes de o Kruskal aceitar uma aresta $e = u$--$v$,
então (P) continua valendo depois. Seja $T^*$ uma árvore geradora mínima que
contém todas as arestas aceitas antes de $e$. Se $e$ está em $T^*$, não há nada a
provar. Se não está, seja $X$ o grupo de $u$ no momento em que $e$ é examinada,
isto é, o conjunto dos vértices ligados a $u$ pelas arestas já aceitas. Como $e$
foi aceita, $v$ não está em $X$. ``Fora de $X$'' quer dizer entre os demais
vértices do grafo.

Como $T^*$ é uma árvore geradora, existe em $T^*$ um caminho de $u$ até $v$. Ele
começa em $u$, que está em $X$, e termina em $v$, que está fora de $X$; então, em
algum ponto, passa de um vértice de $X$ para um vértice fora de $X$. Seja $f$ a
aresta de $T^*$ usada nessa passagem. No exemplo da \cref{fig:troca},
$X = \{A, D\}$ e $f = $ A--E. Note que $f$ não foi aceita pelo Kruskal, porque
toda aresta aceita liga vértices do mesmo grupo, e $f$ liga $X$ ao resto.

Primeiro, $w(f) \ge w(e)$. Suponha, por absurdo, que $w(f) < w(e)$. Então o
Kruskal examinou $f$ antes de $e$. Naquele momento as pontas de $f$ estavam em
grupos diferentes, porque os grupos só se juntam com o tempo e elas continuam
separadas agora. Logo o Kruskal teria aceitado $f$, o que contradiz a observação
anterior.

Agora, seja $T' = T^* - f + e$. Ao retirar $f$, a árvore $T^*$ se divide em dois
pedaços sem ciclo, um contendo $u$ e outro contendo $v$, porque $f$ estava no
único caminho entre eles. A aresta $e = u$--$v$ liga esses dois pedaços
diferentes, então não fecha ciclo (pelo mesmo argumento do \cref{teo:n-1}) e
volta a deixar tudo conectado. Portanto $T'$ é conexo e sem ciclos, isto é, uma
árvore geradora, e
\[
    w(T') = w(T^*) - w(f) + w(e) \le w(T^*).
\]
Como $T^*$ é mínima, $T'$ também é. Além disso, $T'$ contém todas as arestas
aceitas antes de $e$ (nenhuma delas é $f$) e também contém $e$. Logo (P) vale
depois do passo.

\emph{Conclusão:} pelas Afirmações 1 e 2, (P) vale do início ao fim. Como $G$ é
conexo, o Kruskal chega a $n-1$ arestas aceitas: se sobrassem dois grupos ou
mais, alguma aresta de $G$ ligaria grupos diferentes e teria sido aceita quando
foi examinada. Então as $n-1$ arestas de $T_K$ estão contidas numa árvore
geradora mínima $T^*$, que também tem $n-1$ arestas (\cref{teo:n-1}). Logo
$T_K = T^*$, e $T_K$ é uma árvore geradora mínima, que é a tese.
\end{proof}

\textbf{Consequência.} Ordenando as arestas do maior peso para o menor, o mesmo
argumento, com a desigualdade invertida, mostra que o Kruskal encontra a
\textbf{árvore geradora máxima}. No exemplo da \cref{sec:exemplo}, ela custa
$7 + 6 + 6 + 5 + 4 = 28$.

% ============================================================
\section{Complexidade computacional}\label{sec:complexidade}
% ============================================================
O Kruskal com Union-Find otimizado custa $O(E \log E)$ de tempo e $O(V + E)$ de
memória, em que $V$ é o número de vértices e $E$ o de arestas. A ordenação é a
parte mais cara (\cref{tab:etapas}).

\begin{table}[H]
\centering
\caption{Custo de cada etapa do Kruskal}\label{tab:etapas}
\begin{tabular}{lll}
\toprule
Etapa & Custo de tempo & Por quê \\
\midrule
Ordenar as arestas & $O(E \log E)$ & Ordenação por comparação (\texttt{std::sort}) \\
Inicializar o Union-Find & $O(V)$ & Preenche \texttt{pai} e \texttt{tamanhoGrupo} \\
Percorrer as arestas & $O(E \cdot \alpha(V))$ & Uma chamada de \texttt{unir} (dois \texttt{find}) por aresta \\
\midrule
\textbf{Total} & $\mathbf{O(E \log E)}$ & A ordenação domina; $\alpha(V) \le 4$ na prática \\
\bottomrule
\end{tabular}
\fonte{elaborada pelo autor (2026), com base em \textcite[cap.~21 e 23]{cormen2012}.}
\end{table}

Como um grafo simples (sem laços e sem duas arestas ligando o mesmo par de
vértices) tem no máximo $V(V-1)/2$ arestas (no máximo uma para cada par de
vértices, e o número de pares é a combinação de $V$ tomados 2 a 2), temos
$E < V^2$ e, portanto, $\log E < 2 \log V$. Então o total também pode ser escrito
como $O(E \log V)$ \cite[cap.~23]{cormen2012}.

\textbf{Memória.} A lista de arestas ocupa $O(E)$, a cópia ordenada mais $O(E)$, e
os vetores \texttt{pai} e \texttt{tamanhoGrupo} ocupam $O(V)$. O total é
$O(V + E)$.

\begin{complexidade}
Tempo: $T(V,E) = O(E \log E) = O(E \log V)$ \qquad Memória: $S(V,E) = O(V + E)$
\end{complexidade}

\textbf{Comparação com alternativas para o teste de ciclo.} A
\cref{tab:alternativas} compara as formas de responder ``mesmo grupo?'', e a
\cref{fig:custos} mostra como esses custos crescem.

\begin{table}[H]
\centering
\caption{Alternativas para o teste de ciclo}\label{tab:alternativas}
\small
\begin{tabularx}{\textwidth}{>{\raggedright\arraybackslash}X >{\raggedright\arraybackslash}p{3.3cm} >{\raggedright\arraybackslash}p{3.3cm}}
\toprule
Método para ``mesmo grupo?'' & Custo por aresta & Total do Kruskal \\
\midrule
Busca (BFS/DFS) nas arestas já escolhidas & $O(V)$ & $O(E \cdot V)$ \\
Vetor de rótulos \texttt{grupo[v]}, renomeando ao juntar & $O(1)$ consulta, $O(V)$ união & $O(E \log E + V^2)$ \\
Union-Find sem otimizações & até $O(V)$ & $O(E \cdot V)$ no pior caso \\
\textbf{Union-Find com compressão e união por tamanho} & $\mathbf{O(\alpha(V))}$ & $\mathbf{O(E \log E)}$ \\
\bottomrule
\end{tabularx}
\fonte{elaborada pelo autor (2026).}
\end{table}

\begin{figura}{Operações do Kruskal para cada método de teste de ciclo, com $E = 5V$}{fig:custos}{elaborada pelo autor (2026), a partir das fórmulas da \cref{tab:alternativas}.}
\begin{tikzpicture}
\begin{loglogaxis}[
    width=0.82\textwidth, height=7cm,
    xlabel={Número de vértices ($V$)}, ylabel={Operações (aproximadas)},
    xmin=80, xmax=12000, ymin=1e3, ymax=2e9,
    grid=major, grid style={black!10},
    legend pos=north west, legend style={font=\footnotesize, draw=none, fill=white, fill opacity=0.8, text opacity=1},
    domain=100:10000, samples=30,
]
    \addplot[black!55, thick, mark=*, mark size=1.5pt, samples at={100,300,1000,3000,10000}] {5*x*x};
    \addlegendentry{Busca BFS/DFS: $E \cdot V$}
    \addplot[black!35, thick, dashed, mark=*, mark size=1.5pt, samples at={100,300,1000,3000,10000}] {5*x*ln(5*x)/ln(2) + x*x};
    \addlegendentry{Vetor de rótulos: $E \log E + V^2$}
    \addplot[primary, very thick, mark=*, mark size=1.8pt, samples at={100,300,1000,3000,10000}] {5*x*ln(5*x)/ln(2)};
    \addlegendentry{Union-Find: $E \log E$}
\end{loglogaxis}
\end{tikzpicture}
\end{figura}

As três curvas dão o mesmo resultado, mas o trabalho cresce em ritmos bem
diferentes: com 10 mil vértices, a busca já passa da casa das centenas de milhões
de operações, enquanto o Union-Find fica abaixo de um milhão.

A escolha da \textbf{lista de arestas} como representação do grafo também vem da
complexidade: o Kruskal nunca pergunta quem são os vizinhos de um vértice, só
percorre as arestas em ordem de peso. Uma matriz de adjacência gastaria $O(V^2)$
de memória e exigiria varrer a matriz inteira para montar essa lista.

% ============================================================
\section{Aplicações}\label{sec:aplicacoes}
% ============================================================
O Kruskal é usado sempre que é preciso ligar pontos com o menor custo total, e
também em problemas onde não há ``cabo'' nenhum (\cref{tab:aplicacoes}).

\begin{table}[H]
\centering
\caption{Aplicações da árvore geradora mínima}\label{tab:aplicacoes}
\small
\begin{tabularx}{\textwidth}{>{\raggedright\arraybackslash}X >{\raggedright\arraybackslash}X >{\raggedright\arraybackslash}X}
\toprule
Aplicação & Vértices & Arestas e pesos \\
\midrule
Redes elétricas, de água e de fibra óptica & Cidades, bairros ou prédios & Trechos possíveis; custo ou distância \\
Redes de computadores e telecomunicações & Servidores, antenas, estações & Ligações possíveis; custo de instalação \\
Agrupamento de dados (clusterização por ligação simples) & Itens a agrupar & Distância entre itens; para-se o Kruskal com $k$ grupos \\
Segmentação de imagens & Pixels & Diferença de cor entre vizinhos \\
Rastreio de surtos de doenças & Amostras de bactérias ou vírus & Número de diferenças genéticas \\
Geração de labirintos em jogos & Salas de uma grade & Paredes, com pesos aleatórios \\
\bottomrule
\end{tabularx}
\fonte{elaborada pelo autor (2026).}
\end{table}

Duas aplicações menos óbvias merecem destaque. Na \textbf{segmentação de
imagens}, juntar primeiro os pixels de cor mais parecida, com um critério de
parada, divide a foto em regiões como céu, pessoa e chão; é o que faz o
algoritmo de Felzenszwalb e Huttenlocher, que percorre as arestas em ordem
crescente de peso, como o Kruskal, e usa Union-Find \cite{felzenszwalb2004}. No
\textbf{rastreio de epidemias}, laboratórios de vigilância comparam o material
genético de amostras de vários pacientes, e a árvore geradora mínima dessas
amostras ajuda a sugerir quem se contaminou a partir de quem, porque amostras com
poucas diferenças tendem a estar próximas na cadeia de transmissão. A ferramenta
GrapeTree, por exemplo, usa árvores geradoras mínimas para mostrar as relações
genéticas entre mais de 100 mil genomas de bactérias causadoras de doenças
\cite{zhou2018}.

Neste trabalho, o problema principal (\cref{sec:rede}) é uma rede elétrica, e a
extensão (\cref{sec:motoboys}) usa a ideia da clusterização para dividir pedidos
entre entregadores.

% ============================================================
\section{Implementação em C++}\label{sec:implementacao}
% ============================================================
O programa foi dividido em módulos com uma responsabilidade cada, compilado com
\texttt{g++ -std=c++17 -Wall *.cpp -o kruskal.exe}. O Kruskal e o Union-Find
trabalham só com índices inteiros; apenas a leitura e a impressão conhecem os
nomes dos vértices (\cref{tab:modulos}).

O código-fonte completo, os arquivos de dados e as instruções de uso estão no
repositório do projeto:
\url{https://github.com/Nicolasponge/Trabalho-de-Kruskal}.

\begin{table}[H]
\centering
\caption{Módulos do programa}\label{tab:modulos}
\small
\begin{tabularx}{\textwidth}{>{\raggedright\arraybackslash}p{3.2cm} >{\raggedright\arraybackslash}p{3.4cm} >{\raggedright\arraybackslash}X}
\toprule
Arquivo & Responsabilidade & Principais elementos \\
\midrule
\texttt{grafo.hpp/.cpp} & Representar o grafo & \texttt{struct Aresta} (origem, destino, peso); \texttt{struct Grafo} (número de vértices, nomes, lista de arestas); \texttt{inicializaGrafo}, \texttt{defineNomeVertice}, \texttt{adicionaAresta}, \texttt{imprimeGrafo} \\
\texttt{union\_find.hpp/.cpp} & Controlar os grupos de vértices & \texttt{struct UnionFind} (\texttt{pai}, \texttt{tamanhoGrupo}); \texttt{inicializaUF}, \texttt{find} com compressão de caminho, \texttt{unir} com união por tamanho \\
\texttt{kruskal.hpp/.cpp} & O algoritmo & \texttt{struct ResultadoKruskal} (arestas escolhidas, custo total); \texttt{kruskal(grafo, k)} \\
\texttt{leitura.hpp/.cpp} & Ler o grafo de um arquivo \texttt{.txt} & \texttt{lerGrafo}, com um \texttt{unordered\_map} que traduz nome para índice \\
\texttt{visualizacao.hpp/.cpp} & Gerar o desenho do grafo & \texttt{exportarDot}, que escreve um arquivo no formato DOT do Graphviz \\
\texttt{main.cpp} & Menu e apresentação dos resultados & \texttt{executa}, \texttt{imprimeResultado}, menu com as opções \\
\bottomrule
\end{tabularx}
\fonte{elaborada pelo autor (2026).}
\end{table}

\subsection{Decisões de projeto}\label{sec:decisoes}
\begin{itemize}
    \item \textbf{Lista de arestas} em vez de matriz ou lista de adjacência,
    porque o Kruskal só precisa percorrer as arestas em ordem de peso
    (\cref{sec:complexidade}).
    \item \textbf{Peso em \texttt{double}}, porque os problemas usam distâncias
    com casas decimais.
    \item \textbf{Vértices como inteiros de 0 a $n-1$}, para que o vetor
    \texttt{pai} do Union-Find seja acessado diretamente pelo índice. A tradução
    nome para índice existe só na leitura; a tradução índice para nome é o
    próprio vetor \texttt{nomesVertices}.
    \item \textbf{Parâmetro $k$} na função \texttt{kruskal}: o laço para quando a
    árvore tem $n-k$ arestas. Com $k = 1$ sai a árvore geradora mínima completa;
    com $k > 1$ sobram $k$ grupos (\cref{sec:motoboys}).
    \item \textbf{\texttt{unir} devolve \texttt{bool}}: verdadeiro quando junta
    dois grupos, falso quando os vértices já estavam no mesmo grupo. Esse retorno
    já é a decisão de aceitar ou descartar a aresta.
    \item \textbf{Visualização com o Graphviz} \cite{graphviz}, e não desenhada à
    mão. No terminal, o grafo vira uma lista de arestas difícil de imaginar, ainda
    mais com muitos vértices. Desenhar em HTML e CSS exigiria decidir sozinho onde
    colocar cada vértice sem que as linhas se cruzassem, um problema que não é o
    foco do trabalho. O Graphviz já resolve isso: o programa só escreve um
    arquivo de texto no formato DOT, sem nenhuma biblioteca extra no C++, e a
    ferramenta desenha o grafo com a árvore destacada. Essa parte usa uma
    ferramenta externa; o Kruskal e o Union-Find são implementação própria.
    \begin{itemize}
        \item Para ver o desenho há dois caminhos: colar o conteúdo do arquivo
        \texttt{.dot} no Graphviz Online
        (\url{https://dreampuf.github.io/GraphvizOnline}), que desenha no
        navegador sem instalar nada, ou instalar o Graphviz
        (\url{https://graphviz.org/download/}) e rodar
        \texttt{dot -Tpng arquivo.dot -o arquivo.png} no terminal, que gera uma
        imagem PNG.
    \end{itemize}
\end{itemize}

\subsection{Problemas e dúvidas durante o desenvolvimento}

Alguns erros apareceram durante a programação e ajudaram a entender melhor o
algoritmo e a linguagem (\cref{tab:problemas}).

\begin{table}[H]
\centering
\caption{Problemas encontrados e soluções}\label{tab:problemas}
\small
\begin{tabularx}{\textwidth}{>{\raggedright\arraybackslash}p{3.4cm} >{\raggedright\arraybackslash}X >{\raggedright\arraybackslash}X}
\toprule
Problema & Por que acontecia & Como foi resolvido \\
\midrule
Os nomes dos vértices viravam números & \texttt{adicionaAresta} gravava um nome padrão (o índice) por cima do nome já definido & O nome padrão passou a ser dado uma única vez, em \texttt{inicializaGrafo} \\
Escrita fora do vetor de nomes & O vetor era esvaziado com \texttt{clear()} e depois acessado por índice & Usar \texttt{resize(n)} para o vetor já nascer com uma posição por vértice \\
Grupos se ``quebrando'' no Union-Find & Ligar o vértice da aresta, e não a raiz do grupo, tira o vértice do grupo antigo & \texttt{unir} chama \texttt{find} nas duas pontas e liga raiz com raiz \\
Tamanhos antigos no Union-Find & \texttt{resize(n, 1)} só preenche posições novas & Trocar por \texttt{assign(n, 1)}, que recomeça o vetor inteiro \\
\texttt{unir} chamado duas vezes por aresta & O \texttt{if} já executa a união; a segunda chamada não fazia nada & Usar só o retorno de \texttt{unir} como decisão \\
Aviso do compilador na parada do laço & Comparação entre \texttt{size()} (sem sinal) e \texttt{int} (com sinal) & Converter com \texttt{static\_cast<int>} \\
Acentos embaralhados no terminal do Windows & O console não usava UTF-8 & \texttt{SetConsoleOutputCP(CP\_UTF8)} no início do \texttt{main} \\
\bottomrule
\end{tabularx}
\fonte{elaborada pelo autor (2026).}
\end{table}

Também surgiram dúvidas de projeto, que viraram as decisões da \cref{sec:decisoes}:
\begin{itemize}
    \item \textbf{Vértices como letras ou como números?} Números, porque o vetor
    \texttt{pai} é acessado direto pelo índice; letras exigiriam conversão o tempo
    todo e limitariam o grafo a 26 vértices.
    \item \textbf{Precisa de uma tabela hash própria para os nomes?} Não. O índice
    para nome já é o vetor de nomes, e o nome para índice usa o
    \texttt{unordered\_map} da biblioteca padrão, só na leitura do arquivo.
    \item \textbf{Por que \texttt{find} não recebe o Union-Find como
    \texttt{const}?} Porque a compressão de caminho altera o vetor \texttt{pai}
    durante a busca.
\end{itemize}

\subsection{O núcleo do algoritmo}

O \cref{lst:kruskal} mostra a função \texttt{kruskal}, e o \cref{lst:uf}, as
operações do Union-Find.

\codigocpp[caption={Função \texttt{kruskal} (\texttt{kruskal.cpp}).}, label={lst:kruskal}]{code/kruskal.cpp}

\codigocpp[caption={Operações do Union-Find (\texttt{union\_find.cpp}).}, label={lst:uf}]{code/union_find.cpp}

\subsection{Formato de entrada}

Cada grafo fica num arquivo da pasta \texttt{dados/}. A primeira linha tem o
número de vértices e de arestas; a segunda, os nomes dos vértices (com
\texttt{\_} no lugar de espaço); as demais, uma aresta por linha.

\begin{entrada}
6 9\\
A B C D E F\\
A B 4\\
B C 6\\
\ldots
\end{entrada}

O programa mostra o grafo lido, as arestas escolhidas em ordem, o custo total e,
quando $k > 1$, os grupos formados. Também grava em \texttt{saida/} dois arquivos
\texttt{.dot}, um do grafo original e outro com a árvore destacada em vermelho,
que podem ser desenhados com o Graphviz. A visualização usa essa ferramenta
externa; o algoritmo de Kruskal e o Union-Find são implementação própria.

% ============================================================
\section{Problema do cotidiano: rede elétrica na região de Guarapuava}\label{sec:rede}
% ============================================================
Para ligar 11 cidades da região de Guarapuava (PR) numa rede elétrica, o Kruskal
escolheu 10 linhas com \textbf{294,6~km} no total, entre 25 trechos possíveis que
somam 1.016,1~km.

\subsection{Descrição do problema}

Uma concessionária quer interligar as cidades de uma região por linhas de
transmissão, de modo que a energia possa chegar a todas elas, direta ou
indiretamente. Cada quilômetro de linha tem custo de construção e manutenção,
então o objetivo é conectar todas as cidades com a menor extensão total de
linha, sem trechos redundantes.

\begin{figura}{As 11 cidades do problema, no centro-sul do Paraná}{fig:mapa}{elaborada pelo autor (2026), com base na malha municipal do IBGE disponibilizada pelo projeto geodata-br \cite{geodatabr}.}
\resizebox{\textwidth}{!}{\begin{tikzpicture}[x=1cm,y=1cm,font=\footnotesize]
\begin{scope}
\clip (-0.40,0.45) rectangle (11.40,-9.95);
\draw[black!20,line width=0.3pt] (5.04,-9.38) -- (5.18,-9.29) -- (5.36,-9.33) -- (5.39,-9.41) -- (5.63,-9.44) -- (5.58,-9.59) -- (5.79,-9.67) -- (5.74,-9.79) -- (5.84,-9.93) -- (5.94,-9.80) -- (6.16,-9.73) -- (6.28,-9.90) -- (6.53,-9.84) -- (6.75,-9.96) -- (6.78,-9.75) -- (6.92,-9.84) -- (6.97,-10.03) -- (7.03,-10.03) -- (7.08,-9.89) -- (7.20,-9.94) -- (7.21,-10.14) -- (7.29,-10.21) -- (7.50,-10.21) -- (7.66,-10.02) -- (7.77,-10.11) -- (7.62,-10.32) -- (4.75,-10.32) -- (4.84,-10.02) -- (4.55,-9.67) -- (4.66,-9.40) -- (4.54,-9.26) -- (4.60,-9.21) -- (4.52,-9.09) -- (4.75,-8.70) -- (4.84,-8.67) -- (4.94,-8.79) -- (4.73,-9.10) -- cycle;
\draw[black!20,line width=0.3pt] (6.31,0.81) -- (6.47,0.73) -- (6.53,0.63) -- (6.52,0.49) -- (6.69,0.43) -- (6.58,0.37) -- (6.75,0.20) -- (6.77,0.02) -- (6.63,-0.15) -- (6.68,-0.23) -- (6.53,-0.35) -- (6.63,-0.53) -- (6.51,-0.50) -- (6.50,-0.66) -- (6.47,-0.60) -- (6.36,-0.70) -- (6.54,-0.82) -- (6.45,-0.89) -- (6.50,-1.00) -- (6.35,-0.99) -- (6.35,-1.08) -- (6.27,-1.05) -- (6.30,-1.09) -- (6.21,-1.22) -- (6.14,-1.17) -- (6.05,-1.22) -- (6.12,-1.50) -- (6.03,-1.42) -- (5.95,-1.49) -- (5.90,-1.42) -- (5.89,-1.49) -- (5.70,-1.57) -- (5.64,-1.70) -- (5.56,-1.63) -- (5.57,-1.70) -- (5.49,-1.75) -- (5.36,-1.71) -- (5.37,-1.83) -- (5.22,-1.84) -- (5.20,-1.97) -- (4.75,-1.93) -- (4.64,-1.79) -- (4.42,-1.71) -- (4.30,-1.74) -- (4.29,-1.61) -- (4.51,-1.55) -- (4.56,-1.30) -- (4.27,-0.93) -- (4.40,-0.89) -- (4.48,-1.00) -- (4.73,-1.02) -- (4.76,-1.11) -- (5.21,-0.85) -- (5.28,-0.77) -- (5.27,-0.64) -- (5.40,-0.64) -- (5.45,-0.46) -- (5.49,-0.50) -- (5.63,-0.36) -- (5.66,-0.41) -- (5.66,-0.28) -- (5.73,-0.31) -- (5.75,-0.24) -- (5.94,-0.16) -- (5.86,-0.08) -- (5.94,0.01) -- (5.78,0.09) -- (5.79,0.20) -- (5.90,0.12) -- (5.95,0.29) -- (6.04,0.34) -- (6.08,0.76) -- cycle;
\draw[black!20,line width=0.3pt] (8.98,0.82) -- (9.10,0.46) -- (9.08,-0.47) -- (9.33,-0.76) -- (9.54,-0.82) -- (9.78,-1.34) -- (9.67,-1.43) -- (9.63,-1.59) -- (9.53,-1.54) -- (9.47,-1.63) -- (9.33,-1.51) -- (9.13,-1.49) -- (9.17,-1.34) -- (9.27,-1.32) -- (9.24,-1.21) -- (9.20,-1.30) -- (9.05,-1.28) -- (9.05,-1.37) -- (8.91,-1.34) -- (8.86,-1.23) -- (8.90,-1.17) -- (8.97,-1.21) -- (8.95,-1.13) -- (8.71,-1.13) -- (8.87,-0.95) -- (8.74,-0.99) -- (8.83,-0.85) -- (8.65,-0.90) -- (8.65,-0.74) -- (8.54,-0.70) -- (8.56,-0.56) -- (8.48,-0.49) -- (8.52,-0.42) -- (8.39,-0.44) -- (8.29,-0.35) -- (8.30,-0.22) -- (8.20,-0.35) -- (8.10,-0.32) -- (8.16,-0.09) -- (7.94,-0.19) -- (7.95,-0.29) -- (7.78,-0.24) -- (7.81,-0.13) -- (7.63,-0.16) -- (7.67,0.00) -- (7.57,-0.04) -- (7.54,-0.17) -- (7.38,-0.12) -- (7.33,-0.02) -- (7.10,-0.05) -- (7.00,-0.15) -- (6.96,-0.01) -- (6.87,-0.08) -- (6.76,-0.03) -- (6.75,0.20) -- (6.58,0.37) -- (6.69,0.43) -- (6.52,0.49) -- (6.53,0.63) -- (6.35,0.82) -- cycle;
\draw[black!20,line width=0.3pt] (-0.53,-6.93) -- (-0.47,-6.85) -- (-0.36,-6.91) -- (-0.21,-6.87) -- (0.02,-7.06) -- (0.37,-7.11) -- (0.64,-6.90) -- (0.71,-6.99) -- (0.55,-7.23) -- (0.79,-7.13) -- (0.99,-7.25) -- (1.07,-7.40) -- (1.22,-7.37) -- (1.23,-7.43) -- (1.35,-7.49) -- (1.24,-7.55) -- (1.28,-7.76) -- (1.12,-7.91) -- (1.14,-8.00) -- (1.07,-8.05) -- (1.10,-8.09) -- (1.05,-8.17) -- (0.95,-8.15) -- (1.02,-8.24) -- (0.88,-8.29) -- (0.89,-8.45) -- (0.81,-8.43) -- (0.82,-8.49) -- (0.74,-8.50) -- (0.76,-8.56) -- (0.59,-8.58) -- (0.62,-8.70) -- (0.58,-8.66) -- (0.54,-8.74) -- (0.45,-8.72) -- (0.51,-8.77) -- (0.44,-8.87) -- (-0.74,-8.78) -- (-0.74,-6.69) -- (-0.57,-6.78) -- cycle;
\draw[black!20,line width=0.3pt] (2.57,-8.91) -- (2.56,-9.10) -- (2.72,-9.16) -- (2.69,-9.32) -- (2.85,-9.43) -- (2.93,-9.41) -- (3.02,-9.22) -- (3.10,-9.21) -- (3.54,-9.37) -- (3.70,-9.50) -- (4.17,-9.30) -- (4.27,-9.29) -- (4.36,-9.39) -- (4.54,-9.26) -- (4.66,-9.40) -- (4.55,-9.67) -- (4.84,-10.02) -- (4.75,-10.32) -- (2.12,-10.32) -- (2.10,-10.20) -- (2.05,-10.19) -- (2.06,-9.99) -- (2.20,-9.73) -- (2.13,-9.69) -- (2.25,-9.55) -- (2.23,-9.46) -- (2.09,-9.46) -- (2.16,-9.32) -- (2.25,-9.34) -- (2.21,-9.30) -- (2.27,-9.23) -- (2.20,-9.14) -- (2.26,-9.07) -- (2.10,-8.93) -- (2.20,-8.84) -- (2.42,-8.83) -- cycle;
\draw[black!20,line width=0.3pt] (-0.74,-8.78) -- (-0.23,-8.78) -- (0.21,-8.89) -- (0.20,-9.32) -- (0.27,-9.35) -- (0.35,-9.52) -- (0.22,-9.69) -- (0.38,-9.80) -- (0.05,-9.97) -- (-0.01,-10.09) -- (-0.25,-10.08) -- (-0.36,-9.99) -- (-0.44,-10.03) -- (-0.49,-9.94) -- (-0.54,-9.99) -- (-0.68,-9.95) -- (-0.71,-10.10) -- cycle;
\draw[black!20,line width=0.3pt] (9.72,-6.94) -- (9.82,-6.95) -- (9.91,-7.38) -- (9.83,-7.37) -- (9.80,-7.47) -- (9.68,-7.46) -- (9.59,-7.59) -- (9.64,-7.70) -- (9.39,-7.74) -- (9.26,-7.87) -- (9.29,-8.05) -- (9.35,-8.05) -- (9.49,-8.31) -- (9.44,-8.40) -- (9.48,-8.60) -- (9.72,-8.80) -- (9.95,-8.86) -- (9.92,-8.96) -- (9.74,-8.92) -- (9.63,-9.00) -- (9.53,-8.96) -- (9.34,-9.03) -- (9.26,-9.32) -- (9.19,-9.32) -- (9.18,-9.38) -- (9.08,-9.34) -- (9.05,-9.41) -- (8.85,-9.33) -- (8.69,-9.56) -- (8.52,-9.54) -- (8.48,-9.46) -- (8.45,-9.51) -- (8.44,-9.38) -- (8.40,-9.39) -- (8.30,-9.45) -- (8.38,-9.48) -- (8.36,-9.55) -- (8.23,-9.56) -- (8.19,-9.68) -- (8.02,-9.62) -- (7.98,-9.83) -- (7.88,-9.88) -- (7.89,-10.02) -- (7.66,-10.02) -- (7.45,-10.23) -- (7.21,-10.14) -- (7.20,-9.94) -- (7.08,-9.89) -- (7.03,-10.03) -- (6.97,-10.03) -- (6.92,-9.84) -- (6.78,-9.75) -- (6.75,-9.96) -- (6.53,-9.84) -- (6.28,-9.90) -- (6.16,-9.73) -- (5.94,-9.80) -- (5.84,-9.93) -- (5.74,-9.79) -- (5.79,-9.67) -- (5.58,-9.59) -- (5.63,-9.44) -- (5.39,-9.41) -- (5.38,-9.29) -- (5.60,-9.29) -- (5.61,-9.13) -- (5.82,-9.21) -- (5.89,-9.13) -- (5.88,-9.05) -- (6.03,-9.07) -- (6.31,-8.72) -- (6.37,-8.79) -- (6.34,-8.88) -- (6.50,-8.94) -- (6.62,-8.86) -- (6.58,-8.76) -- (6.62,-8.73) -- (6.75,-8.75) -- (6.74,-8.93) -- (6.79,-8.81) -- (6.96,-8.81) -- (6.94,-8.69) -- (7.03,-8.69) -- (6.96,-8.52) -- (7.13,-8.50) -- (7.05,-8.44) -- (7.27,-8.29) -- (7.25,-8.23) -- (7.35,-8.16) -- (7.36,-8.01) -- (7.53,-8.07) -- (7.63,-7.75) -- (7.79,-7.86) -- (7.86,-7.82) -- (7.86,-7.64) -- (7.92,-7.72) -- (8.06,-7.61) -- (8.17,-7.62) -- (8.15,-7.54) -- (8.28,-7.60) -- (8.27,-7.50) -- (8.38,-7.49) -- (8.37,-7.40) -- (8.66,-7.56) -- (8.75,-7.46) -- (8.70,-7.35) -- (8.74,-7.22) -- (8.90,-7.12) -- (9.07,-7.11) -- (9.08,-7.02) -- (9.28,-7.13) -- (9.43,-7.10) -- (9.46,-6.88) -- (9.58,-6.90) -- (9.65,-6.81) -- cycle;
\draw[black!20,line width=0.3pt] (10.23,-2.40) -- (10.57,-2.57) -- (10.76,-2.81) -- (10.80,-3.00) -- (11.09,-3.03) -- (11.36,-3.20) -- (11.39,-3.30) -- (11.25,-3.48) -- (11.20,-3.77) -- (11.13,-3.84) -- (10.92,-3.92) -- (10.85,-3.89) -- (10.75,-3.77) -- (10.80,-3.59) -- (10.51,-3.48) -- (10.35,-3.32) -- (10.27,-3.35) -- (10.18,-3.13) -- (10.09,-3.15) -- (10.05,-3.05) -- (9.97,-3.06) -- (9.97,-3.14) -- (9.92,-3.10) -- (9.87,-3.00) -- (10.02,-2.60) -- (9.95,-2.39) -- cycle;
\draw[black!20,line width=0.3pt] (0.59,-9.81) -- (0.74,-9.65) -- (0.80,-9.69) -- (0.73,-9.88) -- (0.77,-10.07) -- (0.80,-10.13) -- (1.02,-10.11) -- (1.07,-10.32) -- (-0.44,-10.32) -- (-0.59,-10.26) -- (-0.71,-10.10) -- (-0.68,-9.95) -- (-0.54,-9.99) -- (-0.49,-9.94) -- (-0.44,-10.03) -- (-0.36,-9.99) -- (-0.25,-10.08) -- (-0.01,-10.09) -- (0.05,-9.97) -- (0.38,-9.80) -- (0.22,-9.69) -- (0.26,-9.61) -- (0.38,-9.61) -- (0.38,-9.68) -- (0.55,-9.73) -- cycle;
\draw[black!20,line width=0.3pt] (1.12,-10.32) -- (1.12,-10.32) -- (1.12,-10.32) -- cycle;
\draw[black!20,line width=0.3pt] (1.28,-10.31) -- (1.33,-10.32) -- (1.22,-10.32) -- cycle;
\draw[black!20,line width=0.3pt] (11.74,-2.47) -- (11.74,-4.58) -- (11.40,-4.58) -- (11.37,-4.49) -- (11.07,-4.58) -- (10.97,-4.69) -- (10.79,-4.47) -- (10.76,-4.28) -- (10.85,-4.14) -- (10.85,-3.89) -- (10.92,-3.92) -- (11.20,-3.77) -- (11.25,-3.48) -- (11.39,-3.30) -- (11.36,-3.20) -- (11.09,-3.03) -- (11.02,-2.79) -- (11.05,-2.68) -- (11.13,-2.61) -- (11.19,-2.66) -- (11.27,-2.42) -- (11.52,-2.13) -- (11.60,-2.37) -- cycle;
\draw[black!20,line width=0.3pt] (11.71,-2.46) -- (11.60,-2.37) -- (11.56,-1.85) -- (11.18,-1.69) -- (11.20,-1.63) -- (11.02,-1.50) -- (10.91,-1.28) -- (10.83,-1.29) -- (10.79,-1.08) -- (10.54,-0.77) -- (10.61,-0.80) -- (10.74,-0.69) -- (10.99,-0.76) -- (11.22,-0.59) -- (11.33,-0.73) -- (11.46,-0.76) -- (11.69,-0.77) -- (11.74,-0.71) -- cycle;
\draw[black!20,line width=0.3pt] (11.40,-4.58) -- (11.74,-4.58) -- (11.74,-6.09) -- (11.40,-6.01) -- (11.36,-5.94) -- (11.24,-5.99) -- (11.26,-6.05) -- (11.11,-6.11) -- (11.08,-6.28) -- (11.00,-6.25) -- (10.85,-6.38) -- (10.89,-6.47) -- (10.71,-6.47) -- (10.68,-6.34) -- (10.56,-6.30) -- (10.23,-6.37) -- (10.17,-6.21) -- (10.04,-6.35) -- (9.90,-6.36) -- (9.54,-6.23) -- (9.63,-6.12) -- (9.22,-6.07) -- (9.23,-5.99) -- (9.09,-5.95) -- (9.06,-6.01) -- (8.86,-5.97) -- (8.68,-5.74) -- (8.59,-5.78) -- (8.55,-5.63) -- (8.63,-5.58) -- (8.65,-5.70) -- (8.94,-5.61) -- (9.06,-5.43) -- (9.24,-5.37) -- (9.26,-5.28) -- (9.34,-5.34) -- (9.38,-5.24) -- (9.42,-5.27) -- (9.39,-5.15) -- (9.67,-5.05) -- (10.20,-4.68) -- (10.32,-4.72) -- (10.42,-4.63) -- (10.65,-4.62) -- (10.70,-4.71) -- (10.87,-4.78) -- (10.97,-4.95) -- (11.00,-4.83) -- (10.93,-4.74) -- (11.07,-4.58) -- (11.37,-4.49) -- cycle;
\draw[black!20,line width=0.3pt] (10.43,-0.61) -- (10.54,-0.73) -- (10.61,-0.92) -- (10.79,-1.08) -- (10.83,-1.29) -- (10.91,-1.28) -- (11.02,-1.50) -- (11.20,-1.63) -- (11.18,-1.69) -- (11.56,-1.85) -- (11.55,-2.13) -- (11.27,-2.42) -- (11.19,-2.66) -- (11.13,-2.61) -- (11.05,-2.68) -- (11.05,-3.05) -- (10.84,-2.96) -- (10.80,-3.00) -- (10.76,-2.81) -- (10.57,-2.57) -- (10.36,-2.42) -- (10.19,-2.36) -- (9.85,-2.38) -- (9.81,-2.19) -- (9.86,-2.12) -- (9.71,-1.99) -- (9.71,-1.88) -- (9.53,-1.87) -- (9.57,-1.71) -- (9.41,-1.69) -- (9.41,-1.61) -- (9.53,-1.61) -- (9.53,-1.54) -- (9.63,-1.59) -- (9.67,-1.43) -- (9.78,-1.34) -- (9.52,-0.82) -- (9.87,-0.72) -- (9.89,-0.67) -- (10.10,-0.77) -- (10.15,-0.64) -- (10.34,-0.59) -- cycle;
\draw[black!20,line width=0.3pt] (-0.38,0.02) -- (-0.36,0.07) -- (-0.23,-0.01) -- (-0.20,0.03) -- (-0.07,-0.14) -- (-0.12,-0.22) -- (0.06,-0.31) -- (-0.04,-0.36) -- (0.02,-0.43) -- (0.12,-0.41) -- (0.03,-0.46) -- (0.13,-0.54) -- (0.07,-0.62) -- (0.21,-0.74) -- (0.15,-0.76) -- (0.20,-0.78) -- (0.19,-1.00) -- (0.30,-1.11) -- (0.20,-1.17) -- (0.22,-1.56) -- (0.12,-1.66) -- (0.17,-1.73) -- (0.09,-1.78) -- (0.17,-1.81) -- (0.10,-2.01) -- (0.01,-1.95) -- (-0.15,-1.99) -- (-0.24,-2.17) -- (-0.28,-2.06) -- (-0.46,-2.02) -- (-0.48,-1.83) -- (-0.57,-1.96) -- (-0.48,-2.04) -- (-0.73,-2.07) -- (-0.68,-2.22) -- (-0.74,-2.20) -- (-0.74,-0.09) -- (-0.60,0.03) -- (-0.53,0.01) -- (-0.57,0.07) -- (-0.47,0.14) -- cycle;
\draw[black!20,line width=0.3pt] (0.34,-2.86) -- (0.25,-2.99) -- (0.37,-2.97) -- (0.47,-3.22) -- (0.52,-3.17) -- (0.59,-3.28) -- (0.75,-3.32) -- (0.77,-3.52) -- (0.69,-3.77) -- (0.74,-3.84) -- (0.95,-3.74) -- (1.31,-3.75) -- (1.02,-4.00) -- (1.06,-4.16) -- (0.97,-4.25) -- (1.00,-4.31) -- (0.92,-4.32) -- (0.94,-4.38) -- (0.86,-4.36) -- (0.85,-4.47) -- (0.72,-4.46) -- (0.79,-4.60) -- (0.86,-4.55) -- (0.83,-4.65) -- (0.89,-4.70) -- (0.77,-4.78) -- (0.89,-4.79) -- (0.83,-4.86) -- (0.97,-4.96) -- (0.90,-4.96) -- (0.92,-5.06) -- (0.86,-5.08) -- (0.92,-5.14) -- (0.74,-5.10) -- (0.84,-5.19) -- (0.83,-5.26) -- (0.89,-5.27) -- (0.87,-5.20) -- (0.98,-5.25) -- (0.89,-5.32) -- (1.00,-5.38) -- (0.89,-5.37) -- (0.96,-5.46) -- (0.74,-5.43) -- (0.46,-5.53) -- (0.29,-5.45) -- (0.06,-5.46) -- (0.07,-5.55) -- (-0.14,-5.61) -- (-0.74,-5.03) -- (-0.74,-4.49) -- (-0.68,-4.50) -- (-0.70,-4.43) -- (-0.57,-4.23) -- (-0.04,-4.33) -- (-0.08,-4.09) -- (0.03,-3.99) -- (0.03,-3.74) -- (-0.22,-3.59) -- (-0.03,-3.38) -- (-0.06,-3.31) -- (0.00,-3.32) -- (0.04,-3.25) -- (-0.06,-3.26) -- (-0.09,-3.17) -- (0.03,-3.07) -- (-0.02,-2.99) -- (0.05,-3.02) -- (-0.00,-2.97) -- (0.02,-2.86) -- (0.13,-2.92) -- (0.14,-2.82) -- (0.10,-2.83) -- (0.11,-2.80) -- (0.17,-2.83) -- (0.20,-2.70) -- cycle;
\draw[black!20,line width=0.3pt] (9.83,-7.37) -- (10.88,-7.82) -- (10.93,-7.77) -- (11.11,-8.05) -- (11.27,-8.03) -- (11.31,-8.16) -- (11.44,-8.17) -- (11.41,-8.21) -- (11.63,-8.30) -- (11.62,-8.36) -- (11.66,-8.31) -- (11.74,-8.60) -- (11.74,-9.07) -- (11.64,-8.99) -- (11.38,-9.30) -- (11.22,-9.16) -- (10.99,-9.28) -- (10.42,-9.29) -- (10.29,-9.36) -- (10.29,-9.61) -- (10.22,-9.73) -- (10.17,-9.60) -- (9.99,-9.50) -- (9.92,-9.34) -- (9.96,-9.25) -- (9.85,-9.11) -- (10.03,-8.98) -- (10.00,-8.87) -- (9.72,-8.80) -- (9.48,-8.60) -- (9.44,-8.40) -- (9.49,-8.31) -- (9.35,-8.05) -- (9.29,-8.05) -- (9.26,-7.87) -- (9.39,-7.74) -- (9.64,-7.70) -- (9.59,-7.59) -- (9.68,-7.46) -- (9.80,-7.47) -- cycle;
\draw[black!20,line width=0.3pt] (1.82,-7.62) -- (1.96,-7.62) -- (2.04,-7.75) -- (1.87,-7.76) -- (1.68,-7.99) -- (1.70,-8.28) -- (1.79,-8.35) -- (1.82,-8.49) -- (2.03,-8.52) -- (2.20,-8.84) -- (2.10,-8.93) -- (2.26,-9.07) -- (2.20,-9.14) -- (2.25,-9.34) -- (2.16,-9.32) -- (2.09,-9.46) -- (2.23,-9.46) -- (2.25,-9.55) -- (2.13,-9.69) -- (2.20,-9.73) -- (2.06,-9.99) -- (2.05,-10.19) -- (2.10,-10.20) -- (2.12,-10.32) -- (1.07,-10.32) -- (1.02,-10.11) -- (0.80,-10.13) -- (0.77,-10.07) -- (0.73,-9.88) -- (0.80,-9.69) -- (0.74,-9.65) -- (0.59,-9.81) -- (0.55,-9.73) -- (0.38,-9.68) -- (0.38,-9.61) -- (0.26,-9.61) -- (0.35,-9.52) -- (0.27,-9.35) -- (0.20,-9.32) -- (0.21,-8.89) -- (0.44,-8.87) -- (0.51,-8.77) -- (0.45,-8.72) -- (0.54,-8.74) -- (0.58,-8.66) -- (0.62,-8.70) -- (0.59,-8.58) -- (0.76,-8.56) -- (0.74,-8.50) -- (0.82,-8.49) -- (0.81,-8.43) -- (0.89,-8.45) -- (0.88,-8.29) -- (1.02,-8.24) -- (0.95,-8.15) -- (1.05,-8.17) -- (1.12,-7.91) -- (1.28,-7.76) -- (1.23,-7.61) -- (1.27,-7.52) -- (1.35,-7.49) -- (1.50,-7.61) -- (1.67,-7.43) -- cycle;
\draw[black!20,line width=0.3pt] (1.26,-1.77) -- (1.26,-1.72) -- (1.52,-1.85) -- (1.40,-1.73) -- (1.52,-1.71) -- (1.56,-1.79) -- (1.71,-1.82) -- (1.67,-2.02) -- (1.72,-1.98) -- (1.80,-2.16) -- (1.92,-2.16) -- (2.01,-2.24) -- (1.99,-2.46) -- (2.05,-2.62) -- (2.20,-2.60) -- (2.29,-2.79) -- (2.06,-3.01) -- (2.02,-2.97) -- (2.06,-2.91) -- (1.94,-2.91) -- (1.96,-2.85) -- (1.86,-2.88) -- (1.87,-2.94) -- (1.71,-2.87) -- (1.72,-3.06) -- (1.53,-2.93) -- (1.51,-2.99) -- (1.44,-2.93) -- (1.43,-3.02) -- (1.33,-2.97) -- (1.33,-3.06) -- (1.26,-3.02) -- (1.30,-2.95) -- (1.24,-2.85) -- (1.14,-2.94) -- (1.14,-2.88) -- (0.99,-2.86) -- (1.03,-3.01) -- (0.97,-3.03) -- (1.03,-3.08) -- (0.97,-3.11) -- (1.07,-3.18) -- (1.02,-3.23) -- (1.09,-3.27) -- (1.13,-3.46) -- (1.06,-3.51) -- (1.17,-3.70) -- (1.31,-3.75) -- (0.95,-3.74) -- (0.74,-3.84) -- (0.69,-3.77) -- (0.77,-3.52) -- (0.75,-3.32) -- (0.59,-3.28) -- (0.52,-3.17) -- (0.47,-3.22) -- (0.37,-2.97) -- (0.25,-2.99) -- (0.34,-2.86) -- (0.20,-2.70) -- (0.17,-2.83) -- (0.03,-2.82) -- (0.05,-2.70) -- (-0.04,-2.66) -- (0.14,-2.55) -- (-0.05,-2.46) -- (0.08,-2.39) -- (0.06,-2.26) -- (0.21,-2.31) -- (0.21,-2.18) -- (0.25,-2.21) -- (0.23,-2.37) -- (0.31,-2.35) -- (0.32,-2.41) -- (0.42,-2.34) -- (0.30,-2.09) -- (0.44,-2.04) -- (0.44,-2.25) -- (0.57,-2.14) -- (0.50,-2.06) -- (0.56,-2.05) -- (0.52,-1.92) -- (0.83,-2.04) -- (0.67,-1.90) -- (0.75,-1.86) -- (0.77,-1.71) -- (0.87,-1.85) -- (0.97,-1.82) -- (1.01,-1.92) -- (1.14,-1.97) -- (1.15,-1.92) -- (1.27,-1.91) -- (1.09,-1.83) -- (1.17,-1.81) -- (1.11,-1.72) -- (1.17,-1.67) -- cycle;
\draw[black!20,line width=0.3pt] (2.07,0.79) -- (2.17,0.71) -- (2.17,0.58) -- (2.28,0.44) -- (2.17,0.34) -- (2.19,0.21) -- (2.00,-0.07) -- (2.00,-0.19) -- (1.92,-0.17) -- (1.97,-0.26) -- (1.76,-0.21) -- (1.73,-0.32) -- (1.65,-0.33) -- (1.59,-0.29) -- (1.66,-0.25) -- (1.53,-0.18) -- (1.59,-0.16) -- (1.40,-0.14) -- (1.38,-0.03) -- (1.31,-0.04) -- (1.32,-0.14) -- (1.16,-0.12) -- (1.22,-0.21) -- (1.07,-0.14) -- (0.99,-0.22) -- (0.88,-0.20) -- (0.89,-0.04) -- (0.82,-0.11) -- (0.77,-0.06) -- (0.83,0.04) -- (0.69,0.10) -- (0.83,0.16) -- (0.61,0.17) -- (0.56,0.24) -- (0.58,0.07) -- (0.35,-0.03) -- (0.21,0.05) -- (0.30,0.10) -- (0.13,0.17) -- (0.16,0.31) -- (0.30,0.34) -- (0.29,0.45) -- (0.46,0.46) -- (0.45,0.55) -- (0.54,0.50) -- (0.63,0.68) -- (0.78,0.66) -- (0.96,0.82) -- cycle;
\draw[black!20,line width=0.3pt] (-0.07,0.82) -- (-0.01,0.55) -- (-0.07,0.57) -- (-0.10,0.48) -- (-0.03,0.48) -- (-0.04,0.37) -- (-0.17,0.29) -- (-0.05,0.23) -- (-0.17,0.17) -- (-0.09,0.12) -- (-0.17,0.11) -- (-0.16,0.04) -- (-0.36,0.07) -- (-0.38,0.02) -- (-0.47,0.14) -- (-0.57,0.07) -- (-0.53,0.01) -- (-0.60,0.03) -- (-0.73,-0.11) -- (-0.74,0.48) -- (-0.66,0.68) -- (-0.44,0.82) -- cycle;
\draw[black!20,line width=0.3pt] (-0.68,-2.22) -- (-0.73,-2.07) -- (-0.48,-2.04) -- (-0.57,-1.96) -- (-0.48,-1.83) -- (-0.46,-2.02) -- (-0.28,-2.06) -- (-0.24,-2.17) -- (-0.15,-1.99) -- (0.01,-1.95) -- (-0.12,-2.11) -- (0.01,-2.07) -- (0.10,-2.32) -- (-0.05,-2.46) -- (0.14,-2.55) -- (-0.04,-2.66) -- (0.05,-2.70) -- (0.03,-2.82) -- (0.11,-2.80) -- (0.15,-2.87) -- (0.13,-2.92) -- (0.02,-2.86) -- (-0.00,-2.97) -- (0.05,-3.02) -- (-0.02,-2.99) -- (0.03,-3.07) -- (-0.09,-3.17) -- (-0.06,-3.26) -- (0.04,-3.25) -- (0.00,-3.32) -- (-0.06,-3.31) -- (-0.03,-3.38) -- (-0.22,-3.59) -- (0.03,-3.74) -- (0.03,-3.99) -- (-0.08,-4.09) -- (-0.04,-4.33) -- (-0.57,-4.23) -- (-0.65,-4.30) -- (-0.71,-4.52) -- (-0.74,-2.20) -- cycle;
\draw[black!20,line width=0.3pt] (-0.73,-4.57) -- (-0.74,-4.58) -- (-0.74,-4.55) -- cycle;
\draw[black!20,line width=0.3pt] (-0.74,-4.61) -- (-0.74,-4.68) -- (-0.74,-4.68) -- cycle;
\draw[black!20,line width=0.3pt] (0.05,0.17) -- (0.12,0.29) -- (0.13,0.17) -- (0.30,0.10) -- (0.21,0.05) -- (0.25,-0.02) -- (0.35,-0.03) -- (0.58,0.07) -- (0.56,0.24) -- (0.61,0.17) -- (0.83,0.16) -- (0.69,0.10) -- (0.83,0.04) -- (0.77,-0.06) -- (0.82,-0.11) -- (0.89,-0.04) -- (0.88,-0.20) -- (0.99,-0.22) -- (1.07,-0.14) -- (1.22,-0.21) -- (1.16,-0.12) -- (1.32,-0.14) -- (1.31,-0.04) -- (1.38,-0.03) -- (1.40,-0.14) -- (1.59,-0.16) -- (1.53,-0.18) -- (1.66,-0.25) -- (1.59,-0.29) -- (1.65,-0.33) -- (1.73,-0.32) -- (1.76,-0.21) -- (1.87,-0.26) -- (1.98,-0.24) -- (1.92,-0.17) -- (2.03,-0.21) -- (2.11,-0.26) -- (2.15,-0.44) -- (2.33,-0.45) -- (2.33,-0.50) -- (2.42,-0.46) -- (2.40,-0.72) -- (2.52,-0.91) -- (2.15,-0.98) -- (1.99,-1.19) -- (1.80,-1.26) -- (1.85,-1.32) -- (1.79,-1.34) -- (1.78,-1.66) -- (1.65,-1.56) -- (1.59,-1.65) -- (1.52,-1.63) -- (1.52,-1.71) -- (1.40,-1.73) -- (1.52,-1.85) -- (1.26,-1.72) -- (1.26,-1.77) -- (1.17,-1.67) -- (1.11,-1.72) -- (1.17,-1.81) -- (1.09,-1.83) -- (1.27,-1.91) -- (1.15,-1.92) -- (1.14,-1.97) -- (1.01,-1.92) -- (0.97,-1.82) -- (0.87,-1.85) -- (0.77,-1.71) -- (0.75,-1.86) -- (0.67,-1.90) -- (0.83,-2.04) -- (0.52,-1.92) -- (0.56,-2.05) -- (0.50,-2.06) -- (0.57,-2.14) -- (0.44,-2.25) -- (0.44,-2.04) -- (0.30,-2.09) -- (0.42,-2.34) -- (0.32,-2.41) -- (0.31,-2.35) -- (0.23,-2.37) -- (0.24,-2.18) -- (0.18,-2.24) -- (0.21,-2.31) -- (0.06,-2.26) -- (0.01,-2.07) -- (-0.12,-2.11) -- (0.14,-1.96) -- (0.17,-1.81) -- (0.09,-1.78) -- (0.17,-1.73) -- (0.12,-1.66) -- (0.22,-1.56) -- (0.20,-1.17) -- (0.30,-1.11) -- (0.19,-1.00) -- (0.20,-0.78) -- (0.15,-0.76) -- (0.21,-0.74) -- (0.07,-0.62) -- (0.13,-0.54) -- (0.03,-0.46) -- (0.12,-0.41) -- (0.02,-0.43) -- (-0.04,-0.32) -- (0.06,-0.31) -- (-0.12,-0.22) -- (-0.07,-0.14) -- (-0.20,0.03) -- (-0.27,-0.00) -- (-0.09,0.12) -- (-0.17,0.17) -- (-0.05,0.23) -- (-0.15,0.24) -- (-0.11,0.37) -- cycle;
\draw[black!20,line width=0.3pt] (10.21,-10.00) -- (10.48,-10.10) -- (10.76,-9.82) -- (10.86,-9.82) -- (11.17,-10.32) -- (9.78,-10.32) -- (10.22,-9.73) -- cycle;
\draw[black!20,line width=0.3pt] (11.69,-9.08) -- (11.74,-9.07) -- (11.74,-9.98) -- (11.56,-10.28) -- (11.17,-10.32) -- (10.86,-9.82) -- (10.76,-9.82) -- (10.48,-10.10) -- (10.21,-10.00) -- (10.20,-9.80) -- (10.34,-9.33) -- (10.75,-9.24) -- (10.99,-9.28) -- (11.22,-9.16) -- (11.38,-9.30) -- (11.64,-8.99) -- (11.70,-9.01) -- cycle;
\draw[black!20,line width=0.3pt] (6.35,0.79) -- (6.08,0.76) -- (6.04,0.34) -- (5.95,0.29) -- (5.90,0.12) -- (5.79,0.20) -- (5.78,0.09) -- (5.94,0.01) -- (5.86,-0.08) -- (5.94,-0.16) -- (5.75,-0.24) -- (5.73,-0.31) -- (5.66,-0.28) -- (5.66,-0.41) -- (5.63,-0.36) -- (5.49,-0.50) -- (5.45,-0.46) -- (5.40,-0.64) -- (5.27,-0.64) -- (5.28,-0.77) -- (5.21,-0.85) -- (4.76,-1.11) -- (4.73,-1.02) -- (4.48,-1.00) -- (4.40,-0.89) -- (4.27,-0.93) -- (4.03,-0.72) -- (3.84,-0.69) -- (3.68,-0.45) -- (3.64,-0.21) -- (3.44,-0.15) -- (3.33,-0.01) -- (3.18,-0.02) -- (3.18,-0.09) -- (3.12,-0.06) -- (3.12,-0.15) -- (3.05,-0.07) -- (2.75,-0.21) -- (2.78,-0.08) -- (2.58,-0.07) -- (2.55,-0.14) -- (2.51,-0.09) -- (2.52,-0.26) -- (2.49,-0.19) -- (2.34,-0.27) -- (2.36,-0.23) -- (2.27,-0.14) -- (2.14,-0.22) -- (1.98,-0.19) -- (2.09,0.10) -- (2.19,0.21) -- (2.17,0.34) -- (2.28,0.44) -- (2.17,0.58) -- (2.10,0.82) -- cycle;
\draw[black!20,line width=0.3pt] (0.89,-5.52) -- (0.95,-5.48) -- (0.97,-5.64) -- (1.04,-5.63) -- (1.07,-5.71) -- (0.94,-5.76) -- (0.97,-5.83) -- (0.90,-5.90) -- (0.96,-5.97) -- (0.89,-6.03) -- (0.94,-6.15) -- (0.84,-6.40) -- (0.91,-6.41) -- (0.71,-6.41) -- (0.69,-6.59) -- (0.74,-6.66) -- (0.57,-6.75) -- (0.54,-6.66) -- (0.38,-6.65) -- (0.25,-6.80) -- (0.11,-6.70) -- (0.00,-6.81) -- (0.16,-6.97) -- (0.13,-7.07) -- (0.02,-7.06) -- (-0.21,-6.87) -- (-0.36,-6.91) -- (-0.47,-6.85) -- (-0.57,-6.92) -- (-0.57,-6.78) -- (-0.74,-6.69) -- (-0.66,-6.50) -- (-0.70,-6.45) -- (-0.41,-6.34) -- (-0.27,-6.04) -- (-0.20,-6.03) -- (-0.18,-5.86) -- (0.00,-5.79) -- (-0.17,-5.61) -- (0.07,-5.55) -- (0.06,-5.46) -- (0.23,-5.45) -- (0.46,-5.53) -- (0.74,-5.43) -- cycle;
\draw[black!20,line width=0.3pt] (11.40,-6.01) -- (11.74,-6.09) -- (11.74,-6.97) -- (11.66,-7.04) -- (11.65,-6.97) -- (11.48,-6.96) -- (11.45,-6.88) -- (11.55,-6.82) -- (11.43,-6.75) -- (11.28,-6.85) -- (11.19,-6.74) -- (11.03,-6.74) -- (10.98,-6.26) -- (11.08,-6.28) -- (11.11,-6.11) -- (11.26,-6.05) -- (11.24,-5.99) -- (11.36,-5.94) -- cycle;
\draw[black!20,line width=0.3pt] (11.73,-7.39) -- (11.74,-7.36) -- (11.74,-7.40) -- cycle;
\draw[black!20,line width=0.3pt] (11.73,-7.32) -- (11.74,-7.32) -- (11.74,-7.32) -- cycle;
\draw[black!20,line width=0.3pt] (11.74,-7.22) -- (11.68,-7.31) -- (11.71,-7.22) -- cycle;
\draw[black!20,line width=0.3pt] (11.73,-7.14) -- (11.74,-7.14) -- (11.74,-7.21) -- cycle;
\draw[black!20,line width=0.3pt] (11.23,0.80) -- (11.22,0.82) -- (11.23,0.82) -- cycle;
\draw[black!20,line width=0.3pt] (11.17,0.77) -- (11.09,0.78) -- (11.07,0.54) -- (10.99,0.52) -- (11.09,0.33) -- (11.04,0.06) -- (10.97,-0.00) -- (10.80,0.04) -- (10.65,-0.21) -- (10.47,-0.28) -- (10.37,-0.59) -- (10.15,-0.64) -- (10.10,-0.77) -- (9.89,-0.67) -- (9.87,-0.72) -- (9.43,-0.83) -- (9.08,-0.47) -- (9.10,0.46) -- (8.98,0.82) -- (11.13,0.82) -- cycle;
\draw[black!20,line width=0.3pt] (10.23,-6.37) -- (10.56,-6.30) -- (10.68,-6.34) -- (10.71,-6.47) -- (10.89,-6.47) -- (10.85,-6.38) -- (10.96,-6.32) -- (11.03,-6.74) -- (11.19,-6.74) -- (11.28,-6.85) -- (11.43,-6.75) -- (11.55,-6.82) -- (11.45,-6.88) -- (11.48,-6.96) -- (11.65,-6.97) -- (11.66,-7.04) -- (11.74,-6.97) -- (11.73,-7.27) -- (11.68,-7.31) -- (11.74,-7.31) -- (11.74,-8.36) -- (11.66,-8.31) -- (11.62,-8.36) -- (11.63,-8.30) -- (11.41,-8.21) -- (11.44,-8.17) -- (11.31,-8.16) -- (11.27,-8.03) -- (11.11,-8.05) -- (10.93,-7.77) -- (10.88,-7.82) -- (9.91,-7.38) -- (9.82,-6.95) -- (9.71,-6.86) -- (9.69,-6.69) -- (9.60,-6.66) -- (9.65,-6.46) -- (9.56,-6.45) -- (9.50,-6.27) -- (9.69,-6.24) -- (9.90,-6.36) -- (10.04,-6.35) -- (10.17,-6.21) -- cycle;
\draw[black!20,line width=0.3pt] (-0.74,-5.03) -- (0.00,-5.79) -- (-0.18,-5.86) -- (-0.20,-6.03) -- (-0.27,-6.04) -- (-0.41,-6.34) -- (-0.70,-6.45) -- (-0.66,-6.50) -- (-0.70,-6.64) -- (-0.74,-6.64) -- cycle;
\draw[black!20,line width=0.3pt] (-0.72,-6.97) -- (-0.74,-6.99) -- (-0.74,-6.90) -- cycle;
\draw[black!20,line width=0.3pt] (0.96,0.82) -- (0.78,0.66) -- (0.63,0.68) -- (0.54,0.50) -- (0.45,0.55) -- (0.46,0.46) -- (0.29,0.45) -- (0.30,0.34) -- (0.15,0.25) -- (0.12,0.29) -- (0.05,0.17) -- (-0.08,0.30) -- (-0.03,0.48) -- (-0.10,0.51) -- (-0.07,0.57) -- (-0.01,0.55) -- (-0.07,0.64) -- (-0.03,0.71) -- (-0.07,0.82) -- cycle;
\draw[black!20,line width=0.3pt] (3.34,-0.03) -- (3.44,-0.15) -- (3.64,-0.21) -- (3.68,-0.45) -- (3.84,-0.69) -- (4.03,-0.72) -- (4.55,-1.24) -- (4.51,-1.55) -- (4.29,-1.61) -- (4.30,-1.74) -- (4.42,-1.71) -- (4.64,-1.79) -- (4.75,-1.93) -- (5.20,-1.97) -- (5.05,-2.30) -- (4.85,-2.18) -- (4.38,-2.13) -- (4.32,-2.02) -- (3.98,-1.94) -- (4.00,-1.90) -- (3.72,-1.75) -- (3.75,-1.82) -- (3.67,-1.81) -- (3.69,-1.87) -- (3.51,-1.83) -- (3.54,-1.93) -- (3.37,-1.96) -- (3.39,-2.06) -- (3.29,-2.05) -- (3.22,-1.93) -- (3.15,-2.01) -- (3.11,-1.86) -- (2.97,-2.01) -- (2.97,-2.10) -- (2.91,-2.07) -- (2.92,-1.95) -- (2.83,-1.97) -- (2.91,-1.80) -- (2.78,-1.78) -- (2.74,-1.90) -- (2.66,-1.81) -- (2.63,-1.86) -- (2.56,-1.82) -- (2.62,-1.76) -- (2.48,-1.76) -- (2.57,-1.66) -- (2.27,-1.82) -- (2.30,-1.90) -- (2.22,-1.97) -- (2.26,-1.88) -- (2.21,-1.74) -- (2.13,-1.75) -- (2.14,-1.83) -- (2.09,-1.74) -- (1.96,-1.71) -- (2.07,-1.55) -- (1.91,-1.59) -- (1.92,-1.66) -- (1.88,-1.58) -- (1.78,-1.61) -- (1.79,-1.34) -- (1.85,-1.32) -- (1.80,-1.26) -- (1.99,-1.19) -- (2.15,-0.98) -- (2.52,-0.91) -- (2.40,-0.72) -- (2.42,-0.46) -- (2.33,-0.50) -- (2.33,-0.45) -- (2.15,-0.44) -- (2.07,-0.23) -- (2.27,-0.14) -- (2.36,-0.23) -- (2.34,-0.27) -- (2.49,-0.19) -- (2.52,-0.26) -- (2.51,-0.09) -- (2.55,-0.14) -- (2.58,-0.07) -- (2.78,-0.08) -- (2.75,-0.21) -- (3.05,-0.07) -- (3.12,-0.15) -- (3.12,-0.06) -- (3.21,-0.07) -- (3.18,-0.02) -- cycle;
\draw[black!20,line width=0.3pt] (11.69,-0.77) -- (11.33,-0.73) -- (11.22,-0.59) -- (10.99,-0.76) -- (10.74,-0.69) -- (10.61,-0.80) -- (10.37,-0.59) -- (10.47,-0.28) -- (10.65,-0.21) -- (10.80,0.04) -- (10.97,-0.00) -- (11.04,0.06) -- (11.09,0.33) -- (10.99,0.52) -- (11.07,0.54) -- (11.09,0.78) -- (11.17,0.77) -- (11.13,0.82) -- (11.50,0.80) -- (11.72,0.59) -- (11.74,-0.71) -- cycle;
\draw[black!20,line width=0.3pt] (10.03,-8.98) -- (9.85,-9.11) -- (9.96,-9.25) -- (9.92,-9.34) -- (9.99,-9.50) -- (10.17,-9.60) -- (10.22,-9.73) -- (9.78,-10.32) -- (7.62,-10.32) -- (7.77,-10.06) -- (7.89,-10.02) -- (7.88,-9.88) -- (7.98,-9.83) -- (8.02,-9.62) -- (8.19,-9.68) -- (8.23,-9.56) -- (8.36,-9.55) -- (8.38,-9.48) -- (8.30,-9.45) -- (8.40,-9.39) -- (8.44,-9.38) -- (8.45,-9.51) -- (8.48,-9.46) -- (8.52,-9.54) -- (8.69,-9.56) -- (8.85,-9.33) -- (9.05,-9.41) -- (9.08,-9.34) -- (9.18,-9.38) -- (9.19,-9.32) -- (9.26,-9.32) -- (9.34,-9.03) -- (9.53,-8.96) -- (9.63,-9.00) -- (9.74,-8.92) -- (9.92,-8.96) -- (9.95,-8.86) -- cycle;
\draw[black!20,line width=0.3pt] (1.30,-3.81) -- (1.34,-4.14) -- (1.48,-4.20) -- (1.41,-4.31) -- (1.59,-4.47) -- (1.54,-4.50) -- (1.61,-4.52) -- (1.59,-4.62) -- (1.69,-4.83) -- (1.64,-4.88) -- (1.59,-4.82) -- (1.60,-5.04) -- (1.54,-5.04) -- (1.53,-5.15) -- (1.40,-5.24) -- (1.46,-5.29) -- (1.51,-5.25) -- (1.51,-5.51) -- (1.42,-5.48) -- (1.44,-5.44) -- (1.32,-5.46) -- (1.34,-5.56) -- (1.26,-5.53) -- (1.23,-5.69) -- (1.16,-5.72) -- (1.20,-5.82) -- (1.27,-5.81) -- (1.25,-5.91) -- (1.21,-5.96) -- (1.15,-5.87) -- (1.09,-5.94) -- (1.18,-6.03) -- (1.04,-6.04) -- (1.04,-6.14) -- (0.94,-6.15) -- (0.89,-6.03) -- (0.96,-5.97) -- (0.90,-5.90) -- (0.97,-5.83) -- (0.94,-5.76) -- (1.07,-5.71) -- (1.04,-5.63) -- (0.97,-5.64) -- (0.95,-5.48) -- (0.89,-5.52) -- (0.78,-5.45) -- (0.96,-5.46) -- (0.89,-5.37) -- (1.00,-5.38) -- (0.89,-5.32) -- (0.98,-5.25) -- (0.87,-5.20) -- (0.89,-5.27) -- (0.83,-5.26) -- (0.84,-5.19) -- (0.74,-5.10) -- (0.92,-5.14) -- (0.86,-5.08) -- (0.92,-5.06) -- (0.90,-4.96) -- (0.97,-4.96) -- (0.83,-4.86) -- (0.89,-4.79) -- (0.77,-4.78) -- (0.89,-4.70) -- (0.83,-4.65) -- (0.86,-4.55) -- (0.79,-4.60) -- (0.72,-4.46) -- (0.85,-4.47) -- (0.86,-4.36) -- (0.94,-4.38) -- (0.92,-4.32) -- (1.00,-4.31) -- (0.97,-4.25) -- (1.06,-4.16) -- (1.02,-4.00) -- cycle;
\end{scope}
\filldraw[fill=primary!30,draw=black!45,line width=0.4pt] (5.14,-2.53) -- (5.44,-2.78) -- (5.40,-2.94) -- (5.33,-3.01) -- (5.76,-3.16) -- (5.87,-3.06) -- (6.00,-3.05) -- (6.01,-3.09) -- (6.42,-2.93) -- (6.44,-2.99) -- (6.65,-3.01) -- (6.74,-2.85) -- (6.78,-2.92) -- (6.85,-2.92) -- (6.81,-2.84) -- (6.82,-2.79) -- (6.86,-2.80) -- (6.82,-2.72) -- (6.90,-2.76) -- (6.89,-2.69) -- (7.55,-2.65) -- (7.72,-2.78) -- (8.02,-2.85) -- (7.97,-3.00) -- (8.07,-2.98) -- (8.16,-3.04) -- (8.15,-3.12) -- (8.35,-3.16) -- (8.24,-3.23) -- (8.18,-3.37) -- (8.14,-3.32) -- (8.08,-3.37) -- (8.14,-3.40) -- (8.16,-3.48) -- (8.26,-3.42) -- (8.47,-3.57) -- (8.20,-3.54) -- (8.15,-3.58) -- (8.33,-3.71) -- (8.24,-3.75) -- (8.28,-3.82) -- (8.26,-3.87) -- (8.19,-3.81) -- (8.17,-3.84) -- (8.18,-4.07) -- (8.12,-4.00) -- (8.12,-4.07) -- (8.05,-4.03) -- (8.00,-4.06) -- (8.12,-4.11) -- (8.21,-4.23) -- (8.16,-4.25) -- (8.22,-4.28) -- (8.22,-4.32) -- (8.29,-4.24) -- (8.38,-4.30) -- (8.39,-4.46) -- (8.50,-4.48) -- (8.43,-4.39) -- (8.54,-4.27) -- (8.58,-4.43) -- (8.75,-4.45) -- (8.61,-4.44) -- (8.62,-4.59) -- (8.45,-4.62) -- (8.48,-4.63) -- (8.46,-4.69) -- (8.59,-4.66) -- (8.56,-4.71) -- (8.60,-4.78) -- (8.51,-4.78) -- (8.39,-4.93) -- (8.50,-4.87) -- (8.58,-4.87) -- (8.57,-4.92) -- (8.67,-4.91) -- (8.60,-5.01) -- (8.53,-5.00) -- (8.54,-5.06) -- (8.49,-5.11) -- (8.68,-5.07) -- (8.69,-5.23) -- (8.74,-5.17) -- (8.77,-5.22) -- (8.80,-5.02) -- (8.88,-5.05) -- (8.85,-5.13) -- (8.92,-5.17) -- (8.90,-5.23) -- (8.85,-5.22) -- (8.77,-5.38) -- (8.67,-5.39) -- (8.66,-5.43) -- (8.76,-5.44) -- (8.76,-5.50) -- (8.82,-5.51) -- (8.83,-5.55) -- (8.65,-5.70) -- (8.59,-5.54) -- (8.61,-5.40) -- (8.50,-5.46) -- (8.39,-5.35) -- (8.34,-5.41) -- (8.27,-5.41) -- (8.20,-5.47) -- (8.13,-5.34) -- (8.08,-5.35) -- (8.08,-5.40) -- (8.00,-5.29) -- (7.94,-5.37) -- (7.88,-5.34) -- (7.55,-5.41) -- (7.49,-5.30) -- (7.49,-5.35) -- (7.35,-5.38) -- (7.34,-5.30) -- (7.28,-5.33) -- (7.19,-5.23) -- (7.17,-5.28) -- (7.04,-5.30) -- (7.04,-5.38) -- (7.11,-5.47) -- (7.27,-5.54) -- (7.28,-5.64) -- (7.19,-5.68) -- (7.26,-5.78) -- (7.38,-5.81) -- (7.55,-5.98) -- (7.71,-5.94) -- (7.70,-6.26) -- (7.56,-6.47) -- (7.57,-6.57) -- (7.53,-6.60) -- (7.41,-6.52) -- (7.38,-6.56) -- (7.35,-6.51) -- (7.26,-6.56) -- (7.06,-6.53) -- (6.97,-6.64) -- (6.84,-6.51) -- (6.77,-6.60) -- (6.80,-6.65) -- (6.74,-6.60) -- (6.67,-6.63) -- (6.69,-6.71) -- (6.65,-6.77) -- (6.57,-6.78) -- (6.52,-6.73) -- (6.35,-6.78) -- (6.37,-6.81) -- (6.33,-6.84) -- (6.29,-6.78) -- (6.26,-6.83) -- (6.28,-6.93) -- (6.11,-6.87) -- (6.09,-6.78) -- (6.17,-6.72) -- (6.17,-6.55) -- (6.32,-6.55) -- (6.24,-6.38) -- (6.21,-6.41) -- (6.24,-6.45) -- (6.17,-6.51) -- (6.07,-6.47) -- (6.07,-6.41) -- (5.99,-6.44) -- (6.01,-6.52) -- (5.92,-6.45) -- (5.82,-6.47) -- (5.75,-6.51) -- (5.74,-6.59) -- (5.69,-6.58) -- (5.60,-6.70) -- (5.55,-6.68) -- (5.54,-6.76) -- (5.47,-6.66) -- (5.43,-6.42) -- (5.31,-6.35) -- (5.31,-6.28) -- (5.23,-6.27) -- (5.20,-6.33) -- (5.23,-6.40) -- (5.18,-6.44) -- (5.11,-6.40) -- (5.12,-6.32) -- (5.07,-6.35) -- (4.78,-6.27) -- (4.75,-6.36) -- (4.78,-6.43) -- (4.66,-6.37) -- (4.55,-6.42) -- (4.53,-6.34) -- (4.44,-6.30) -- (4.47,-6.39) -- (4.36,-6.41) -- (4.27,-6.34) -- (4.32,-6.24) -- (4.17,-6.32) -- (4.15,-6.24) -- (4.20,-6.20) -- (4.18,-6.12) -- (4.05,-6.00) -- (4.01,-5.84) -- (4.10,-5.90) -- (4.20,-5.89) -- (4.21,-5.81) -- (4.30,-5.70) -- (4.21,-5.63) -- (4.28,-5.60) -- (4.28,-5.51) -- (4.34,-5.54) -- (4.32,-5.51) -- (4.45,-5.38) -- (4.45,-5.22) -- (4.33,-5.15) -- (4.40,-4.99) -- (4.38,-4.94) -- (4.24,-4.96) -- (4.22,-4.88) -- (4.42,-4.69) -- (4.37,-4.53) -- (4.49,-4.33) -- (4.55,-4.09) -- (4.41,-4.08) -- (4.13,-4.28) -- (4.00,-4.26) -- (3.88,-4.17) -- (3.72,-4.18) -- (3.60,-4.29) -- (3.22,-4.31) -- (3.00,-4.48) -- (2.91,-4.44) -- (2.73,-4.56) -- (2.69,-4.45) -- (2.73,-4.40) -- (2.63,-4.35) -- (2.69,-4.30) -- (2.69,-4.21) -- (2.77,-4.15) -- (2.86,-4.20) -- (2.81,-4.13) -- (2.86,-4.11) -- (2.94,-4.13) -- (2.93,-4.05) -- (2.98,-4.07) -- (3.01,-4.21) -- (3.07,-4.14) -- (3.16,-4.20) -- (3.21,-4.14) -- (3.20,-4.10) -- (3.25,-4.18) -- (3.30,-4.18) -- (3.28,-4.10) -- (3.33,-4.14) -- (3.39,-4.12) -- (3.39,-4.06) -- (3.44,-4.04) -- (3.38,-3.98) -- (3.47,-3.97) -- (3.49,-3.92) -- (3.58,-4.00) -- (3.72,-3.83) -- (3.64,-3.74) -- (3.56,-3.73) -- (3.66,-3.56) -- (3.64,-3.50) -- (4.22,-3.58) -- (4.50,-3.35) -- (4.74,-3.56) -- (4.87,-3.50) -- (5.07,-3.52) -- (5.15,-3.46) -- (5.11,-3.33) -- (5.21,-3.35) -- (5.23,-3.28) -- (5.10,-3.10) -- (5.12,-2.89) -- (4.96,-2.61) -- (5.01,-2.54) -- cycle;
\filldraw[fill=primary!10,draw=black!45,line width=0.4pt] (4.33,-6.25) -- (4.27,-6.34) -- (4.40,-6.43) -- (4.47,-6.39) -- (4.44,-6.30) -- (4.53,-6.34) -- (4.55,-6.42) -- (4.66,-6.37) -- (4.78,-6.43) -- (4.75,-6.36) -- (4.78,-6.27) -- (5.07,-6.35) -- (5.12,-6.32) -- (5.11,-6.40) -- (5.21,-6.43) -- (5.23,-6.27) -- (5.31,-6.28) -- (5.31,-6.35) -- (5.43,-6.42) -- (5.47,-6.66) -- (5.54,-6.76) -- (5.55,-6.68) -- (5.60,-6.70) -- (5.69,-6.58) -- (5.74,-6.59) -- (5.75,-6.51) -- (5.82,-6.47) -- (5.92,-6.45) -- (6.01,-6.52) -- (5.99,-6.44) -- (6.07,-6.41) -- (6.07,-6.47) -- (6.17,-6.51) -- (6.24,-6.45) -- (6.21,-6.41) -- (6.24,-6.38) -- (6.32,-6.55) -- (6.17,-6.55) -- (6.17,-6.72) -- (6.09,-6.78) -- (6.11,-6.87) -- (6.20,-6.93) -- (6.28,-6.93) -- (6.26,-6.83) -- (6.29,-6.78) -- (6.33,-6.84) -- (6.37,-6.81) -- (6.35,-6.78) -- (6.52,-6.73) -- (6.57,-6.78) -- (6.65,-6.77) -- (6.69,-6.71) -- (6.67,-6.63) -- (6.74,-6.60) -- (6.99,-6.73) -- (7.01,-6.81) -- (7.20,-6.97) -- (7.12,-7.17) -- (7.14,-7.21) -- (7.06,-7.33) -- (7.18,-7.61) -- (7.25,-7.63) -- (7.40,-7.76) -- (7.40,-7.84) -- (7.48,-7.93) -- (7.49,-8.06) -- (7.42,-8.00) -- (7.36,-8.01) -- (7.35,-8.16) -- (7.25,-8.23) -- (7.27,-8.29) -- (7.21,-8.28) -- (7.18,-8.39) -- (7.05,-8.44) -- (7.13,-8.50) -- (7.08,-8.55) -- (6.96,-8.52) -- (7.03,-8.69) -- (6.94,-8.69) -- (6.96,-8.81) -- (6.90,-8.84) -- (6.79,-8.81) -- (6.79,-8.90) -- (6.74,-8.93) -- (6.75,-8.75) -- (6.69,-8.71) -- (6.65,-8.75) -- (6.59,-8.74) -- (6.62,-8.86) -- (6.50,-8.94) -- (6.42,-8.88) -- (6.34,-8.88) -- (6.37,-8.79) -- (6.31,-8.72) -- (6.14,-8.89) -- (6.15,-8.97) -- (6.08,-8.95) -- (6.08,-9.03) -- (6.03,-9.07) -- (5.88,-9.05) -- (5.89,-9.13) -- (5.82,-9.21) -- (5.61,-9.13) -- (5.60,-9.29) -- (5.38,-9.29) -- (5.36,-9.33) -- (5.28,-9.34) -- (5.18,-9.29) -- (5.04,-9.38) -- (4.99,-9.36) -- (4.95,-9.26) -- (4.81,-9.20) -- (4.73,-9.10) -- (4.94,-8.79) -- (4.84,-8.67) -- (4.75,-8.70) -- (4.66,-8.79) -- (4.60,-8.98) -- (4.52,-9.09) -- (4.60,-9.21) -- (4.42,-9.37) -- (4.36,-9.39) -- (4.27,-9.29) -- (4.11,-9.29) -- (4.08,-9.25) -- (4.12,-9.21) -- (4.00,-9.08) -- (4.03,-9.03) -- (3.99,-8.96) -- (4.07,-8.92) -- (4.03,-8.85) -- (4.11,-8.68) -- (4.05,-8.35) -- (4.11,-8.26) -- (4.06,-8.17) -- (4.11,-8.13) -- (3.96,-7.86) -- (3.88,-7.79) -- (3.88,-7.57) -- (3.78,-7.49) -- (3.80,-7.41) -- (3.70,-7.40) -- (3.65,-7.45) -- (3.55,-7.21) -- (3.44,-7.34) -- (3.35,-7.27) -- (3.20,-7.33) -- (3.13,-7.43) -- (3.04,-7.36) -- (3.05,-7.41) -- (3.00,-7.45) -- (2.91,-7.50) -- (2.86,-7.46) -- (2.80,-7.52) -- (2.64,-7.51) -- (2.59,-7.34) -- (2.65,-7.30) -- (2.64,-7.19) -- (2.86,-7.21) -- (2.77,-7.07) -- (2.93,-7.05) -- (2.85,-6.94) -- (3.03,-6.76) -- (2.97,-6.73) -- (2.97,-6.59) -- (3.18,-6.71) -- (3.27,-6.55) -- (3.35,-6.57) -- (3.32,-6.76) -- (3.42,-6.78) -- (3.49,-6.61) -- (3.59,-6.66) -- (3.60,-6.45) -- (3.67,-6.37) -- (3.74,-6.36) -- (3.79,-6.40) -- (3.80,-6.56) -- (3.89,-6.60) -- (3.90,-6.40) -- (4.09,-6.40) -- (4.01,-6.30) -- (4.16,-6.29) -- (4.17,-6.32) -- (4.26,-6.24) -- cycle;
\filldraw[fill=primary!10,draw=black!45,line width=0.4pt] (4.50,-4.21) -- (4.49,-4.33) -- (4.37,-4.53) -- (4.42,-4.69) -- (4.22,-4.88) -- (4.24,-4.96) -- (4.38,-4.94) -- (4.40,-4.99) -- (4.33,-5.15) -- (4.38,-5.21) -- (4.43,-5.19) -- (4.47,-5.30) -- (4.32,-5.51) -- (4.34,-5.54) -- (4.28,-5.51) -- (4.28,-5.60) -- (4.21,-5.63) -- (4.30,-5.70) -- (4.21,-5.81) -- (4.20,-5.89) -- (4.10,-5.90) -- (4.01,-5.84) -- (4.05,-6.00) -- (4.18,-6.12) -- (4.20,-6.20) -- (4.15,-6.24) -- (4.16,-6.29) -- (4.01,-6.30) -- (4.08,-6.43) -- (3.99,-6.38) -- (3.90,-6.40) -- (3.89,-6.60) -- (3.80,-6.56) -- (3.81,-6.45) -- (3.74,-6.36) -- (3.67,-6.37) -- (3.60,-6.45) -- (3.59,-6.66) -- (3.49,-6.61) -- (3.42,-6.78) -- (3.32,-6.76) -- (3.35,-6.57) -- (3.27,-6.55) -- (3.18,-6.71) -- (2.97,-6.59) -- (2.91,-6.75) -- (2.82,-6.73) -- (2.81,-6.58) -- (2.46,-6.50) -- (2.42,-6.38) -- (2.28,-6.50) -- (2.13,-6.55) -- (1.94,-6.87) -- (1.79,-6.94) -- (1.59,-7.14) -- (1.12,-7.40) -- (1.02,-7.35) -- (0.99,-7.25) -- (0.79,-7.13) -- (0.70,-7.15) -- (0.61,-7.24) -- (0.55,-7.23) -- (0.71,-6.99) -- (0.70,-6.93) -- (0.64,-6.90) -- (0.54,-6.93) -- (0.37,-7.11) -- (0.13,-7.07) -- (0.16,-6.97) -- (0.00,-6.81) -- (0.11,-6.70) -- (0.17,-6.78) -- (0.25,-6.80) -- (0.38,-6.65) -- (0.54,-6.66) -- (0.57,-6.75) -- (0.74,-6.66) -- (0.69,-6.59) -- (0.71,-6.41) -- (0.85,-6.45) -- (0.91,-6.41) -- (0.84,-6.40) -- (0.95,-6.12) -- (1.04,-6.14) -- (1.04,-6.04) -- (1.12,-6.08) -- (1.18,-6.03) -- (1.10,-6.00) -- (1.09,-5.94) -- (1.15,-5.87) -- (1.21,-5.96) -- (1.20,-5.91) -- (1.25,-5.91) -- (1.27,-5.81) -- (1.20,-5.82) -- (1.16,-5.72) -- (1.19,-5.67) -- (1.23,-5.69) -- (1.22,-5.58) -- (1.26,-5.53) -- (1.36,-5.55) -- (1.32,-5.46) -- (1.44,-5.44) -- (1.42,-5.48) -- (1.51,-5.51) -- (1.49,-5.41) -- (1.63,-5.42) -- (1.60,-5.36) -- (1.65,-5.31) -- (1.65,-5.37) -- (1.70,-5.37) -- (1.73,-5.26) -- (1.84,-5.26) -- (1.93,-5.35) -- (1.98,-5.34) -- (1.94,-5.25) -- (2.04,-5.19) -- (1.98,-5.16) -- (2.04,-5.05) -- (2.09,-5.12) -- (2.18,-5.06) -- (2.17,-4.95) -- (2.20,-5.02) -- (2.28,-4.99) -- (2.34,-5.02) -- (2.29,-4.85) -- (2.40,-4.81) -- (2.37,-4.75) -- (2.55,-4.76) -- (2.56,-4.68) -- (2.64,-4.67) -- (2.63,-4.58) -- (2.70,-4.59) -- (2.91,-4.44) -- (3.00,-4.48) -- (3.22,-4.31) -- (3.60,-4.29) -- (3.72,-4.18) -- (3.88,-4.17) -- (4.00,-4.26) -- (4.13,-4.28) -- (4.41,-4.08) -- (4.55,-4.09) -- (4.56,-4.15) -- cycle;
\filldraw[fill=primary!10,draw=black!45,line width=0.4pt] (6.96,-0.01) -- (7.00,-0.15) -- (7.10,-0.05) -- (7.18,-0.08) -- (7.29,-0.05) -- (7.40,-0.22) -- (7.33,-0.30) -- (7.31,-0.26) -- (7.26,-0.31) -- (7.15,-0.27) -- (7.18,-0.34) -- (7.06,-0.42) -- (7.09,-0.50) -- (7.02,-0.55) -- (7.05,-0.61) -- (6.99,-0.60) -- (6.96,-0.68) -- (7.07,-0.82) -- (7.03,-0.82) -- (7.03,-0.90) -- (6.88,-0.93) -- (6.88,-0.99) -- (6.81,-1.07) -- (6.84,-1.17) -- (6.76,-1.23) -- (6.92,-1.27) -- (6.90,-1.31) -- (6.96,-1.39) -- (7.11,-1.50) -- (7.32,-1.56) -- (7.27,-1.66) -- (7.34,-1.74) -- (7.27,-1.81) -- (7.33,-1.87) -- (7.31,-2.34) -- (7.40,-2.51) -- (7.51,-2.53) -- (7.55,-2.65) -- (6.89,-2.69) -- (6.90,-2.76) -- (6.82,-2.72) -- (6.86,-2.80) -- (6.82,-2.79) -- (6.81,-2.84) -- (6.85,-2.92) -- (6.78,-2.92) -- (6.74,-2.85) -- (6.65,-3.01) -- (6.44,-2.99) -- (6.42,-2.93) -- (6.01,-3.09) -- (6.00,-3.05) -- (5.87,-3.06) -- (5.76,-3.16) -- (5.33,-3.01) -- (5.40,-2.94) -- (5.44,-2.78) -- (5.14,-2.53) -- (5.04,-2.56) -- (4.94,-2.45) -- (4.86,-2.50) -- (4.79,-2.38) -- (4.60,-2.44) -- (4.51,-2.32) -- (4.46,-2.36) -- (4.43,-2.31) -- (4.37,-2.33) -- (4.35,-2.26) -- (4.30,-2.29) -- (4.31,-2.14) -- (4.26,-2.17) -- (4.24,-2.12) -- (4.32,-2.02) -- (4.38,-2.13) -- (4.42,-2.10) -- (4.48,-2.16) -- (4.85,-2.18) -- (5.05,-2.30) -- (5.09,-2.13) -- (5.16,-2.11) -- (5.22,-1.84) -- (5.37,-1.83) -- (5.33,-1.78) -- (5.36,-1.71) -- (5.49,-1.75) -- (5.57,-1.70) -- (5.56,-1.63) -- (5.64,-1.70) -- (5.70,-1.57) -- (5.89,-1.49) -- (5.90,-1.42) -- (5.95,-1.49) -- (6.03,-1.42) -- (6.12,-1.50) -- (6.05,-1.22) -- (6.14,-1.17) -- (6.21,-1.22) -- (6.30,-1.09) -- (6.27,-1.05) -- (6.35,-1.08) -- (6.35,-0.99) -- (6.50,-1.00) -- (6.44,-0.95) -- (6.45,-0.89) -- (6.51,-0.88) -- (6.54,-0.82) -- (6.36,-0.70) -- (6.47,-0.60) -- (6.50,-0.66) -- (6.47,-0.55) -- (6.52,-0.54) -- (6.51,-0.50) -- (6.63,-0.53) -- (6.56,-0.45) -- (6.53,-0.35) -- (6.68,-0.23) -- (6.63,-0.15) -- (6.69,-0.02) -- (6.87,-0.08) -- cycle;
\filldraw[fill=primary!10,draw=black!45,line width=0.4pt] (7.67,0.00) -- (7.63,-0.16) -- (7.81,-0.13) -- (7.78,-0.24) -- (7.95,-0.29) -- (7.94,-0.19) -- (8.16,-0.09) -- (8.10,-0.32) -- (8.20,-0.35) -- (8.24,-0.24) -- (8.30,-0.22) -- (8.29,-0.35) -- (8.34,-0.35) -- (8.39,-0.44) -- (8.52,-0.42) -- (8.48,-0.49) -- (8.56,-0.56) -- (8.54,-0.70) -- (8.65,-0.74) -- (8.61,-0.83) -- (8.65,-0.90) -- (8.83,-0.85) -- (8.74,-0.99) -- (8.87,-0.95) -- (8.71,-1.13) -- (8.95,-1.13) -- (8.97,-1.21) -- (8.90,-1.17) -- (8.86,-1.23) -- (8.95,-1.37) -- (9.05,-1.37) -- (9.05,-1.28) -- (9.20,-1.30) -- (9.24,-1.21) -- (9.27,-1.32) -- (9.17,-1.34) -- (9.13,-1.49) -- (9.33,-1.51) -- (9.44,-1.72) -- (9.57,-1.71) -- (9.53,-1.87) -- (9.71,-1.88) -- (9.74,-1.95) -- (9.71,-1.99) -- (9.86,-2.12) -- (9.81,-2.19) -- (9.85,-2.38) -- (9.95,-2.39) -- (9.98,-2.58) -- (10.02,-2.60) -- (9.93,-2.74) -- (9.93,-2.89) -- (9.87,-3.00) -- (9.92,-3.10) -- (9.97,-3.14) -- (9.97,-3.06) -- (10.05,-3.05) -- (10.09,-3.15) -- (10.18,-3.13) -- (10.27,-3.35) -- (10.35,-3.32) -- (10.37,-3.38) -- (10.46,-3.41) -- (10.51,-3.48) -- (10.80,-3.59) -- (10.74,-3.68) -- (10.76,-3.82) -- (10.88,-3.97) -- (10.82,-4.22) -- (10.76,-4.28) -- (10.82,-4.39) -- (10.79,-4.47) -- (10.97,-4.69) -- (10.93,-4.74) -- (11.00,-4.83) -- (10.96,-4.87) -- (10.97,-4.95) -- (10.89,-4.86) -- (10.87,-4.78) -- (10.70,-4.71) -- (10.65,-4.62) -- (10.48,-4.67) -- (10.42,-4.63) -- (10.32,-4.72) -- (10.27,-4.66) -- (10.20,-4.68) -- (9.67,-5.05) -- (9.58,-5.11) -- (9.52,-5.08) -- (9.46,-5.16) -- (9.39,-5.15) -- (9.43,-5.24) -- (9.42,-5.27) -- (9.38,-5.24) -- (9.34,-5.34) -- (9.26,-5.28) -- (9.22,-5.31) -- (9.24,-5.37) -- (9.16,-5.36) -- (9.06,-5.43) -- (8.99,-5.57) -- (8.73,-5.66) -- (8.83,-5.55) -- (8.82,-5.51) -- (8.76,-5.50) -- (8.76,-5.44) -- (8.66,-5.43) -- (8.67,-5.39) -- (8.77,-5.38) -- (8.85,-5.22) -- (8.90,-5.23) -- (8.92,-5.17) -- (8.85,-5.13) -- (8.88,-5.05) -- (8.80,-5.02) -- (8.77,-5.22) -- (8.74,-5.17) -- (8.69,-5.23) -- (8.68,-5.07) -- (8.49,-5.11) -- (8.54,-5.06) -- (8.53,-5.00) -- (8.60,-5.01) -- (8.67,-4.91) -- (8.57,-4.92) -- (8.58,-4.87) -- (8.50,-4.87) -- (8.39,-4.93) -- (8.51,-4.78) -- (8.60,-4.78) -- (8.56,-4.71) -- (8.59,-4.66) -- (8.46,-4.69) -- (8.48,-4.63) -- (8.45,-4.62) -- (8.62,-4.59) -- (8.61,-4.44) -- (8.75,-4.45) -- (8.58,-4.43) -- (8.54,-4.27) -- (8.43,-4.39) -- (8.50,-4.48) -- (8.39,-4.46) -- (8.38,-4.30) -- (8.29,-4.24) -- (8.22,-4.32) -- (8.22,-4.28) -- (8.16,-4.25) -- (8.21,-4.23) -- (8.12,-4.11) -- (8.00,-4.06) -- (8.05,-4.03) -- (8.12,-4.07) -- (8.12,-4.00) -- (8.18,-4.07) -- (8.17,-3.84) -- (8.19,-3.81) -- (8.26,-3.87) -- (8.28,-3.82) -- (8.24,-3.75) -- (8.33,-3.71) -- (8.15,-3.58) -- (8.20,-3.54) -- (8.47,-3.57) -- (8.26,-3.42) -- (8.16,-3.48) -- (8.14,-3.40) -- (8.08,-3.37) -- (8.14,-3.32) -- (8.18,-3.37) -- (8.24,-3.23) -- (8.35,-3.16) -- (8.15,-3.12) -- (8.16,-3.04) -- (8.07,-2.98) -- (7.97,-3.00) -- (8.02,-2.85) -- (7.72,-2.78) -- (7.51,-2.61) -- (7.49,-2.51) -- (7.40,-2.51) -- (7.31,-2.34) -- (7.33,-1.87) -- (7.27,-1.81) -- (7.34,-1.74) -- (7.27,-1.66) -- (7.32,-1.56) -- (7.11,-1.50) -- (6.96,-1.39) -- (6.90,-1.31) -- (6.92,-1.27) -- (6.76,-1.23) -- (6.84,-1.17) -- (6.81,-1.07) -- (6.88,-0.99) -- (6.88,-0.93) -- (7.03,-0.90) -- (7.03,-0.82) -- (7.07,-0.82) -- (6.96,-0.68) -- (6.99,-0.60) -- (7.05,-0.61) -- (7.02,-0.55) -- (7.09,-0.50) -- (7.06,-0.42) -- (7.18,-0.34) -- (7.15,-0.27) -- (7.26,-0.31) -- (7.31,-0.26) -- (7.33,-0.30) -- (7.40,-0.22) -- (7.29,-0.02) -- (7.33,-0.02) -- (7.38,-0.12) -- (7.47,-0.11) -- (7.54,-0.17) -- (7.57,-0.04) -- cycle;
\filldraw[fill=primary!10,draw=black!45,line width=0.4pt] (7.19,-5.23) -- (7.28,-5.33) -- (7.34,-5.30) -- (7.35,-5.38) -- (7.49,-5.35) -- (7.49,-5.30) -- (7.55,-5.41) -- (7.88,-5.34) -- (7.94,-5.37) -- (8.00,-5.29) -- (8.08,-5.40) -- (8.08,-5.35) -- (8.13,-5.34) -- (8.20,-5.47) -- (8.27,-5.41) -- (8.34,-5.41) -- (8.39,-5.35) -- (8.50,-5.46) -- (8.61,-5.40) -- (8.59,-5.54) -- (8.63,-5.58) -- (8.55,-5.63) -- (8.59,-5.78) -- (8.68,-5.74) -- (8.84,-5.88) -- (8.86,-5.97) -- (9.06,-6.01) -- (9.09,-5.95) -- (9.23,-5.99) -- (9.22,-6.07) -- (9.29,-6.04) -- (9.32,-6.08) -- (9.63,-6.12) -- (9.50,-6.27) -- (9.48,-6.35) -- (9.56,-6.45) -- (9.65,-6.46) -- (9.66,-6.54) -- (9.61,-6.56) -- (9.60,-6.66) -- (9.69,-6.69) -- (9.71,-6.86) -- (9.80,-6.95) -- (9.72,-6.94) -- (9.65,-6.81) -- (9.58,-6.90) -- (9.46,-6.88) -- (9.43,-7.10) -- (9.33,-7.08) -- (9.28,-7.13) -- (9.20,-7.05) -- (9.08,-7.02) -- (9.07,-7.11) -- (8.98,-7.08) -- (8.74,-7.22) -- (8.70,-7.35) -- (8.75,-7.46) -- (8.68,-7.48) -- (8.66,-7.56) -- (8.44,-7.48) -- (8.49,-7.43) -- (8.34,-7.41) -- (8.38,-7.49) -- (8.27,-7.50) -- (8.28,-7.60) -- (8.15,-7.54) -- (8.17,-7.62) -- (8.06,-7.61) -- (8.03,-7.68) -- (7.92,-7.72) -- (7.90,-7.64) -- (7.86,-7.64) -- (7.86,-7.82) -- (7.79,-7.86) -- (7.63,-7.75) -- (7.55,-7.92) -- (7.55,-8.05) -- (7.49,-8.06) -- (7.48,-7.93) -- (7.40,-7.84) -- (7.40,-7.76) -- (7.25,-7.63) -- (7.18,-7.61) -- (7.06,-7.33) -- (7.14,-7.21) -- (7.12,-7.17) -- (7.20,-6.97) -- (7.01,-6.81) -- (6.99,-6.73) -- (6.77,-6.60) -- (6.84,-6.51) -- (6.97,-6.64) -- (7.03,-6.59) -- (7.02,-6.53) -- (7.26,-6.56) -- (7.35,-6.51) -- (7.38,-6.56) -- (7.41,-6.52) -- (7.53,-6.60) -- (7.57,-6.57) -- (7.56,-6.47) -- (7.70,-6.26) -- (7.71,-5.94) -- (7.55,-5.98) -- (7.38,-5.81) -- (7.26,-5.78) -- (7.19,-5.68) -- (7.28,-5.64) -- (7.27,-5.54) -- (7.11,-5.47) -- (7.04,-5.38) -- (7.04,-5.30) -- (7.17,-5.28) -- cycle;
\filldraw[fill=primary!10,draw=black!45,line width=0.4pt] (1.25,-2.92) -- (1.30,-2.95) -- (1.26,-3.02) -- (1.33,-3.06) -- (1.33,-2.97) -- (1.43,-3.02) -- (1.44,-2.93) -- (1.51,-2.99) -- (1.53,-2.93) -- (1.59,-2.99) -- (1.66,-2.98) -- (1.72,-3.06) -- (1.75,-2.96) -- (1.71,-2.87) -- (1.77,-2.85) -- (1.77,-2.92) -- (1.87,-2.94) -- (1.86,-2.88) -- (1.96,-2.85) -- (1.92,-2.87) -- (1.94,-2.91) -- (2.06,-2.91) -- (2.02,-2.97) -- (2.06,-3.01) -- (2.12,-2.92) -- (2.16,-2.94) -- (2.15,-2.88) -- (2.24,-2.87) -- (2.20,-2.93) -- (2.28,-3.03) -- (2.30,-2.98) -- (2.32,-2.99) -- (2.40,-2.88) -- (2.43,-2.88) -- (2.49,-2.96) -- (2.40,-2.89) -- (2.29,-3.05) -- (2.23,-3.05) -- (2.30,-3.23) -- (2.26,-3.30) -- (2.35,-3.52) -- (2.34,-3.59) -- (2.40,-3.65) -- (2.62,-3.68) -- (2.63,-3.75) -- (2.71,-3.79) -- (2.69,-3.88) -- (2.77,-4.02) -- (2.82,-4.00) -- (2.63,-4.35) -- (2.73,-4.40) -- (2.69,-4.45) -- (2.73,-4.56) -- (2.67,-4.61) -- (2.63,-4.58) -- (2.64,-4.67) -- (2.56,-4.68) -- (2.55,-4.76) -- (2.37,-4.75) -- (2.40,-4.81) -- (2.29,-4.85) -- (2.34,-5.02) -- (2.28,-4.99) -- (2.20,-5.02) -- (2.17,-4.95) -- (2.18,-5.06) -- (2.09,-5.12) -- (2.04,-5.05) -- (1.98,-5.16) -- (2.04,-5.19) -- (1.94,-5.25) -- (1.98,-5.34) -- (1.93,-5.35) -- (1.84,-5.26) -- (1.73,-5.26) -- (1.70,-5.37) -- (1.65,-5.37) -- (1.65,-5.31) -- (1.60,-5.36) -- (1.63,-5.42) -- (1.49,-5.41) -- (1.51,-5.25) -- (1.45,-5.24) -- (1.46,-5.29) -- (1.40,-5.24) -- (1.47,-5.21) -- (1.48,-5.13) -- (1.53,-5.15) -- (1.50,-5.09) -- (1.55,-5.10) -- (1.54,-5.04) -- (1.60,-5.04) -- (1.57,-5.00) -- (1.62,-4.95) -- (1.58,-4.95) -- (1.56,-4.89) -- (1.59,-4.82) -- (1.64,-4.88) -- (1.64,-4.83) -- (1.69,-4.83) -- (1.68,-4.75) -- (1.62,-4.72) -- (1.63,-4.64) -- (1.59,-4.62) -- (1.63,-4.58) -- (1.61,-4.52) -- (1.55,-4.54) -- (1.54,-4.50) -- (1.59,-4.47) -- (1.41,-4.31) -- (1.48,-4.20) -- (1.43,-4.14) -- (1.34,-4.14) -- (1.35,-4.07) -- (1.30,-4.02) -- (1.31,-3.75) -- (1.17,-3.70) -- (1.06,-3.51) -- (1.13,-3.46) -- (1.12,-3.41) -- (1.08,-3.37) -- (1.09,-3.27) -- (1.02,-3.23) -- (1.07,-3.22) -- (1.07,-3.18) -- (0.97,-3.11) -- (1.03,-3.08) -- (0.97,-3.03) -- (1.03,-3.01) -- (0.99,-2.86) -- (1.14,-2.88) -- (1.14,-2.94) -- (1.20,-2.84) -- (1.24,-2.85) -- cycle;
\filldraw[fill=primary!10,draw=black!45,line width=0.4pt] (2.81,-6.58) -- (2.82,-6.73) -- (2.91,-6.75) -- (2.92,-6.69) -- (3.03,-6.75) -- (2.99,-6.84) -- (2.85,-6.94) -- (2.93,-7.05) -- (2.77,-7.07) -- (2.86,-7.21) -- (2.64,-7.19) -- (2.65,-7.30) -- (2.61,-7.32) -- (2.51,-7.23) -- (2.34,-7.23) -- (2.33,-7.33) -- (2.16,-7.61) -- (1.82,-7.62) -- (1.74,-7.47) -- (1.67,-7.43) -- (1.50,-7.61) -- (1.23,-7.43) -- (1.22,-7.37) -- (1.28,-7.29) -- (1.59,-7.14) -- (1.79,-6.94) -- (1.94,-6.87) -- (2.13,-6.55) -- (2.28,-6.50) -- (2.42,-6.38) -- (2.46,-6.50) -- cycle;
\filldraw[fill=primary!10,draw=black!45,line width=0.4pt] (2.60,-7.46) -- (2.64,-7.51) -- (2.80,-7.52) -- (2.86,-7.46) -- (2.91,-7.50) -- (3.05,-7.41) -- (3.04,-7.36) -- (3.13,-7.43) -- (3.20,-7.33) -- (3.35,-7.27) -- (3.44,-7.34) -- (3.55,-7.21) -- (3.65,-7.45) -- (3.70,-7.40) -- (3.80,-7.41) -- (3.78,-7.49) -- (3.88,-7.57) -- (3.88,-7.79) -- (3.96,-7.86) -- (4.11,-8.13) -- (4.06,-8.17) -- (4.11,-8.26) -- (4.05,-8.35) -- (4.11,-8.68) -- (4.03,-8.85) -- (4.07,-8.92) -- (3.99,-8.96) -- (4.03,-9.03) -- (4.00,-9.08) -- (4.12,-9.21) -- (4.08,-9.25) -- (4.17,-9.30) -- (3.99,-9.40) -- (3.86,-9.38) -- (3.70,-9.50) -- (3.61,-9.48) -- (3.54,-9.37) -- (3.10,-9.21) -- (3.02,-9.22) -- (2.93,-9.41) -- (2.85,-9.43) -- (2.69,-9.32) -- (2.72,-9.16) -- (2.56,-9.10) -- (2.59,-8.99) -- (2.57,-8.91) -- (2.42,-8.83) -- (2.20,-8.84) -- (2.03,-8.52) -- (1.82,-8.49) -- (1.79,-8.35) -- (1.70,-8.28) -- (1.65,-8.03) -- (1.74,-7.97) -- (1.77,-7.87) -- (1.87,-7.76) -- (2.04,-7.75) -- (1.96,-7.62) -- (2.16,-7.61) -- (2.33,-7.33) -- (2.34,-7.23) -- (2.51,-7.23) -- (2.61,-7.32) -- cycle;
\filldraw[fill=primary!10,draw=black!45,line width=0.4pt] (2.07,-1.55) -- (1.96,-1.71) -- (2.09,-1.74) -- (2.14,-1.83) -- (2.18,-1.80) -- (2.13,-1.75) -- (2.19,-1.77) -- (2.17,-1.72) -- (2.21,-1.74) -- (2.26,-1.88) -- (2.22,-1.97) -- (2.30,-1.90) -- (2.27,-1.82) -- (2.38,-1.81) -- (2.46,-1.69) -- (2.57,-1.66) -- (2.48,-1.76) -- (2.62,-1.76) -- (2.56,-1.82) -- (2.63,-1.86) -- (2.66,-1.81) -- (2.74,-1.90) -- (2.78,-1.78) -- (2.84,-1.82) -- (2.91,-1.80) -- (2.83,-1.97) -- (2.87,-2.00) -- (2.92,-1.95) -- (2.91,-2.07) -- (2.97,-2.10) -- (3.01,-2.08) -- (2.97,-2.01) -- (3.11,-1.86) -- (3.15,-1.87) -- (3.15,-2.01) -- (3.22,-1.93) -- (3.27,-1.95) -- (3.29,-2.05) -- (3.33,-2.02) -- (3.39,-2.06) -- (3.37,-1.96) -- (3.43,-1.95) -- (3.61,-2.03) -- (3.60,-2.09) -- (3.65,-2.13) -- (3.53,-2.18) -- (3.58,-2.20) -- (3.57,-2.33) -- (3.60,-2.29) -- (3.64,-2.33) -- (3.58,-2.36) -- (3.65,-2.40) -- (3.60,-2.52) -- (3.65,-2.61) -- (3.55,-3.14) -- (3.66,-3.56) -- (3.56,-3.73) -- (3.64,-3.74) -- (3.72,-3.83) -- (3.58,-4.00) -- (3.49,-3.92) -- (3.47,-3.97) -- (3.38,-3.98) -- (3.44,-4.04) -- (3.39,-4.06) -- (3.39,-4.12) -- (3.33,-4.14) -- (3.28,-4.10) -- (3.30,-4.18) -- (3.25,-4.18) -- (3.20,-4.10) -- (3.21,-4.14) -- (3.16,-4.20) -- (3.07,-4.14) -- (3.01,-4.21) -- (2.98,-4.07) -- (2.93,-4.05) -- (2.94,-4.13) -- (2.81,-4.13) -- (2.86,-4.20) -- (2.82,-4.20) -- (2.76,-4.13) -- (2.82,-4.00) -- (2.77,-4.02) -- (2.69,-3.88) -- (2.71,-3.79) -- (2.63,-3.75) -- (2.62,-3.68) -- (2.40,-3.65) -- (2.34,-3.59) -- (2.35,-3.52) -- (2.26,-3.30) -- (2.30,-3.23) -- (2.23,-3.05) -- (2.29,-3.05) -- (2.40,-2.89) -- (2.46,-2.95) -- (2.49,-2.93) -- (2.40,-2.88) -- (2.30,-3.02) -- (2.22,-2.99) -- (2.20,-2.93) -- (2.29,-2.74) -- (2.20,-2.60) -- (2.05,-2.62) -- (1.99,-2.46) -- (2.01,-2.24) -- (1.92,-2.21) -- (1.92,-2.16) -- (1.80,-2.16) -- (1.72,-1.98) -- (1.67,-2.02) -- (1.66,-1.95) -- (1.71,-1.96) -- (1.72,-1.91) -- (1.67,-1.87) -- (1.71,-1.82) -- (1.64,-1.77) -- (1.56,-1.79) -- (1.57,-1.73) -- (1.52,-1.71) -- (1.52,-1.63) -- (1.59,-1.65) -- (1.65,-1.56) -- (1.74,-1.67) -- (1.82,-1.58) -- (1.88,-1.58) -- (1.92,-1.66) -- (1.91,-1.59) -- cycle;
\filldraw[fill=primary!10,draw=black!45,line width=0.4pt] (3.80,-1.78) -- (3.85,-1.86) -- (3.88,-1.82) -- (4.00,-1.90) -- (3.98,-1.94) -- (4.32,-2.02) -- (4.24,-2.12) -- (4.26,-2.17) -- (4.31,-2.14) -- (4.30,-2.29) -- (4.35,-2.26) -- (4.37,-2.33) -- (4.43,-2.31) -- (4.46,-2.36) -- (4.51,-2.32) -- (4.60,-2.44) -- (4.79,-2.38) -- (4.86,-2.50) -- (4.97,-2.46) -- (5.01,-2.54) -- (4.96,-2.61) -- (5.12,-2.89) -- (5.10,-3.10) -- (5.23,-3.28) -- (5.21,-3.35) -- (5.11,-3.33) -- (5.15,-3.46) -- (5.07,-3.52) -- (4.87,-3.50) -- (4.74,-3.56) -- (4.50,-3.35) -- (4.22,-3.58) -- (3.64,-3.50) -- (3.55,-3.14) -- (3.65,-2.61) -- (3.60,-2.52) -- (3.65,-2.40) -- (3.58,-2.36) -- (3.64,-2.33) -- (3.60,-2.29) -- (3.57,-2.33) -- (3.58,-2.20) -- (3.53,-2.18) -- (3.65,-2.13) -- (3.60,-2.09) -- (3.61,-2.03) -- (3.50,-1.96) -- (3.54,-1.93) -- (3.50,-1.85) -- (3.69,-1.87) -- (3.67,-1.81) -- (3.75,-1.82) -- (3.72,-1.75) -- cycle;
\fill (6.44,-4.79) circle (1.6pt) node[below=1pt] {\bfseries Guarapuava};
\fill (5.08,-7.04) circle (1.6pt) node[below=1pt] {Pinhão};
\fill (2.51,-6.14) circle (1.6pt) node[below=1pt] {Candói};
\fill (5.97,-2.17) circle (1.6pt) node[above=1pt] {Turvo};
\fill (9.67,-3.44) circle (1.6pt) node[below=1pt] {Prudentópolis};
\fill (9.01,-6.10) circle (1.6pt) node[below=1pt] {Inácio Martins};
\fill (1.94,-4.63) circle (1.6pt) node[above=1pt] {Cantagalo};
\fill (1.99,-7.35) circle (1.6pt) node[above=1pt] {Foz do Jordão};
\fill (2.61,-8.05) circle (1.6pt) node[below=1pt] {Reserva do Iguaçu};
\fill (2.85,-3.28) circle (1.6pt) node[above=1pt] {Goioxim};
\fill (3.98,-2.44) circle (1.6pt) node[above=1pt] {Campina do Simão};
\draw[line width=0.6pt] (0.2,-9.80) -- ++(0,-0.1) -- ++(2.68,0) -- ++(0,0.1); \node[above] at (1.54,-9.85) {40 km};
\draw[black!60,line width=0.4pt] (12.23,-2.58) -- (12.13,-2.56) -- (12.08,-2.48) -- (12.08,-2.43) -- (12.03,-2.39) -- (12.04,-2.29) -- (12.00,-2.18) -- (11.96,-2.19) -- (11.93,-2.14) -- (11.89,-2.14) -- (11.88,-2.17) -- (11.88,-2.10) -- (11.83,-2.13) -- (11.84,-2.15) -- (11.81,-2.14) -- (11.80,-2.16) -- (11.78,-2.14) -- (11.73,-2.16) -- (11.70,-2.22) -- (11.66,-2.17) -- (11.61,-2.16) -- (11.60,-2.08) -- (11.70,-1.90) -- (11.70,-1.77) -- (11.76,-1.60) -- (11.75,-1.50) -- (11.80,-1.41) -- (11.75,-1.28) -- (11.88,-1.17) -- (11.95,-0.87) -- (12.08,-0.79) -- (12.14,-0.67) -- (12.14,-0.59) -- (12.19,-0.53) -- (12.50,-0.33) -- (12.65,-0.37) -- (12.71,-0.33) -- (12.76,-0.37) -- (12.89,-0.36) -- (12.91,-0.40) -- (12.94,-0.38) -- (12.97,-0.30) -- (13.10,-0.37) -- (13.16,-0.36) -- (13.18,-0.39) -- (13.43,-0.39) -- (13.49,-0.44) -- (13.64,-0.47) -- (13.69,-0.53) -- (13.68,-0.56) -- (13.72,-0.57) -- (13.76,-0.53) -- (13.99,-0.56) -- (14.13,-0.53) -- (14.18,-0.59) -- (14.17,-0.62) -- (14.26,-0.65) -- (14.30,-0.69) -- (14.32,-0.80) -- (14.34,-0.82) -- (14.33,-0.83) -- (14.36,-0.85) -- (14.32,-0.90) -- (14.33,-0.98) -- (14.37,-1.02) -- (14.36,-1.09) -- (14.33,-1.11) -- (14.39,-1.15) -- (14.43,-1.25) -- (14.48,-1.28) -- (14.47,-1.33) -- (14.51,-1.38) -- (14.56,-1.40) -- (14.49,-1.53) -- (14.50,-1.60) -- (14.63,-1.61) -- (14.66,-1.58) -- (14.69,-1.61) -- (14.76,-1.59) -- (14.85,-1.62) -- (14.90,-1.60) -- (14.94,-1.64) -- (14.88,-1.81) -- (14.91,-1.85) -- (14.99,-1.79) -- (15.03,-1.84) -- (15.08,-1.79) -- (15.13,-1.89) -- (15.11,-1.92) -- (15.15,-1.95) -- (15.20,-1.94) -- (15.10,-2.08) -- (15.04,-2.10) -- (15.05,-2.14) -- (14.97,-2.20) -- (14.89,-2.39) -- (14.57,-2.41) -- (14.36,-2.55) -- (14.27,-2.51) -- (14.21,-2.44) -- (14.15,-2.41) -- (14.14,-2.44) -- (14.13,-2.41) -- (14.06,-2.44) -- (14.03,-2.42) -- (14.02,-2.45) -- (13.98,-2.42) -- (13.94,-2.46) -- (13.94,-2.49) -- (13.88,-2.43) -- (13.81,-2.41) -- (13.80,-2.43) -- (13.82,-2.44) -- (13.78,-2.44) -- (13.72,-2.53) -- (13.73,-2.55) -- (13.65,-2.55) -- (13.63,-2.58) -- (13.61,-2.55) -- (13.61,-2.58) -- (13.53,-2.54) -- (13.44,-2.60) -- (13.41,-2.66) -- (13.45,-2.78) -- (13.36,-2.81) -- (13.35,-2.84) -- (13.30,-2.76) -- (13.10,-2.77) -- (13.02,-2.75) -- (12.93,-2.67) -- (12.70,-2.66) -- (12.63,-2.61) -- (12.43,-2.64) -- (12.33,-2.55) -- cycle;
\filldraw[fill=primary!60,draw=primary] (13.16,-2.41) -- (13.10,-2.42) -- (13.05,-2.40) -- (13.02,-2.41) -- (13.01,-2.37) -- (12.98,-2.37) -- (12.94,-2.32) -- (12.94,-2.30) -- (12.97,-2.27) -- (12.94,-2.25) -- (12.92,-2.27) -- (12.87,-2.23) -- (12.85,-2.24) -- (12.86,-2.21) -- (12.81,-2.22) -- (12.80,-2.20) -- (12.86,-2.19) -- (12.86,-2.17) -- (12.90,-2.13) -- (12.90,-2.10) -- (12.93,-2.10) -- (12.92,-2.08) -- (12.94,-2.04) -- (12.89,-1.94) -- (12.88,-1.88) -- (12.94,-1.90) -- (12.94,-1.88) -- (12.97,-1.90) -- (12.99,-1.88) -- (12.93,-1.78) -- (12.97,-1.78) -- (12.96,-1.79) -- (12.98,-1.79) -- (12.98,-1.81) -- (12.99,-1.80) -- (13.01,-1.79) -- (13.02,-1.80) -- (13.04,-1.80) -- (13.04,-1.82) -- (13.06,-1.80) -- (13.08,-1.82) -- (13.10,-1.79) -- (13.21,-1.84) -- (13.24,-1.79) -- (13.30,-1.77) -- (13.29,-1.75) -- (13.33,-1.73) -- (13.32,-1.71) -- (13.34,-1.69) -- (13.35,-1.65) -- (13.41,-1.67) -- (13.43,-1.65) -- (13.42,-1.66) -- (13.45,-1.68) -- (13.46,-1.66) -- (13.46,-1.68) -- (13.48,-1.67) -- (13.50,-1.72) -- (13.52,-1.72) -- (13.51,-1.74) -- (13.53,-1.74) -- (13.53,-1.76) -- (13.55,-1.75) -- (13.54,-1.77) -- (13.60,-1.82) -- (13.62,-1.86) -- (13.61,-1.90) -- (13.68,-1.94) -- (13.69,-2.05) -- (13.67,-2.03) -- (13.63,-2.03) -- (13.53,-2.11) -- (13.50,-2.11) -- (13.52,-2.13) -- (13.58,-2.15) -- (13.57,-2.17) -- (13.60,-2.21) -- (13.57,-2.21) -- (13.57,-2.23) -- (13.51,-2.24) -- (13.50,-2.26) -- (13.48,-2.25) -- (13.43,-2.29) -- (13.42,-2.28) -- (13.37,-2.34) -- (13.37,-2.37) -- (13.35,-2.37) -- (13.34,-2.36) -- (13.33,-2.38) -- (13.31,-2.36) -- (13.27,-2.40) -- (13.21,-2.41) -- (13.19,-2.39) -- (13.20,-2.35) -- cycle;
\fill (14.52,-2.06) circle (1.2pt) node[above] {\scriptsize Curitiba};
\node[anchor=south west,font=\small\bfseries] at (11.60,-0.20) {Paraná};
\node[anchor=north west,font=\scriptsize,text width=3.6cm] at (11.60,-3.04) {Em destaque, a região do mapa ao lado.};
\draw[-{Stealth[length=6pt]},line width=0.8pt] (13.40,-8.90) -- ++(0,0.8) node[above] {N};
\end{tikzpicture}}
\par\smallskip{\footnotesize Cada área colorida é um município; o ponto marca a sede, onde a rede precisa chegar.}
\end{figura}

Entre Cantagalo, no oeste, e Prudentópolis, no leste, são mais de 110~km em linha
reta. A pergunta é como levar energia a todas essas sedes gastando o mínimo de
linha.

\subsection{Modelagem como grafo}

\begin{table}[H]
\centering
\caption{Modelagem da rede elétrica como grafo}\label{tab:modelo-rede}
\begin{tabularx}{\textwidth}{l >{\raggedright\arraybackslash}X}
\toprule
Elemento & Significado no problema \\
\midrule
Vértice & Uma cidade (11 no total) \\
Aresta & Um trecho de linha de transmissão que pode ser construído entre duas cidades \\
Peso & Comprimento do trecho, em km \\
Direção & Não direcionado: uma linha conduz energia nos dois sentidos \\
Resultado desejado & Árvore geradora mínima: o conjunto de trechos que liga todas as cidades com a menor extensão total \\
\bottomrule
\end{tabularx}
\fonte{elaborada pelo autor (2026).}
\end{table}

Cidades: Guarapuava, Pinhão, Candói, Turvo, Prudentópolis, Inácio Martins,
Cantagalo, Foz do Jordão, Reserva do Iguaçu, Goioxim e Campina do Simão.

\textbf{Dados.} Os pesos são distâncias em linha reta entre as sedes dos
municípios, calculadas a partir de coordenadas geográficas aproximadas pela
fórmula de Haversine. Foram considerados possíveis apenas os trechos de até
60~km, o que resultou em 25 arestas e num grafo conexo. As distâncias são
aproximadas e servem para demonstrar o método; um projeto real usaria o traçado
e o custo de cada trecho.

A fórmula de Haversine calcula a distância entre dois pontos da Terra a partir da
latitude e da longitude de cada um, levando em conta que a Terra é redonda. Aqui
ela serve apenas para gerar os pesos: por exemplo, Guarapuava e Pinhão ficam a
39,1~km uma da outra.

\begin{figura}{As 10 linhas escolhidas pelo Kruskal somam 294,6 km}{fig:rede-final}{elaborada pelo autor (2026), com base na malha municipal do IBGE \cite{geodatabr} e na saída do programa.}
\resizebox{0.92\textwidth}{!}{\begin{tikzpicture}[x=1cm,y=1cm,font=\footnotesize]
\filldraw[fill=primary!6,draw=black!25,line width=0.3pt] (5.14,-2.53) -- (5.44,-2.78) -- (5.40,-2.94) -- (5.33,-3.01) -- (5.76,-3.16) -- (5.87,-3.06) -- (6.00,-3.05) -- (6.01,-3.09) -- (6.42,-2.93) -- (6.44,-2.99) -- (6.65,-3.01) -- (6.74,-2.85) -- (6.78,-2.92) -- (6.85,-2.92) -- (6.81,-2.84) -- (6.82,-2.79) -- (6.86,-2.80) -- (6.82,-2.72) -- (6.90,-2.76) -- (6.89,-2.69) -- (7.55,-2.65) -- (7.72,-2.78) -- (8.02,-2.85) -- (7.97,-3.00) -- (8.07,-2.98) -- (8.16,-3.04) -- (8.15,-3.12) -- (8.35,-3.16) -- (8.24,-3.23) -- (8.18,-3.37) -- (8.14,-3.32) -- (8.08,-3.37) -- (8.14,-3.40) -- (8.16,-3.48) -- (8.26,-3.42) -- (8.47,-3.57) -- (8.20,-3.54) -- (8.15,-3.58) -- (8.33,-3.71) -- (8.24,-3.75) -- (8.28,-3.82) -- (8.26,-3.87) -- (8.19,-3.81) -- (8.17,-3.84) -- (8.18,-4.07) -- (8.12,-4.00) -- (8.12,-4.07) -- (8.05,-4.03) -- (8.00,-4.06) -- (8.12,-4.11) -- (8.21,-4.23) -- (8.16,-4.25) -- (8.22,-4.28) -- (8.22,-4.32) -- (8.29,-4.24) -- (8.38,-4.30) -- (8.39,-4.46) -- (8.50,-4.48) -- (8.43,-4.39) -- (8.54,-4.27) -- (8.58,-4.43) -- (8.75,-4.45) -- (8.61,-4.44) -- (8.62,-4.59) -- (8.45,-4.62) -- (8.48,-4.63) -- (8.46,-4.69) -- (8.59,-4.66) -- (8.56,-4.71) -- (8.60,-4.78) -- (8.51,-4.78) -- (8.39,-4.93) -- (8.50,-4.87) -- (8.58,-4.87) -- (8.57,-4.92) -- (8.67,-4.91) -- (8.60,-5.01) -- (8.53,-5.00) -- (8.54,-5.06) -- (8.49,-5.11) -- (8.68,-5.07) -- (8.69,-5.23) -- (8.74,-5.17) -- (8.77,-5.22) -- (8.80,-5.02) -- (8.88,-5.05) -- (8.85,-5.13) -- (8.92,-5.17) -- (8.90,-5.23) -- (8.85,-5.22) -- (8.77,-5.38) -- (8.67,-5.39) -- (8.66,-5.43) -- (8.76,-5.44) -- (8.76,-5.50) -- (8.82,-5.51) -- (8.83,-5.55) -- (8.65,-5.70) -- (8.59,-5.54) -- (8.61,-5.40) -- (8.50,-5.46) -- (8.39,-5.35) -- (8.34,-5.41) -- (8.27,-5.41) -- (8.20,-5.47) -- (8.13,-5.34) -- (8.08,-5.35) -- (8.08,-5.40) -- (8.00,-5.29) -- (7.94,-5.37) -- (7.88,-5.34) -- (7.55,-5.41) -- (7.49,-5.30) -- (7.49,-5.35) -- (7.35,-5.38) -- (7.34,-5.30) -- (7.28,-5.33) -- (7.19,-5.23) -- (7.17,-5.28) -- (7.04,-5.30) -- (7.04,-5.38) -- (7.11,-5.47) -- (7.27,-5.54) -- (7.28,-5.64) -- (7.19,-5.68) -- (7.26,-5.78) -- (7.38,-5.81) -- (7.55,-5.98) -- (7.71,-5.94) -- (7.70,-6.26) -- (7.56,-6.47) -- (7.57,-6.57) -- (7.53,-6.60) -- (7.41,-6.52) -- (7.38,-6.56) -- (7.35,-6.51) -- (7.26,-6.56) -- (7.06,-6.53) -- (6.97,-6.64) -- (6.84,-6.51) -- (6.77,-6.60) -- (6.80,-6.65) -- (6.74,-6.60) -- (6.67,-6.63) -- (6.69,-6.71) -- (6.65,-6.77) -- (6.57,-6.78) -- (6.52,-6.73) -- (6.35,-6.78) -- (6.37,-6.81) -- (6.33,-6.84) -- (6.29,-6.78) -- (6.26,-6.83) -- (6.28,-6.93) -- (6.11,-6.87) -- (6.09,-6.78) -- (6.17,-6.72) -- (6.17,-6.55) -- (6.32,-6.55) -- (6.24,-6.38) -- (6.21,-6.41) -- (6.24,-6.45) -- (6.17,-6.51) -- (6.07,-6.47) -- (6.07,-6.41) -- (5.99,-6.44) -- (6.01,-6.52) -- (5.92,-6.45) -- (5.82,-6.47) -- (5.75,-6.51) -- (5.74,-6.59) -- (5.69,-6.58) -- (5.60,-6.70) -- (5.55,-6.68) -- (5.54,-6.76) -- (5.47,-6.66) -- (5.43,-6.42) -- (5.31,-6.35) -- (5.31,-6.28) -- (5.23,-6.27) -- (5.20,-6.33) -- (5.23,-6.40) -- (5.18,-6.44) -- (5.11,-6.40) -- (5.12,-6.32) -- (5.07,-6.35) -- (4.78,-6.27) -- (4.75,-6.36) -- (4.78,-6.43) -- (4.66,-6.37) -- (4.55,-6.42) -- (4.53,-6.34) -- (4.44,-6.30) -- (4.47,-6.39) -- (4.36,-6.41) -- (4.27,-6.34) -- (4.32,-6.24) -- (4.17,-6.32) -- (4.15,-6.24) -- (4.20,-6.20) -- (4.18,-6.12) -- (4.05,-6.00) -- (4.01,-5.84) -- (4.10,-5.90) -- (4.20,-5.89) -- (4.21,-5.81) -- (4.30,-5.70) -- (4.21,-5.63) -- (4.28,-5.60) -- (4.28,-5.51) -- (4.34,-5.54) -- (4.32,-5.51) -- (4.45,-5.38) -- (4.45,-5.22) -- (4.33,-5.15) -- (4.40,-4.99) -- (4.38,-4.94) -- (4.24,-4.96) -- (4.22,-4.88) -- (4.42,-4.69) -- (4.37,-4.53) -- (4.49,-4.33) -- (4.55,-4.09) -- (4.41,-4.08) -- (4.13,-4.28) -- (4.00,-4.26) -- (3.88,-4.17) -- (3.72,-4.18) -- (3.60,-4.29) -- (3.22,-4.31) -- (3.00,-4.48) -- (2.91,-4.44) -- (2.73,-4.56) -- (2.69,-4.45) -- (2.73,-4.40) -- (2.63,-4.35) -- (2.69,-4.30) -- (2.69,-4.21) -- (2.77,-4.15) -- (2.86,-4.20) -- (2.81,-4.13) -- (2.86,-4.11) -- (2.94,-4.13) -- (2.93,-4.05) -- (2.98,-4.07) -- (3.01,-4.21) -- (3.07,-4.14) -- (3.16,-4.20) -- (3.21,-4.14) -- (3.20,-4.10) -- (3.25,-4.18) -- (3.30,-4.18) -- (3.28,-4.10) -- (3.33,-4.14) -- (3.39,-4.12) -- (3.39,-4.06) -- (3.44,-4.04) -- (3.38,-3.98) -- (3.47,-3.97) -- (3.49,-3.92) -- (3.58,-4.00) -- (3.72,-3.83) -- (3.64,-3.74) -- (3.56,-3.73) -- (3.66,-3.56) -- (3.64,-3.50) -- (4.22,-3.58) -- (4.50,-3.35) -- (4.74,-3.56) -- (4.87,-3.50) -- (5.07,-3.52) -- (5.15,-3.46) -- (5.11,-3.33) -- (5.21,-3.35) -- (5.23,-3.28) -- (5.10,-3.10) -- (5.12,-2.89) -- (4.96,-2.61) -- (5.01,-2.54) -- cycle;
\filldraw[fill=primary!6,draw=black!25,line width=0.3pt] (4.33,-6.25) -- (4.27,-6.34) -- (4.40,-6.43) -- (4.47,-6.39) -- (4.44,-6.30) -- (4.53,-6.34) -- (4.55,-6.42) -- (4.66,-6.37) -- (4.78,-6.43) -- (4.75,-6.36) -- (4.78,-6.27) -- (5.07,-6.35) -- (5.12,-6.32) -- (5.11,-6.40) -- (5.21,-6.43) -- (5.23,-6.27) -- (5.31,-6.28) -- (5.31,-6.35) -- (5.43,-6.42) -- (5.47,-6.66) -- (5.54,-6.76) -- (5.55,-6.68) -- (5.60,-6.70) -- (5.69,-6.58) -- (5.74,-6.59) -- (5.75,-6.51) -- (5.82,-6.47) -- (5.92,-6.45) -- (6.01,-6.52) -- (5.99,-6.44) -- (6.07,-6.41) -- (6.07,-6.47) -- (6.17,-6.51) -- (6.24,-6.45) -- (6.21,-6.41) -- (6.24,-6.38) -- (6.32,-6.55) -- (6.17,-6.55) -- (6.17,-6.72) -- (6.09,-6.78) -- (6.11,-6.87) -- (6.20,-6.93) -- (6.28,-6.93) -- (6.26,-6.83) -- (6.29,-6.78) -- (6.33,-6.84) -- (6.37,-6.81) -- (6.35,-6.78) -- (6.52,-6.73) -- (6.57,-6.78) -- (6.65,-6.77) -- (6.69,-6.71) -- (6.67,-6.63) -- (6.74,-6.60) -- (6.99,-6.73) -- (7.01,-6.81) -- (7.20,-6.97) -- (7.12,-7.17) -- (7.14,-7.21) -- (7.06,-7.33) -- (7.18,-7.61) -- (7.25,-7.63) -- (7.40,-7.76) -- (7.40,-7.84) -- (7.48,-7.93) -- (7.49,-8.06) -- (7.42,-8.00) -- (7.36,-8.01) -- (7.35,-8.16) -- (7.25,-8.23) -- (7.27,-8.29) -- (7.21,-8.28) -- (7.18,-8.39) -- (7.05,-8.44) -- (7.13,-8.50) -- (7.08,-8.55) -- (6.96,-8.52) -- (7.03,-8.69) -- (6.94,-8.69) -- (6.96,-8.81) -- (6.90,-8.84) -- (6.79,-8.81) -- (6.79,-8.90) -- (6.74,-8.93) -- (6.75,-8.75) -- (6.69,-8.71) -- (6.65,-8.75) -- (6.59,-8.74) -- (6.62,-8.86) -- (6.50,-8.94) -- (6.42,-8.88) -- (6.34,-8.88) -- (6.37,-8.79) -- (6.31,-8.72) -- (6.14,-8.89) -- (6.15,-8.97) -- (6.08,-8.95) -- (6.08,-9.03) -- (6.03,-9.07) -- (5.88,-9.05) -- (5.89,-9.13) -- (5.82,-9.21) -- (5.61,-9.13) -- (5.60,-9.29) -- (5.38,-9.29) -- (5.36,-9.33) -- (5.28,-9.34) -- (5.18,-9.29) -- (5.04,-9.38) -- (4.99,-9.36) -- (4.95,-9.26) -- (4.81,-9.20) -- (4.73,-9.10) -- (4.94,-8.79) -- (4.84,-8.67) -- (4.75,-8.70) -- (4.66,-8.79) -- (4.60,-8.98) -- (4.52,-9.09) -- (4.60,-9.21) -- (4.42,-9.37) -- (4.36,-9.39) -- (4.27,-9.29) -- (4.11,-9.29) -- (4.08,-9.25) -- (4.12,-9.21) -- (4.00,-9.08) -- (4.03,-9.03) -- (3.99,-8.96) -- (4.07,-8.92) -- (4.03,-8.85) -- (4.11,-8.68) -- (4.05,-8.35) -- (4.11,-8.26) -- (4.06,-8.17) -- (4.11,-8.13) -- (3.96,-7.86) -- (3.88,-7.79) -- (3.88,-7.57) -- (3.78,-7.49) -- (3.80,-7.41) -- (3.70,-7.40) -- (3.65,-7.45) -- (3.55,-7.21) -- (3.44,-7.34) -- (3.35,-7.27) -- (3.20,-7.33) -- (3.13,-7.43) -- (3.04,-7.36) -- (3.05,-7.41) -- (3.00,-7.45) -- (2.91,-7.50) -- (2.86,-7.46) -- (2.80,-7.52) -- (2.64,-7.51) -- (2.59,-7.34) -- (2.65,-7.30) -- (2.64,-7.19) -- (2.86,-7.21) -- (2.77,-7.07) -- (2.93,-7.05) -- (2.85,-6.94) -- (3.03,-6.76) -- (2.97,-6.73) -- (2.97,-6.59) -- (3.18,-6.71) -- (3.27,-6.55) -- (3.35,-6.57) -- (3.32,-6.76) -- (3.42,-6.78) -- (3.49,-6.61) -- (3.59,-6.66) -- (3.60,-6.45) -- (3.67,-6.37) -- (3.74,-6.36) -- (3.79,-6.40) -- (3.80,-6.56) -- (3.89,-6.60) -- (3.90,-6.40) -- (4.09,-6.40) -- (4.01,-6.30) -- (4.16,-6.29) -- (4.17,-6.32) -- (4.26,-6.24) -- cycle;
\filldraw[fill=primary!6,draw=black!25,line width=0.3pt] (4.50,-4.21) -- (4.49,-4.33) -- (4.37,-4.53) -- (4.42,-4.69) -- (4.22,-4.88) -- (4.24,-4.96) -- (4.38,-4.94) -- (4.40,-4.99) -- (4.33,-5.15) -- (4.38,-5.21) -- (4.43,-5.19) -- (4.47,-5.30) -- (4.32,-5.51) -- (4.34,-5.54) -- (4.28,-5.51) -- (4.28,-5.60) -- (4.21,-5.63) -- (4.30,-5.70) -- (4.21,-5.81) -- (4.20,-5.89) -- (4.10,-5.90) -- (4.01,-5.84) -- (4.05,-6.00) -- (4.18,-6.12) -- (4.20,-6.20) -- (4.15,-6.24) -- (4.16,-6.29) -- (4.01,-6.30) -- (4.08,-6.43) -- (3.99,-6.38) -- (3.90,-6.40) -- (3.89,-6.60) -- (3.80,-6.56) -- (3.81,-6.45) -- (3.74,-6.36) -- (3.67,-6.37) -- (3.60,-6.45) -- (3.59,-6.66) -- (3.49,-6.61) -- (3.42,-6.78) -- (3.32,-6.76) -- (3.35,-6.57) -- (3.27,-6.55) -- (3.18,-6.71) -- (2.97,-6.59) -- (2.91,-6.75) -- (2.82,-6.73) -- (2.81,-6.58) -- (2.46,-6.50) -- (2.42,-6.38) -- (2.28,-6.50) -- (2.13,-6.55) -- (1.94,-6.87) -- (1.79,-6.94) -- (1.59,-7.14) -- (1.12,-7.40) -- (1.02,-7.35) -- (0.99,-7.25) -- (0.79,-7.13) -- (0.70,-7.15) -- (0.61,-7.24) -- (0.55,-7.23) -- (0.71,-6.99) -- (0.70,-6.93) -- (0.64,-6.90) -- (0.54,-6.93) -- (0.37,-7.11) -- (0.13,-7.07) -- (0.16,-6.97) -- (0.00,-6.81) -- (0.11,-6.70) -- (0.17,-6.78) -- (0.25,-6.80) -- (0.38,-6.65) -- (0.54,-6.66) -- (0.57,-6.75) -- (0.74,-6.66) -- (0.69,-6.59) -- (0.71,-6.41) -- (0.85,-6.45) -- (0.91,-6.41) -- (0.84,-6.40) -- (0.95,-6.12) -- (1.04,-6.14) -- (1.04,-6.04) -- (1.12,-6.08) -- (1.18,-6.03) -- (1.10,-6.00) -- (1.09,-5.94) -- (1.15,-5.87) -- (1.21,-5.96) -- (1.20,-5.91) -- (1.25,-5.91) -- (1.27,-5.81) -- (1.20,-5.82) -- (1.16,-5.72) -- (1.19,-5.67) -- (1.23,-5.69) -- (1.22,-5.58) -- (1.26,-5.53) -- (1.36,-5.55) -- (1.32,-5.46) -- (1.44,-5.44) -- (1.42,-5.48) -- (1.51,-5.51) -- (1.49,-5.41) -- (1.63,-5.42) -- (1.60,-5.36) -- (1.65,-5.31) -- (1.65,-5.37) -- (1.70,-5.37) -- (1.73,-5.26) -- (1.84,-5.26) -- (1.93,-5.35) -- (1.98,-5.34) -- (1.94,-5.25) -- (2.04,-5.19) -- (1.98,-5.16) -- (2.04,-5.05) -- (2.09,-5.12) -- (2.18,-5.06) -- (2.17,-4.95) -- (2.20,-5.02) -- (2.28,-4.99) -- (2.34,-5.02) -- (2.29,-4.85) -- (2.40,-4.81) -- (2.37,-4.75) -- (2.55,-4.76) -- (2.56,-4.68) -- (2.64,-4.67) -- (2.63,-4.58) -- (2.70,-4.59) -- (2.91,-4.44) -- (3.00,-4.48) -- (3.22,-4.31) -- (3.60,-4.29) -- (3.72,-4.18) -- (3.88,-4.17) -- (4.00,-4.26) -- (4.13,-4.28) -- (4.41,-4.08) -- (4.55,-4.09) -- (4.56,-4.15) -- cycle;
\filldraw[fill=primary!6,draw=black!25,line width=0.3pt] (6.96,-0.01) -- (7.00,-0.15) -- (7.10,-0.05) -- (7.18,-0.08) -- (7.29,-0.05) -- (7.40,-0.22) -- (7.33,-0.30) -- (7.31,-0.26) -- (7.26,-0.31) -- (7.15,-0.27) -- (7.18,-0.34) -- (7.06,-0.42) -- (7.09,-0.50) -- (7.02,-0.55) -- (7.05,-0.61) -- (6.99,-0.60) -- (6.96,-0.68) -- (7.07,-0.82) -- (7.03,-0.82) -- (7.03,-0.90) -- (6.88,-0.93) -- (6.88,-0.99) -- (6.81,-1.07) -- (6.84,-1.17) -- (6.76,-1.23) -- (6.92,-1.27) -- (6.90,-1.31) -- (6.96,-1.39) -- (7.11,-1.50) -- (7.32,-1.56) -- (7.27,-1.66) -- (7.34,-1.74) -- (7.27,-1.81) -- (7.33,-1.87) -- (7.31,-2.34) -- (7.40,-2.51) -- (7.51,-2.53) -- (7.55,-2.65) -- (6.89,-2.69) -- (6.90,-2.76) -- (6.82,-2.72) -- (6.86,-2.80) -- (6.82,-2.79) -- (6.81,-2.84) -- (6.85,-2.92) -- (6.78,-2.92) -- (6.74,-2.85) -- (6.65,-3.01) -- (6.44,-2.99) -- (6.42,-2.93) -- (6.01,-3.09) -- (6.00,-3.05) -- (5.87,-3.06) -- (5.76,-3.16) -- (5.33,-3.01) -- (5.40,-2.94) -- (5.44,-2.78) -- (5.14,-2.53) -- (5.04,-2.56) -- (4.94,-2.45) -- (4.86,-2.50) -- (4.79,-2.38) -- (4.60,-2.44) -- (4.51,-2.32) -- (4.46,-2.36) -- (4.43,-2.31) -- (4.37,-2.33) -- (4.35,-2.26) -- (4.30,-2.29) -- (4.31,-2.14) -- (4.26,-2.17) -- (4.24,-2.12) -- (4.32,-2.02) -- (4.38,-2.13) -- (4.42,-2.10) -- (4.48,-2.16) -- (4.85,-2.18) -- (5.05,-2.30) -- (5.09,-2.13) -- (5.16,-2.11) -- (5.22,-1.84) -- (5.37,-1.83) -- (5.33,-1.78) -- (5.36,-1.71) -- (5.49,-1.75) -- (5.57,-1.70) -- (5.56,-1.63) -- (5.64,-1.70) -- (5.70,-1.57) -- (5.89,-1.49) -- (5.90,-1.42) -- (5.95,-1.49) -- (6.03,-1.42) -- (6.12,-1.50) -- (6.05,-1.22) -- (6.14,-1.17) -- (6.21,-1.22) -- (6.30,-1.09) -- (6.27,-1.05) -- (6.35,-1.08) -- (6.35,-0.99) -- (6.50,-1.00) -- (6.44,-0.95) -- (6.45,-0.89) -- (6.51,-0.88) -- (6.54,-0.82) -- (6.36,-0.70) -- (6.47,-0.60) -- (6.50,-0.66) -- (6.47,-0.55) -- (6.52,-0.54) -- (6.51,-0.50) -- (6.63,-0.53) -- (6.56,-0.45) -- (6.53,-0.35) -- (6.68,-0.23) -- (6.63,-0.15) -- (6.69,-0.02) -- (6.87,-0.08) -- cycle;
\filldraw[fill=primary!6,draw=black!25,line width=0.3pt] (7.67,0.00) -- (7.63,-0.16) -- (7.81,-0.13) -- (7.78,-0.24) -- (7.95,-0.29) -- (7.94,-0.19) -- (8.16,-0.09) -- (8.10,-0.32) -- (8.20,-0.35) -- (8.24,-0.24) -- (8.30,-0.22) -- (8.29,-0.35) -- (8.34,-0.35) -- (8.39,-0.44) -- (8.52,-0.42) -- (8.48,-0.49) -- (8.56,-0.56) -- (8.54,-0.70) -- (8.65,-0.74) -- (8.61,-0.83) -- (8.65,-0.90) -- (8.83,-0.85) -- (8.74,-0.99) -- (8.87,-0.95) -- (8.71,-1.13) -- (8.95,-1.13) -- (8.97,-1.21) -- (8.90,-1.17) -- (8.86,-1.23) -- (8.95,-1.37) -- (9.05,-1.37) -- (9.05,-1.28) -- (9.20,-1.30) -- (9.24,-1.21) -- (9.27,-1.32) -- (9.17,-1.34) -- (9.13,-1.49) -- (9.33,-1.51) -- (9.44,-1.72) -- (9.57,-1.71) -- (9.53,-1.87) -- (9.71,-1.88) -- (9.74,-1.95) -- (9.71,-1.99) -- (9.86,-2.12) -- (9.81,-2.19) -- (9.85,-2.38) -- (9.95,-2.39) -- (9.98,-2.58) -- (10.02,-2.60) -- (9.93,-2.74) -- (9.93,-2.89) -- (9.87,-3.00) -- (9.92,-3.10) -- (9.97,-3.14) -- (9.97,-3.06) -- (10.05,-3.05) -- (10.09,-3.15) -- (10.18,-3.13) -- (10.27,-3.35) -- (10.35,-3.32) -- (10.37,-3.38) -- (10.46,-3.41) -- (10.51,-3.48) -- (10.80,-3.59) -- (10.74,-3.68) -- (10.76,-3.82) -- (10.88,-3.97) -- (10.82,-4.22) -- (10.76,-4.28) -- (10.82,-4.39) -- (10.79,-4.47) -- (10.97,-4.69) -- (10.93,-4.74) -- (11.00,-4.83) -- (10.96,-4.87) -- (10.97,-4.95) -- (10.89,-4.86) -- (10.87,-4.78) -- (10.70,-4.71) -- (10.65,-4.62) -- (10.48,-4.67) -- (10.42,-4.63) -- (10.32,-4.72) -- (10.27,-4.66) -- (10.20,-4.68) -- (9.67,-5.05) -- (9.58,-5.11) -- (9.52,-5.08) -- (9.46,-5.16) -- (9.39,-5.15) -- (9.43,-5.24) -- (9.42,-5.27) -- (9.38,-5.24) -- (9.34,-5.34) -- (9.26,-5.28) -- (9.22,-5.31) -- (9.24,-5.37) -- (9.16,-5.36) -- (9.06,-5.43) -- (8.99,-5.57) -- (8.73,-5.66) -- (8.83,-5.55) -- (8.82,-5.51) -- (8.76,-5.50) -- (8.76,-5.44) -- (8.66,-5.43) -- (8.67,-5.39) -- (8.77,-5.38) -- (8.85,-5.22) -- (8.90,-5.23) -- (8.92,-5.17) -- (8.85,-5.13) -- (8.88,-5.05) -- (8.80,-5.02) -- (8.77,-5.22) -- (8.74,-5.17) -- (8.69,-5.23) -- (8.68,-5.07) -- (8.49,-5.11) -- (8.54,-5.06) -- (8.53,-5.00) -- (8.60,-5.01) -- (8.67,-4.91) -- (8.57,-4.92) -- (8.58,-4.87) -- (8.50,-4.87) -- (8.39,-4.93) -- (8.51,-4.78) -- (8.60,-4.78) -- (8.56,-4.71) -- (8.59,-4.66) -- (8.46,-4.69) -- (8.48,-4.63) -- (8.45,-4.62) -- (8.62,-4.59) -- (8.61,-4.44) -- (8.75,-4.45) -- (8.58,-4.43) -- (8.54,-4.27) -- (8.43,-4.39) -- (8.50,-4.48) -- (8.39,-4.46) -- (8.38,-4.30) -- (8.29,-4.24) -- (8.22,-4.32) -- (8.22,-4.28) -- (8.16,-4.25) -- (8.21,-4.23) -- (8.12,-4.11) -- (8.00,-4.06) -- (8.05,-4.03) -- (8.12,-4.07) -- (8.12,-4.00) -- (8.18,-4.07) -- (8.17,-3.84) -- (8.19,-3.81) -- (8.26,-3.87) -- (8.28,-3.82) -- (8.24,-3.75) -- (8.33,-3.71) -- (8.15,-3.58) -- (8.20,-3.54) -- (8.47,-3.57) -- (8.26,-3.42) -- (8.16,-3.48) -- (8.14,-3.40) -- (8.08,-3.37) -- (8.14,-3.32) -- (8.18,-3.37) -- (8.24,-3.23) -- (8.35,-3.16) -- (8.15,-3.12) -- (8.16,-3.04) -- (8.07,-2.98) -- (7.97,-3.00) -- (8.02,-2.85) -- (7.72,-2.78) -- (7.51,-2.61) -- (7.49,-2.51) -- (7.40,-2.51) -- (7.31,-2.34) -- (7.33,-1.87) -- (7.27,-1.81) -- (7.34,-1.74) -- (7.27,-1.66) -- (7.32,-1.56) -- (7.11,-1.50) -- (6.96,-1.39) -- (6.90,-1.31) -- (6.92,-1.27) -- (6.76,-1.23) -- (6.84,-1.17) -- (6.81,-1.07) -- (6.88,-0.99) -- (6.88,-0.93) -- (7.03,-0.90) -- (7.03,-0.82) -- (7.07,-0.82) -- (6.96,-0.68) -- (6.99,-0.60) -- (7.05,-0.61) -- (7.02,-0.55) -- (7.09,-0.50) -- (7.06,-0.42) -- (7.18,-0.34) -- (7.15,-0.27) -- (7.26,-0.31) -- (7.31,-0.26) -- (7.33,-0.30) -- (7.40,-0.22) -- (7.29,-0.02) -- (7.33,-0.02) -- (7.38,-0.12) -- (7.47,-0.11) -- (7.54,-0.17) -- (7.57,-0.04) -- cycle;
\filldraw[fill=primary!6,draw=black!25,line width=0.3pt] (7.19,-5.23) -- (7.28,-5.33) -- (7.34,-5.30) -- (7.35,-5.38) -- (7.49,-5.35) -- (7.49,-5.30) -- (7.55,-5.41) -- (7.88,-5.34) -- (7.94,-5.37) -- (8.00,-5.29) -- (8.08,-5.40) -- (8.08,-5.35) -- (8.13,-5.34) -- (8.20,-5.47) -- (8.27,-5.41) -- (8.34,-5.41) -- (8.39,-5.35) -- (8.50,-5.46) -- (8.61,-5.40) -- (8.59,-5.54) -- (8.63,-5.58) -- (8.55,-5.63) -- (8.59,-5.78) -- (8.68,-5.74) -- (8.84,-5.88) -- (8.86,-5.97) -- (9.06,-6.01) -- (9.09,-5.95) -- (9.23,-5.99) -- (9.22,-6.07) -- (9.29,-6.04) -- (9.32,-6.08) -- (9.63,-6.12) -- (9.50,-6.27) -- (9.48,-6.35) -- (9.56,-6.45) -- (9.65,-6.46) -- (9.66,-6.54) -- (9.61,-6.56) -- (9.60,-6.66) -- (9.69,-6.69) -- (9.71,-6.86) -- (9.80,-6.95) -- (9.72,-6.94) -- (9.65,-6.81) -- (9.58,-6.90) -- (9.46,-6.88) -- (9.43,-7.10) -- (9.33,-7.08) -- (9.28,-7.13) -- (9.20,-7.05) -- (9.08,-7.02) -- (9.07,-7.11) -- (8.98,-7.08) -- (8.74,-7.22) -- (8.70,-7.35) -- (8.75,-7.46) -- (8.68,-7.48) -- (8.66,-7.56) -- (8.44,-7.48) -- (8.49,-7.43) -- (8.34,-7.41) -- (8.38,-7.49) -- (8.27,-7.50) -- (8.28,-7.60) -- (8.15,-7.54) -- (8.17,-7.62) -- (8.06,-7.61) -- (8.03,-7.68) -- (7.92,-7.72) -- (7.90,-7.64) -- (7.86,-7.64) -- (7.86,-7.82) -- (7.79,-7.86) -- (7.63,-7.75) -- (7.55,-7.92) -- (7.55,-8.05) -- (7.49,-8.06) -- (7.48,-7.93) -- (7.40,-7.84) -- (7.40,-7.76) -- (7.25,-7.63) -- (7.18,-7.61) -- (7.06,-7.33) -- (7.14,-7.21) -- (7.12,-7.17) -- (7.20,-6.97) -- (7.01,-6.81) -- (6.99,-6.73) -- (6.77,-6.60) -- (6.84,-6.51) -- (6.97,-6.64) -- (7.03,-6.59) -- (7.02,-6.53) -- (7.26,-6.56) -- (7.35,-6.51) -- (7.38,-6.56) -- (7.41,-6.52) -- (7.53,-6.60) -- (7.57,-6.57) -- (7.56,-6.47) -- (7.70,-6.26) -- (7.71,-5.94) -- (7.55,-5.98) -- (7.38,-5.81) -- (7.26,-5.78) -- (7.19,-5.68) -- (7.28,-5.64) -- (7.27,-5.54) -- (7.11,-5.47) -- (7.04,-5.38) -- (7.04,-5.30) -- (7.17,-5.28) -- cycle;
\filldraw[fill=primary!6,draw=black!25,line width=0.3pt] (1.25,-2.92) -- (1.30,-2.95) -- (1.26,-3.02) -- (1.33,-3.06) -- (1.33,-2.97) -- (1.43,-3.02) -- (1.44,-2.93) -- (1.51,-2.99) -- (1.53,-2.93) -- (1.59,-2.99) -- (1.66,-2.98) -- (1.72,-3.06) -- (1.75,-2.96) -- (1.71,-2.87) -- (1.77,-2.85) -- (1.77,-2.92) -- (1.87,-2.94) -- (1.86,-2.88) -- (1.96,-2.85) -- (1.92,-2.87) -- (1.94,-2.91) -- (2.06,-2.91) -- (2.02,-2.97) -- (2.06,-3.01) -- (2.12,-2.92) -- (2.16,-2.94) -- (2.15,-2.88) -- (2.24,-2.87) -- (2.20,-2.93) -- (2.28,-3.03) -- (2.30,-2.98) -- (2.32,-2.99) -- (2.40,-2.88) -- (2.43,-2.88) -- (2.49,-2.96) -- (2.40,-2.89) -- (2.29,-3.05) -- (2.23,-3.05) -- (2.30,-3.23) -- (2.26,-3.30) -- (2.35,-3.52) -- (2.34,-3.59) -- (2.40,-3.65) -- (2.62,-3.68) -- (2.63,-3.75) -- (2.71,-3.79) -- (2.69,-3.88) -- (2.77,-4.02) -- (2.82,-4.00) -- (2.63,-4.35) -- (2.73,-4.40) -- (2.69,-4.45) -- (2.73,-4.56) -- (2.67,-4.61) -- (2.63,-4.58) -- (2.64,-4.67) -- (2.56,-4.68) -- (2.55,-4.76) -- (2.37,-4.75) -- (2.40,-4.81) -- (2.29,-4.85) -- (2.34,-5.02) -- (2.28,-4.99) -- (2.20,-5.02) -- (2.17,-4.95) -- (2.18,-5.06) -- (2.09,-5.12) -- (2.04,-5.05) -- (1.98,-5.16) -- (2.04,-5.19) -- (1.94,-5.25) -- (1.98,-5.34) -- (1.93,-5.35) -- (1.84,-5.26) -- (1.73,-5.26) -- (1.70,-5.37) -- (1.65,-5.37) -- (1.65,-5.31) -- (1.60,-5.36) -- (1.63,-5.42) -- (1.49,-5.41) -- (1.51,-5.25) -- (1.45,-5.24) -- (1.46,-5.29) -- (1.40,-5.24) -- (1.47,-5.21) -- (1.48,-5.13) -- (1.53,-5.15) -- (1.50,-5.09) -- (1.55,-5.10) -- (1.54,-5.04) -- (1.60,-5.04) -- (1.57,-5.00) -- (1.62,-4.95) -- (1.58,-4.95) -- (1.56,-4.89) -- (1.59,-4.82) -- (1.64,-4.88) -- (1.64,-4.83) -- (1.69,-4.83) -- (1.68,-4.75) -- (1.62,-4.72) -- (1.63,-4.64) -- (1.59,-4.62) -- (1.63,-4.58) -- (1.61,-4.52) -- (1.55,-4.54) -- (1.54,-4.50) -- (1.59,-4.47) -- (1.41,-4.31) -- (1.48,-4.20) -- (1.43,-4.14) -- (1.34,-4.14) -- (1.35,-4.07) -- (1.30,-4.02) -- (1.31,-3.75) -- (1.17,-3.70) -- (1.06,-3.51) -- (1.13,-3.46) -- (1.12,-3.41) -- (1.08,-3.37) -- (1.09,-3.27) -- (1.02,-3.23) -- (1.07,-3.22) -- (1.07,-3.18) -- (0.97,-3.11) -- (1.03,-3.08) -- (0.97,-3.03) -- (1.03,-3.01) -- (0.99,-2.86) -- (1.14,-2.88) -- (1.14,-2.94) -- (1.20,-2.84) -- (1.24,-2.85) -- cycle;
\filldraw[fill=primary!6,draw=black!25,line width=0.3pt] (2.81,-6.58) -- (2.82,-6.73) -- (2.91,-6.75) -- (2.92,-6.69) -- (3.03,-6.75) -- (2.99,-6.84) -- (2.85,-6.94) -- (2.93,-7.05) -- (2.77,-7.07) -- (2.86,-7.21) -- (2.64,-7.19) -- (2.65,-7.30) -- (2.61,-7.32) -- (2.51,-7.23) -- (2.34,-7.23) -- (2.33,-7.33) -- (2.16,-7.61) -- (1.82,-7.62) -- (1.74,-7.47) -- (1.67,-7.43) -- (1.50,-7.61) -- (1.23,-7.43) -- (1.22,-7.37) -- (1.28,-7.29) -- (1.59,-7.14) -- (1.79,-6.94) -- (1.94,-6.87) -- (2.13,-6.55) -- (2.28,-6.50) -- (2.42,-6.38) -- (2.46,-6.50) -- cycle;
\filldraw[fill=primary!6,draw=black!25,line width=0.3pt] (2.60,-7.46) -- (2.64,-7.51) -- (2.80,-7.52) -- (2.86,-7.46) -- (2.91,-7.50) -- (3.05,-7.41) -- (3.04,-7.36) -- (3.13,-7.43) -- (3.20,-7.33) -- (3.35,-7.27) -- (3.44,-7.34) -- (3.55,-7.21) -- (3.65,-7.45) -- (3.70,-7.40) -- (3.80,-7.41) -- (3.78,-7.49) -- (3.88,-7.57) -- (3.88,-7.79) -- (3.96,-7.86) -- (4.11,-8.13) -- (4.06,-8.17) -- (4.11,-8.26) -- (4.05,-8.35) -- (4.11,-8.68) -- (4.03,-8.85) -- (4.07,-8.92) -- (3.99,-8.96) -- (4.03,-9.03) -- (4.00,-9.08) -- (4.12,-9.21) -- (4.08,-9.25) -- (4.17,-9.30) -- (3.99,-9.40) -- (3.86,-9.38) -- (3.70,-9.50) -- (3.61,-9.48) -- (3.54,-9.37) -- (3.10,-9.21) -- (3.02,-9.22) -- (2.93,-9.41) -- (2.85,-9.43) -- (2.69,-9.32) -- (2.72,-9.16) -- (2.56,-9.10) -- (2.59,-8.99) -- (2.57,-8.91) -- (2.42,-8.83) -- (2.20,-8.84) -- (2.03,-8.52) -- (1.82,-8.49) -- (1.79,-8.35) -- (1.70,-8.28) -- (1.65,-8.03) -- (1.74,-7.97) -- (1.77,-7.87) -- (1.87,-7.76) -- (2.04,-7.75) -- (1.96,-7.62) -- (2.16,-7.61) -- (2.33,-7.33) -- (2.34,-7.23) -- (2.51,-7.23) -- (2.61,-7.32) -- cycle;
\filldraw[fill=primary!6,draw=black!25,line width=0.3pt] (2.07,-1.55) -- (1.96,-1.71) -- (2.09,-1.74) -- (2.14,-1.83) -- (2.18,-1.80) -- (2.13,-1.75) -- (2.19,-1.77) -- (2.17,-1.72) -- (2.21,-1.74) -- (2.26,-1.88) -- (2.22,-1.97) -- (2.30,-1.90) -- (2.27,-1.82) -- (2.38,-1.81) -- (2.46,-1.69) -- (2.57,-1.66) -- (2.48,-1.76) -- (2.62,-1.76) -- (2.56,-1.82) -- (2.63,-1.86) -- (2.66,-1.81) -- (2.74,-1.90) -- (2.78,-1.78) -- (2.84,-1.82) -- (2.91,-1.80) -- (2.83,-1.97) -- (2.87,-2.00) -- (2.92,-1.95) -- (2.91,-2.07) -- (2.97,-2.10) -- (3.01,-2.08) -- (2.97,-2.01) -- (3.11,-1.86) -- (3.15,-1.87) -- (3.15,-2.01) -- (3.22,-1.93) -- (3.27,-1.95) -- (3.29,-2.05) -- (3.33,-2.02) -- (3.39,-2.06) -- (3.37,-1.96) -- (3.43,-1.95) -- (3.61,-2.03) -- (3.60,-2.09) -- (3.65,-2.13) -- (3.53,-2.18) -- (3.58,-2.20) -- (3.57,-2.33) -- (3.60,-2.29) -- (3.64,-2.33) -- (3.58,-2.36) -- (3.65,-2.40) -- (3.60,-2.52) -- (3.65,-2.61) -- (3.55,-3.14) -- (3.66,-3.56) -- (3.56,-3.73) -- (3.64,-3.74) -- (3.72,-3.83) -- (3.58,-4.00) -- (3.49,-3.92) -- (3.47,-3.97) -- (3.38,-3.98) -- (3.44,-4.04) -- (3.39,-4.06) -- (3.39,-4.12) -- (3.33,-4.14) -- (3.28,-4.10) -- (3.30,-4.18) -- (3.25,-4.18) -- (3.20,-4.10) -- (3.21,-4.14) -- (3.16,-4.20) -- (3.07,-4.14) -- (3.01,-4.21) -- (2.98,-4.07) -- (2.93,-4.05) -- (2.94,-4.13) -- (2.81,-4.13) -- (2.86,-4.20) -- (2.82,-4.20) -- (2.76,-4.13) -- (2.82,-4.00) -- (2.77,-4.02) -- (2.69,-3.88) -- (2.71,-3.79) -- (2.63,-3.75) -- (2.62,-3.68) -- (2.40,-3.65) -- (2.34,-3.59) -- (2.35,-3.52) -- (2.26,-3.30) -- (2.30,-3.23) -- (2.23,-3.05) -- (2.29,-3.05) -- (2.40,-2.89) -- (2.46,-2.95) -- (2.49,-2.93) -- (2.40,-2.88) -- (2.30,-3.02) -- (2.22,-2.99) -- (2.20,-2.93) -- (2.29,-2.74) -- (2.20,-2.60) -- (2.05,-2.62) -- (1.99,-2.46) -- (2.01,-2.24) -- (1.92,-2.21) -- (1.92,-2.16) -- (1.80,-2.16) -- (1.72,-1.98) -- (1.67,-2.02) -- (1.66,-1.95) -- (1.71,-1.96) -- (1.72,-1.91) -- (1.67,-1.87) -- (1.71,-1.82) -- (1.64,-1.77) -- (1.56,-1.79) -- (1.57,-1.73) -- (1.52,-1.71) -- (1.52,-1.63) -- (1.59,-1.65) -- (1.65,-1.56) -- (1.74,-1.67) -- (1.82,-1.58) -- (1.88,-1.58) -- (1.92,-1.66) -- (1.91,-1.59) -- cycle;
\filldraw[fill=primary!6,draw=black!25,line width=0.3pt] (3.80,-1.78) -- (3.85,-1.86) -- (3.88,-1.82) -- (4.00,-1.90) -- (3.98,-1.94) -- (4.32,-2.02) -- (4.24,-2.12) -- (4.26,-2.17) -- (4.31,-2.14) -- (4.30,-2.29) -- (4.35,-2.26) -- (4.37,-2.33) -- (4.43,-2.31) -- (4.46,-2.36) -- (4.51,-2.32) -- (4.60,-2.44) -- (4.79,-2.38) -- (4.86,-2.50) -- (4.97,-2.46) -- (5.01,-2.54) -- (4.96,-2.61) -- (5.12,-2.89) -- (5.10,-3.10) -- (5.23,-3.28) -- (5.21,-3.35) -- (5.11,-3.33) -- (5.15,-3.46) -- (5.07,-3.52) -- (4.87,-3.50) -- (4.74,-3.56) -- (4.50,-3.35) -- (4.22,-3.58) -- (3.64,-3.50) -- (3.55,-3.14) -- (3.65,-2.61) -- (3.60,-2.52) -- (3.65,-2.40) -- (3.58,-2.36) -- (3.64,-2.33) -- (3.60,-2.29) -- (3.57,-2.33) -- (3.58,-2.20) -- (3.53,-2.18) -- (3.65,-2.13) -- (3.60,-2.09) -- (3.61,-2.03) -- (3.50,-1.96) -- (3.54,-1.93) -- (3.50,-1.85) -- (3.69,-1.87) -- (3.67,-1.81) -- (3.75,-1.82) -- (3.72,-1.75) -- cycle;
\draw[black!35,dashed,line width=0.4pt] (6.44,-4.79) -- (9.67,-3.44);
\draw[black!35,dashed,line width=0.4pt] (6.44,-4.79) -- (2.85,-3.28);
\draw[black!35,dashed,line width=0.4pt] (6.44,-4.79) -- (3.98,-2.44);
\draw[black!35,dashed,line width=0.4pt] (5.08,-7.04) -- (2.51,-6.14);
\draw[black!35,dashed,line width=0.4pt] (5.08,-7.04) -- (1.94,-4.63);
\draw[black!35,dashed,line width=0.4pt] (5.08,-7.04) -- (1.99,-7.35);
\draw[black!35,dashed,line width=0.4pt] (5.08,-7.04) -- (2.61,-8.05);
\draw[black!35,dashed,line width=0.4pt] (2.51,-6.14) -- (2.61,-8.05);
\draw[black!35,dashed,line width=0.4pt] (2.51,-6.14) -- (2.85,-3.28);
\draw[black!35,dashed,line width=0.4pt] (2.51,-6.14) -- (3.98,-2.44);
\draw[black!35,dashed,line width=0.4pt] (5.97,-2.17) -- (9.67,-3.44);
\draw[black!35,dashed,line width=0.4pt] (5.97,-2.17) -- (2.85,-3.28);
\draw[black!35,dashed,line width=0.4pt] (1.94,-4.63) -- (1.99,-7.35);
\draw[black!35,dashed,line width=0.4pt] (1.94,-4.63) -- (2.61,-8.05);
\draw[black!35,dashed,line width=0.4pt] (1.94,-4.63) -- (3.98,-2.44);
\draw[primary,line width=1.6pt] (6.44,-4.79) -- node[fill=white,inner sep=1pt,font=\scriptsize\bfseries,text=primary] {39,1} (5.08,-7.04);
\draw[primary,line width=1.6pt] (6.44,-4.79) -- node[fill=white,inner sep=1pt,font=\scriptsize\bfseries,text=primary] {39,7} (5.97,-2.17);
\draw[primary,line width=1.6pt] (6.44,-4.79) -- node[fill=white,inner sep=1pt,font=\scriptsize\bfseries,text=primary] {42,9} (9.01,-6.10);
\draw[primary,line width=1.6pt] (2.51,-6.14) -- node[fill=white,inner sep=1pt,font=\scriptsize\bfseries,text=primary] {23,7} (1.94,-4.63);
\draw[primary,line width=1.6pt] (2.51,-6.14) -- node[fill=white,inner sep=1pt,font=\scriptsize\bfseries,text=primary] {19,6} (1.99,-7.35);
\draw[primary,line width=1.6pt] (5.97,-2.17) -- node[fill=white,inner sep=1pt,font=\scriptsize\bfseries,text=primary] {29,7} (3.98,-2.44);
\draw[primary,line width=1.6pt] (9.67,-3.44) -- node[fill=white,inner sep=1pt,font=\scriptsize\bfseries,text=primary] {40,9} (9.01,-6.10);
\draw[primary,line width=1.6pt] (1.94,-4.63) -- node[fill=white,inner sep=1pt,font=\scriptsize\bfseries,text=primary] {23,9} (2.85,-3.28);
\draw[primary,line width=1.6pt] (1.99,-7.35) -- node[fill=white,inner sep=1pt,font=\scriptsize\bfseries,text=primary] {13,8} (2.61,-8.05);
\draw[primary,line width=1.6pt] (2.85,-3.28) -- node[fill=white,inner sep=1pt,font=\scriptsize\bfseries,text=primary] {21,3} (3.98,-2.44);
\fill (6.44,-4.79) circle (1.8pt) node[below=1pt] {\bfseries Guarapuava};
\fill (5.08,-7.04) circle (1.8pt) node[below=1pt] {Pinhão};
\fill (2.51,-6.14) circle (1.8pt) node[below=1pt] {Candói};
\fill (5.97,-2.17) circle (1.8pt) node[above=1pt] {Turvo};
\fill (9.67,-3.44) circle (1.8pt) node[below=1pt] {Prudentópolis};
\fill (9.01,-6.10) circle (1.8pt) node[below=1pt] {Inácio Martins};
\fill (1.94,-4.63) circle (1.8pt) node[above=1pt] {Cantagalo};
\fill (1.99,-7.35) circle (1.8pt) node[above=1pt] {Foz do Jordão};
\fill (2.61,-8.05) circle (1.8pt) node[below=1pt] {Reserva do Iguaçu};
\fill (2.85,-3.28) circle (1.8pt) node[above=1pt] {Goioxim};
\fill (3.98,-2.44) circle (1.8pt) node[above=1pt] {Campina do Simão};
\draw[primary,line width=1.6pt] (0.2,-9.80) -- ++(0.6,0) node[right,text=black] {linha escolhida};
\draw[black!35,dashed,line width=0.4pt] (3.6,-9.80) -- ++(0.6,0) node[right,text=black] {trecho descartado};
\draw[line width=0.6pt] (8,-9.75) -- ++(0,-0.1) -- ++(2.68,0) -- ++(0,0.1); \node[above] at (9.34,-9.80) {40 km};
\end{tikzpicture}}
\par\smallskip{\footnotesize Ao fundo, os 11 municípios. Distâncias em km.}
\end{figura}

Com essas distâncias, o Kruskal chega à rede da \cref{fig:rede-final}: 10 linhas
que ligam as 11 cidades com 294,6~km no total. A seção seguinte mostra a execução
passo a passo.

\subsection{Execução}

Entrada: arquivo \texttt{dados/rede\_eletrica.txt} (opção 2 do menu). A
\cref{tab:execucao-rede} mostra a saída do programa, na ordem em que o Kruskal
aceitou os trechos.

\begin{table}[H]
\centering
\caption{Trechos aceitos pelo Kruskal na rede elétrica}\label{tab:execucao-rede}
\begin{tabular}{clr}
\toprule
Ordem & Trecho & Extensão (km) \\
\midrule
1 & Foz do Jordão -- Reserva do Iguaçu & 13,8 \\
2 & Candói -- Foz do Jordão & 19,6 \\
3 & Goioxim -- Campina do Simão & 21,3 \\
4 & Candói -- Cantagalo & 23,7 \\
5 & Cantagalo -- Goioxim & 23,9 \\
6 & Turvo -- Campina do Simão & 29,7 \\
7 & Guarapuava -- Pinhão & 39,1 \\
8 & Guarapuava -- Turvo & 39,7 \\
9 & Prudentópolis -- Inácio Martins & 40,9 \\
10 & Guarapuava -- Inácio Martins & 42,9 \\
\midrule
 & \textbf{Total} & \textbf{294,6} \\
\bottomrule
\end{tabular}
\fonte{saída do programa elaborado pelo autor (2026).}
\end{table}

\subsection{Análise dos resultados}

A rede tem 10 trechos para 11 cidades, exatamente $n-1$, como garante o
\cref{teo:n-1}: todas as cidades ficam ligadas e não há nenhum trecho
redundante. Pelo \cref{teo:kruskal}, nenhuma outra combinação dos 25 trechos
possíveis liga as 11 cidades com menos de 294,6~km, menos de um terço dos
1.016,1~km que custaria construir todos.

O resultado também mostra o caráter ``de ilhas'' do Kruskal. Os primeiros trechos
aceitos formam um grupo a oeste (Foz do Jordão, Reserva do Iguaçu, Candói,
Cantagalo, Goioxim, Campina do Simão), enquanto Guarapuava só entra no sétimo
passo. O trecho Pinhão -- Reserva do Iguaçu (40,1~km) foi descartado porque,
quando chegou a vez dele, as duas cidades já estavam no mesmo grupo: ligá-las
fecharia um ciclo. Guarapuava acaba como o centro da rede, com três trechos, o
que é coerente com seu papel de polo regional.

% ============================================================
\section{Extensão: dividir pedidos entre motoboys}\label{sec:motoboys}
% ============================================================
Parando o Kruskal quando restam $k$ grupos, em vez de 1, o mesmo código divide 12
pedidos de entrega entre 3 motoboys, cada um com 4 pedidos próximos entre si.

\subsection{Descrição do problema}

Um restaurante tem 12 pedidos para entregar à noite e 3 motoboys. Quer dividir os
pedidos em 3 regiões, de forma que cada motoboy atenda endereços próximos uns dos
outros.

\begin{figura}{12 pedidos espalhados pela cidade}{fig:pedidos}{elaborada pelo autor (2026), com dados fictícios.}
\resizebox{0.78\textwidth}{!}{\begin{tikzpicture}[x=1cm,y=1cm,font=\footnotesize]
\fill[black!5,rounded corners=4pt] (0,0) rectangle (11.2,-11.2);
\begin{scope}
\clip[rounded corners=4pt] (0,0) rectangle (11.2,-11.2);
\draw[black!18,line width=4.0mm] (2.10,0.00) -- (2.10,-11.20);
\draw[black!18,line width=2.6mm] (4.22,-3.20) -- (4.22,-11.20);
\draw[black!18,line width=2.6mm] (6.18,0.00) -- (6.18,-11.20);
\draw[black!18,line width=2.6mm] (7.18,-5.14) -- (7.18,-9.68);
\draw[black!18,line width=2.6mm] (8.10,0.00) -- (8.10,-9.68);
\draw[black!18,line width=2.6mm] (9.46,-3.20) -- (9.46,-9.68);
\draw[black!18,line width=2.6mm] (0.90,-5.14) -- (0.90,-11.20);
\draw[black!18,line width=2.6mm] (10.30,0.00) -- (10.30,-5.14);
\draw[black!18,line width=2.6mm] (6.18,-1.62) -- (11.20,-1.62);
\draw[black!18,line width=4.0mm] (0.00,-3.20) -- (11.20,-3.20);
\draw[black!18,line width=2.6mm] (0.00,-5.14) -- (11.20,-5.14);
\draw[black!18,line width=2.6mm] (2.10,-6.98) -- (11.20,-6.98);
\draw[black!18,line width=2.6mm] (6.18,-8.08) -- (11.20,-8.08);
\draw[black!18,line width=2.6mm] (0.00,-9.68) -- (11.20,-9.68);
\draw[black!18,line width=2.6mm] (0.00,-0.88) -- (6.18,-0.88);
\end{scope}
\draw[black!40,rounded corners=4pt] (0,0) rectangle (11.2,-11.2);
\node[draw,thick,fill=white,minimum size=5mm,inner sep=1pt,font=\scriptsize\bfseries] (R) at (6.18,-6.98) {R};
\node[below=1pt of R,font=\scriptsize\bfseries,fill=black!5,inner sep=1pt] {Restaurante};
\draw[line width=0.6pt] (9.4,-10.7) -- ++(0,0.1) -- ++(1,0) -- ++(0,-0.1); \node[above] at (9.9,-10.6) {1 km};
\node[pino=black,label={[font=\scriptsize\bfseries]right:P01}] at (7.44,-3.20) {};
\node[pino=black,label={[font=\scriptsize\bfseries]left:P02}] at (1.76,-3.18) {};
\node[pino=black,label={[font=\scriptsize\bfseries]right:P03}] at (9.46,-6.98) {};
\node[pino=black,label={[font=\scriptsize\bfseries]right:P04}] at (8.10,-3.70) {};
\node[pino=black,label={[font=\scriptsize\bfseries]right:P05}] at (2.18,-4.12) {};
\node[pino=black,label={[font=\scriptsize\bfseries]left:P06}] at (7.18,-6.42) {};
\node[pino=black,label={[font=\scriptsize\bfseries]right:P07}] at (7.80,-1.62) {};
\node[pino=black,label={[font=\scriptsize\bfseries]right:P08}] at (4.22,-5.14) {};
\node[pino=black,label={[font=\scriptsize\bfseries]right:P09}] at (8.10,-6.08) {};
\node[pino=black,label={[font=\scriptsize\bfseries]left:P10}] at (6.18,-3.34) {};
\node[pino=black,label={[font=\scriptsize\bfseries]left:P11}] at (2.10,-4.52) {};
\node[pino=black,label={[font=\scriptsize\bfseries]left:P12}] at (8.16,-8.08) {};
\end{tikzpicture}}
\par\smallskip{\footnotesize Mapa ilustrativo em escala. Faixas cinza são ruas; P01 a P12 são os pedidos; R é o restaurante.}
\end{figura}

Só olhando, não dá para dizer com certeza onde termina um grupo e começa outro:
pedidos como o P08 ficam no meio do caminho. A pergunta é como o Kruskal decide,
usando apenas as distâncias.

\subsection{Modelagem como grafo}

\begin{table}[H]
\centering
\caption{Modelagem da divisão de pedidos como grafo}\label{tab:modelo-motoboys}
\begin{tabularx}{\textwidth}{l >{\raggedright\arraybackslash}X}
\toprule
Elemento & Significado no problema \\
\midrule
Vértice & Um pedido de entrega, com posição $(x, y)$ em km num mapa da cidade (12 no total) \\
Aresta & A ligação entre dois pedidos; todo pedido liga a todo pedido (66 arestas) \\
Peso & Distância em linha reta entre os dois pedidos, em km \\
Direção & Não direcionado: a distância de A até B é a mesma de B até A \\
Resultado desejado & Floresta com $k = 3$ árvores: cada árvore é o grupo de pedidos de um motoboy \\
\bottomrule
\end{tabularx}
\fonte{elaborada pelo autor (2026).}
\end{table}

\textbf{Dados.} Os 12 pedidos e suas posições são fictícios, sorteados em torno de
três regiões próximas da cidade (arquivo \texttt{dados/motoboys\_coordenadas.txt}).

\textbf{Por que funciona.} O Kruskal junta sempre os dois pedidos mais próximos
que ainda estão em grupos diferentes. Cada aresta aceita reduz o número de grupos
em 1 (\cref{teo:n-1}), então parar com $n - k = 12 - 3 = 9$ arestas deixa
exatamente 3 grupos. O que fica separado no fim é o que está mais distante. Em
estatística esse método se chama clusterização por ligação simples e equivale a
construir a árvore geradora mínima e remover suas $k-1$ arestas mais pesadas.

\subsection{Execução}

Entrada: arquivo \texttt{dados/motoboys.txt} (opção 3 do menu, com $k = 3$). O
programa aceitou 9 arestas (11,71~km somados) e formou os grupos da
\cref{tab:grupos}.

\begin{table}[H]
\centering
\caption{Grupos formados com $k = 3$}\label{tab:grupos}
\begin{tabular}{cl}
\toprule
Motoboy & Pedidos \\
\midrule
1 & P01, P04, P07, P10 \\
2 & P02, P05, P08, P11 \\
3 & P03, P06, P09, P12 \\
\bottomrule
\end{tabular}
\fonte{saída do programa elaborado pelo autor (2026).}
\end{table}

\begin{figura}{Com $k = 3$, o Kruskal forma 3 grupos, um por motoboy}{fig:grupos}{elaborada pelo autor (2026), com dados fictícios e a saída do programa.}
\resizebox{0.78\textwidth}{!}{\begin{tikzpicture}[x=1cm,y=1cm,font=\footnotesize]
\fill[black!5,rounded corners=4pt] (0,0) rectangle (11.2,-11.2);
\begin{scope}
\clip[rounded corners=4pt] (0,0) rectangle (11.2,-11.2);
\draw[black!18,line width=4.0mm] (2.10,0.00) -- (2.10,-11.20);
\draw[black!18,line width=2.6mm] (4.22,-3.20) -- (4.22,-11.20);
\draw[black!18,line width=2.6mm] (6.18,0.00) -- (6.18,-11.20);
\draw[black!18,line width=2.6mm] (7.18,-5.14) -- (7.18,-9.68);
\draw[black!18,line width=2.6mm] (8.10,0.00) -- (8.10,-9.68);
\draw[black!18,line width=2.6mm] (9.46,-3.20) -- (9.46,-9.68);
\draw[black!18,line width=2.6mm] (0.90,-5.14) -- (0.90,-11.20);
\draw[black!18,line width=2.6mm] (10.30,0.00) -- (10.30,-5.14);
\draw[black!18,line width=2.6mm] (6.18,-1.62) -- (11.20,-1.62);
\draw[black!18,line width=4.0mm] (0.00,-3.20) -- (11.20,-3.20);
\draw[black!18,line width=2.6mm] (0.00,-5.14) -- (11.20,-5.14);
\draw[black!18,line width=2.6mm] (2.10,-6.98) -- (11.20,-6.98);
\draw[black!18,line width=2.6mm] (6.18,-8.08) -- (11.20,-8.08);
\draw[black!18,line width=2.6mm] (0.00,-9.68) -- (11.20,-9.68);
\draw[black!18,line width=2.6mm] (0.00,-0.88) -- (6.18,-0.88);
\end{scope}
\draw[black!40,rounded corners=4pt] (0,0) rectangle (11.2,-11.2);
\node[draw,thick,fill=white,minimum size=5mm,inner sep=1pt,font=\scriptsize\bfseries] (R) at (6.18,-6.98) {R};
\node[below=1pt of R,font=\scriptsize\bfseries,fill=black!5,inner sep=1pt] {Restaurante};
\draw[line width=0.6pt] (9.4,-10.7) -- ++(0,0.1) -- ++(1,0) -- ++(0,-0.1); \node[above] at (9.9,-10.6) {1 km};
\draw[grupoB,line width=1.4pt] (2.18,-4.12) -- node[right,font=\scriptsize\bfseries,text=grupoB,inner sep=1.5pt] {0,41} (2.10,-4.52);
\draw[grupoA,line width=1.4pt] (7.44,-3.20) -- node[below left,font=\scriptsize\bfseries,text=grupoA,inner sep=1.5pt] {0,82} (8.10,-3.70);
\draw[grupoC,line width=1.4pt] (7.18,-6.42) -- node[below,font=\scriptsize\bfseries,text=grupoC,inner sep=1.5pt] {0,98} (8.10,-6.08);
\draw[grupoB,line width=1.4pt] (1.76,-3.18) -- node[left,font=\scriptsize\bfseries,text=grupoB,inner sep=1.5pt] {1,04} (2.18,-4.12);
\draw[grupoA,line width=1.4pt] (7.44,-3.20) -- node[below,font=\scriptsize\bfseries,text=grupoA,inner sep=1.5pt] {1,29} (6.18,-3.34);
\draw[grupoA,line width=1.4pt] (7.44,-3.20) -- node[right,font=\scriptsize\bfseries,text=grupoA,inner sep=1.5pt] {1,61} (7.80,-1.62);
\draw[grupoC,line width=1.4pt] (9.46,-6.98) -- node[above right,font=\scriptsize\bfseries,text=grupoC,inner sep=1.5pt] {1,64} (8.10,-6.08);
\draw[grupoC,line width=1.4pt] (9.46,-6.98) -- node[right,font=\scriptsize\bfseries,text=grupoC,inner sep=1.5pt] {1,70} (8.16,-8.08);
\draw[grupoB,line width=1.4pt] (4.22,-5.14) -- node[below,font=\scriptsize\bfseries,text=grupoB,inner sep=1.5pt] {2,22} (2.10,-4.52);
\node[pino=grupoA,label={[font=\scriptsize\bfseries]right:P01}] at (7.44,-3.20) {};
\node[pino=grupoB,label={[font=\scriptsize\bfseries]left:P02}] at (1.76,-3.18) {};
\node[pino=grupoC,label={[font=\scriptsize\bfseries]right:P03}] at (9.46,-6.98) {};
\node[pino=grupoA,label={[font=\scriptsize\bfseries]right:P04}] at (8.10,-3.70) {};
\node[pino=grupoB,label={[font=\scriptsize\bfseries]right:P05}] at (2.18,-4.12) {};
\node[pino=grupoC,label={[font=\scriptsize\bfseries]left:P06}] at (7.18,-6.42) {};
\node[pino=grupoA,label={[font=\scriptsize\bfseries]right:P07}] at (7.80,-1.62) {};
\node[pino=grupoB,label={[font=\scriptsize\bfseries]right:P08}] at (4.22,-5.14) {};
\node[pino=grupoC,label={[font=\scriptsize\bfseries]right:P09}] at (8.10,-6.08) {};
\node[pino=grupoA,label={[font=\scriptsize\bfseries]left:P10}] at (6.18,-3.34) {};
\node[pino=grupoB,label={[font=\scriptsize\bfseries]left:P11}] at (2.10,-4.52) {};
\node[pino=grupoC,label={[font=\scriptsize\bfseries]left:P12}] at (8.16,-8.08) {};
\end{tikzpicture}}
\par\smallskip{\footnotesize Cada cor é um motoboy. Números coloridos: distância em linha reta de cada ligação, em km (total 11,71~km).}
\end{figura}

\subsection{Análise}

Os três grupos coincidem com as três regiões de onde os pedidos foram sorteados,
mesmo com elas próximas. A decisão foi apertada: a maior ligação aceita mede
2,22~km (P08--P11), e a próxima que juntaria dois grupos, P04--P09, mede 2,38~km,
apenas 160 metros a mais. A extensão mostra que a função \texttt{kruskal(g, k)}
resolve dois problemas diferentes trocando apenas o valor de $k$. Há, porém, uma
limitação: o método não equilibra o número de pedidos por motoboy; aqui cada
grupo ficou com 4 pedidos, mas com outros dados um motoboy poderia ficar com bem
mais pedidos que os outros.

% ============================================================
\section{Conclusão}
% ============================================================
O algoritmo de Kruskal resolve o problema da árvore geradora mínima com uma regra
simples, pegar sempre a aresta mais barata que não fecha ciclo, e o argumento de
troca garante que essa escolha gulosa sempre leva ao custo mínimo. Com o
Union-Find usando compressão de caminho e união por tamanho, o teste de ciclo
fica praticamente instantâneo e o custo total é dominado pela ordenação,
$O(E \log E)$.

A implementação em C++ aplicou o método a uma rede elétrica entre 11 cidades da
região de Guarapuava, ligando todas com 294,6~km de linha, e, com uma única
mudança no critério de parada, dividiu pedidos de entrega entre motoboys. Isso
mostra que o mesmo algoritmo atende problemas que, à primeira vista, parecem bem
diferentes.

\nocite{fiset2017intro,fiset2017ops,fiset2017path,fiset2017kruskal}
\references

\end{document}
